\documentclass[12pt,a4paper,openany]{book}
\maxdeadcycles=1000

\usepackage{fontspec}
\setmainfont{Latin Modern Roman}
\setmonofont{Latin Modern Mono}
\usepackage{amsmath,amssymb,amsthm,amsfonts}
\usepackage{mathtools}
\usepackage{tikz}
\usetikzlibrary{arrows.meta,positioning,calc}
\usepackage{pgfplots}
\pgfplotsset{compat=1.18}
\usepackage{tcolorbox}
\tcbuselibrary{theorems,skins,breakable}
\usepackage{enumitem}
\usepackage{fancyhdr}
\usepackage{titlesec}
\usepackage{geometry}
\usepackage{hyperref}
\usepackage{xcolor}
\usepackage{booktabs}
\usepackage{multicol}

\geometry{a4paper,margin=2.5cm,top=3cm,bottom=3cm}

% B&W-friendly colors
\definecolor{deepblue}{HTML}{1a1a6e}
\definecolor{accentblue}{HTML}{333366}
\definecolor{lightblue}{HTML}{f0f0f0}
\definecolor{prereqbg}{HTML}{f5f5f0}
\definecolor{prereqborder}{HTML}{444444}
\definecolor{examplebg}{HTML}{f0f0f0}
\definecolor{exampleborder}{HTML}{333333}
\definecolor{warningbg}{HTML}{f5f0f0}
\definecolor{warningborder}{HTML}{222222}
\definecolor{thmblue}{HTML}{222255}
\definecolor{defgreen}{HTML}{224422}

\pagestyle{fancy}
\fancyhf{}
\fancyhead[LE]{\textit{\nouppercase{\leftmark}}}
\fancyhead[RO]{\textit{\nouppercase{\rightmark}}}
\fancyfoot[C]{\thepage}
\renewcommand{\headrulewidth}{0.4pt}

\titleformat{\chapter}[display]
  {\normalfont\Huge\bfseries\color{deepblue}}
  {\chaptertitlename\ \thechapter}{20pt}{\Huge}
\titleformat{\section}
  {\normalfont\Large\bfseries\color{accentblue}}
  {\thesection}{1em}{}
\titleformat{\subsection}
  {\normalfont\large\bfseries\color{accentblue!80!black}}
  {\thesubsection}{1em}{}

\newtcolorbox{prereqbox}{}
  colback=prereqbg, colframe=prereqborder,
  title={\textbf{Prerequisites -- What You Need to Know}},
  fonttitle=\bfseries\color{black},
  breakable, left=8pt, right=8pt, boxrule=2pt, arc=4pt,
}

\renewcommand{\qedsymbol}{$\blacksquare$}
\hypersetup{colorlinks=true,linkcolor=deepblue,urlcolor=accentblue}

\begin{document}

% Cover
\begin{titlepage}
\begin{tikzpicture}[remember picture,overlay]
  \fill[deepblue] (current page.south west) rectangle (current page.north east);
  \node[anchor=center,white,font=\fontsize{48}{56}\selectfont\bfseries] 
    at ([yshift=4cm]current page.center) {MSc Mathematics};
  \node[anchor=center,white,font=\fontsize{28}{34}\selectfont] 
    at ([yshift=1.5cm]current page.center) {Semester 1 --- Complete Study Guide};
  \draw[white,line width=2pt] ([yshift=-0.5cm,xshift=-5cm]current page.center) -- 
    ([yshift=-0.5cm,xshift=5cm]current page.center);
  \node[anchor=center,white!80,font=\fontsize{14}{18}\selectfont] 
    at ([yshift=-2cm]current page.center) {Manipal University Jaipur --- Distance Education};
  \node[anchor=center,white!70,font=\fontsize{12}{16}\selectfont\itshape] 
    at ([yshift=-3.5cm]current page.center) {Prepared for CSIR NET Qualification \& Assistant Professorship};
  \node[anchor=center,white!60,font=\fontsize{10}{14}\selectfont] 
    at ([yshift=-5.5cm]current page.center) {6 Subjects $\cdot$ 60 Units $\cdot$ 35+ Proofs $\cdot$ 120+ Examples};
\end{tikzpicture}
\end{titlepage}

\tableofcontents
\newpage

\chapter{Prerequisites --- Class 12 Refresher}


\begin{quote}
\textbf{Read this chapter FIRST if it's been years since you studied Class 12 Maths.}
\textbf{This covers everything you need before starting MSc Mathematics.}
\end{quote}

\bigskip\noindent\rule{\textwidth}{0.4pt}\bigskip


A \textbf{set} is a well-defined collection of objects (called \textbf{elements}).

\textbf{Notation:}
\begin{itemize}[leftmargin=*]
  \item x $\in$ A means "x belongs to set A"
  \item x $\notin$ A means "x does NOT belong to A"
  \item A $\subset$ B means "A is a subset of B" (every element of A is in B)
  \item A $\cup$ B = union (elements in A OR B)
  \item A $\cap$ B = intersection (elements in A AND B)
  \item A \ B = difference (elements in A but NOT in B)
  \item $\emptyset$ = empty set (no elements)
\end{itemize}

\begin{tcolorbox}[colback=warningbg!30,colframe=warningborder,title=\textbf{Important -- Exam Alert},breakable,boxrule=1.5pt]
\textbf{Important Number Sets:}
```
$\mathbb{N}$ = {1, 2, 3, ...}           Natural numbers
$\mathbb{Z}$ = {..., -2, -1, 0, 1, 2, ...} Integers
$\mathbb{Q}$ = {p/q : p,q $\in\mathbb{Z}$, q $\neq$ 0}   Rational numbers (fractions)
$\mathbb{R}$ = all real numbers           (includes irrationals like $\sqrt{2}$, $\pi$)
$\mathbb{C}$ = {a + bi : a,b $\in\mathbb{R}$}        Complex numbers (i$^2$ = -1)
```
\end{tcolorbox}

\textbf{Relationship:} $\mathbb{N}\subset\mathbb{Z}\subset\mathbb{Q}\subset\mathbb{R}\subset\mathbb{C}$

\bigskip\noindent\rule{\textwidth}{0.4pt}\bigskip


A \textbf{function} f: A  B assigns exactly one element f(x) $\in$ B to each x $\in$ A.

\begin{itemize}[leftmargin=*]
  \item \textbf{Domain} = set A (inputs)
  \item \textbf{Codomain} = set B (possible outputs)  
  \item \textbf{Range} = {f(x) : x $\in$ A} (actual outputs) $\subseteq$ B
\end{itemize}

\textbf{Types:}
\begin{itemize}[leftmargin=*]
  \item \textbf{Injective (one-to-one):} f(x$_1$) = f(x$_2$) $\Longrightarrow$ x$_1$ = x$_2$ (different inputs  different outputs)
  \item \textbf{Surjective (onto):} every b $\in$ B has some a with f(a) = b
  \item \textbf{Bijective:} both injective AND surjective (perfect pairing)
\end{itemize}

\begin{tcolorbox}[colback=examplebg!30,colframe=exampleborder,title=\textbf{Worked Example},breakable,boxrule=1pt]
\textbf{Example:} f(x) = x$^2$ on $\mathbb{R}\mathbb{R}$ is NOT injective (f(2) = f(-2) = 4) and NOT surjective (no x gives f(x) = -1).
\end{tcolorbox}

\bigskip\noindent\rule{\textwidth}{0.4pt}\bigskip


A \textbf{matrix} is a rectangular array of numbers. An n$\times$n matrix is "square."

\begin{verbatim}
A = | a  b |      2×2 matrix
    | c  d |
\end{verbatim}

\textbf{Operations:}
\begin{itemize}[leftmargin=*]
  \item \textbf{Addition:} add corresponding entries
  \item \textbf{Scalar multiplication:} multiply every entry by a number
  \item \textbf{Matrix multiplication:} (AB)$_{ij}$ = $\Sigma_k$ a$_{ik}$b$_{kj}$ (row $\times$ column)
\end{itemize}

\textbf{Determinant} (2$\times$2): det(A) = ad - bc

\textbf{Determinant} (3$\times$3): expand along first row:
det = a$_{11}$(a$_{22}$a$_{33}$ - a$_{23}$a$_{32}$) - a$_{12}$(a$_{21}$a$_{33}$ - a$_{23}$a$_{31}$) + a$_{13}$(a$_{21}$a$_{32}$ - a$_{22}$a$_{31}$)

\begin{tcolorbox}[colback=warningbg!30,colframe=warningborder,title=\textbf{Important -- Exam Alert},breakable,boxrule=1.5pt]
\textbf{Key facts:}
- det(A) $\neq$ 0 $\Longleftrightarrow$ A is invertible
- A$^{-}^1$ exists only when det(A) $\neq$ 0
- (AB)$^{-}^1$ = B$^{-}^1$A$^{-}^1$ (reverse order!)
\end{tcolorbox}

\textbf{Identity matrix:} I = diag(1,1,...,1). AI = IA = A.

\bigskip\noindent\rule{\textwidth}{0.4pt}\bigskip


A \textbf{vector} in $\mathbb{R}^n$ is an ordered list of n real numbers: v = (v$_1$, v$_2$, ..., v$_n$).

\textbf{Dot product:} u $\cdot$ v = u$_1$v$_1$ + u$_2$v$_2$ + ... + u$_n$v$_n$ (gives a NUMBER)

\textbf{Length (norm):} ||v|| = $\sqrt{v $\cdot$ v}$ = $\sqrt{v$_1^2$ + v$_2^2$ + ... + v$_n^2$}$

\textbf{Orthogonal:} u $\perp$ v means u $\cdot$ v = 0 (perpendicular)

\textbf{Linear combination:} c$_1$v$_1$ + c$_2$v$_2$ + ... + c$_k$v$_k$ (scalars $\times$ vectors, added up)

\textbf{Linearly independent:} no vector can be written as a combination of the others.

\textbf{Basis:} a set of linearly independent vectors that span the whole space.

$\mathbb{R}^n$ has dimension n. Standard basis for $\mathbb{R}^3$: e$_1$ = (1,0,0), e$_2$ = (0,1,0), e$_3$ = (0,0,1).

\bigskip\noindent\rule{\textwidth}{0.4pt}\bigskip


\subsection{Limits}

lim(xa) f(x) = L means f(x) gets arbitrarily close to L as x approaches a.

\begin{tcolorbox}[colback=warningbg!30,colframe=warningborder,title=\textbf{Important -- Exam Alert},breakable,boxrule=1.5pt]
\textbf{Key limits to remember:}
- lim(x0) sin(x)/x = 1
- lim(x$\infty$) (1 + 1/x)$^x$ = e $\approx$ 2.718
- lim(x0) (e$^x$ - 1)/x = 1
\end{tcolorbox}

\subsection{Derivatives}

f'(x) = lim(h0) [f(x+h) - f(x)] / h

\textbf{Must-know derivatives:}
\begin{verbatim}
d/dx [xⁿ] = nxⁿ⁻¹           d/dx [eˣ] = eˣ
d/dx [sin x] = cos x         d/dx [cos x] = -sin x
d/dx [ln x] = 1/x            d/dx [tan x] = sec²x
\end{verbatim}

\textbf{Chain rule:} d/dx [f(g(x))] = f'(g(x)) $\cdot$ g'(x)

\textbf{Product rule:} (fg)' = f'g + fg'

\subsection{Integration}

$\int$ is the "reverse" of differentiation.

\textbf{Must-know integrals:}
\begin{verbatim}
∫ xⁿ dx = xⁿ⁺¹/(n+1) + C    (n ≠ -1)
∫ 1/x dx = ln|x| + C
∫ eˣ dx = eˣ + C
∫ sin x dx = -cos x + C
∫ cos x dx = sin x + C
\end{verbatim}

\textbf{Fundamental Theorem:} $\int$ₐᵇ f(x)dx = F(b) - F(a) where F' = f.

\textbf{Integration by parts:} $\int$ u dv = uv - $\int$ v du

\bigskip\noindent\rule{\textwidth}{0.4pt}\bigskip


\textbf{Sequence:} ordered list a$_1$, a$_2$, a$_3$, ... (one number for each n $\in\mathbb{N}$)

\textbf{Converges:} a$_n$  L means a$_n$ gets arbitrarily close to L as n  $\infty$

\textbf{Series:} $\Sigma_n$₌$_1$^$\infty$ a$_n$ = a$_1$ + a$_2$ + a$_3$ + ... (infinite sum)

\textbf{Converges} if the partial sums S$_n$ = a$_1$ + ... + a$_n$ approach a finite limit.

\begin{tcolorbox}[colback=warningbg!30,colframe=warningborder,title=\textbf{Important -- Exam Alert},breakable,boxrule=1.5pt]
\textbf{Key tests:}
- \textbf{Geometric:} $\Sigma$ r$^n$ converges iff |r| < 1, sum = 1/(1-r)
- \textbf{p-series:} $\Sigma$ 1/nᵖ converges iff p > 1
- \textbf{Comparison:} if 0 $\leq$ a$_n\leq$ b$_n$ and $\Sigma$b$_n$ converges, then $\Sigma$a$_n$ converges
- \textbf{Ratio test:} if lim |a$_n$₊$_1$/a$_n$| < 1, converges; > 1, diverges
\end{tcolorbox}

\bigskip\noindent\rule{\textwidth}{0.4pt}\bigskip


\textbf{i = $\sqrt{-1}$}, so i$^2$ = -1.

A complex number: z = a + bi where a = Re(z), b = Im(z).

\textbf{Conjugate:} z̄ = a - bi

\textbf{Modulus:} |z| = $\sqrt{a$^2$ + b$^2$}$

\textbf{Polar form:} z = r(cos$\theta$ + i sin$\theta$) = re^(i$\theta$) where r = |z|, $\theta$ = arg(z)

\textbf{Euler's formula:} e^(i$\theta$) = cos$\theta$ + i sin$\theta$

\textbf{De Moivre:} (cos$\theta$ + i sin$\theta$)$^n$ = cos(n$\theta$) + i sin(n$\theta$)

\begin{tcolorbox}[colback=warningbg!30,colframe=warningborder,title=\textbf{Important -- Exam Alert},breakable,boxrule=1.5pt]
\textbf{Key:} Every polynomial of degree n has exactly n roots in $\mathbb{C}$ (counting multiplicity).
\end{tcolorbox}

\bigskip\noindent\rule{\textwidth}{0.4pt}\bigskip


In MSc Maths, you'll write proofs constantly. Here are the main methods:

\subsection{Direct Proof}
Assume P, derive Q step by step.
\textit{Example:} "If n is even, then n$^2$ is even." Let n = 2k. Then n$^2$ = 4k$^2$ = 2(2k$^2$), which is even. $\blacksquare$

\subsection{Proof by Contradiction}
Assume the statement is FALSE, derive a contradiction.
\textit{Example:} "$\sqrt{2}$ is irrational." Assume $\sqrt{2}$ = p/q (fully reduced). Then 2q$^2$ = p$^2$, so p is even. Write p = 2k: 2q$^2$ = 4k$^2$, q$^2$ = 2k$^2$, so q is even. But p,q both even contradicts "fully reduced." $\blacksquare$

\subsection{Proof by Induction}
\begin{itemize}
  \item \textbf{Base case:} Prove for n = 1.
  \item \textbf{Inductive step:} Assume true for n = k, prove for n = k+1.
\textit{Example:} 1 + 2 + ... + n = n(n+1)/2. Base: n=1: 1 = 1$\cdot$$\frac{2}{2}$ . Step: Assume for k. Then 1+...+k+(k+1) = k(k+1)/2 + (k+1) = (k+1)(k+2)/2 . $\blacksquare$
\end{itemize}

\subsection{Proof by Contrapositive}
Instead of "P $\Longrightarrow$ Q", prove "not Q $\Longrightarrow$ not P" (logically equivalent).
\textit{Example:} "If n$^2$ is even, then n is even." Contrapositive: "If n is odd, then n$^2$ is odd." Let n = 2k+1. n$^2$ = 4k$^2$+4k+1 = 2(2k$^2$+2k)+1, which is odd.  $\blacksquare$

\bigskip\noindent\rule{\textwidth}{0.4pt}\bigskip


\begin{center}
\begin{tabular}{|l|l|l|}
\hline
\textbf{Symbol} & \textbf{Meaning} & \textbf{Example} \\
\hline
$\forall$ & "for all" & $\forall$x $\in\mathbb{R}$ means "for every real x" \\
\hline
$\exists$ & "there exists" & $\exists$x such that x$^2$ = 2 \\
\hline
$\Longrightarrow$ & "implies" & P $\Longrightarrow$ Q means "if P then Q" \\
\hline
$\Longleftrightarrow$ & "if and only if" & P $\Longleftrightarrow$ Q means P and Q are equivalent \\
\hline
$\blacksquare$ or □ & "end of proof" & Written at end of every proof \\
\hline
:= & "defined as" & f(x) := x$^2$ means "f is defined as x$^2$" \\
\hline
$\Sigma$ & summation & $\Sigma_i$₌$_1^n$ i = 1+2+...+n \\
\hline
$\Pi$ & product & $\Pi_i$₌$_1^n$ i = 1$\cdot$2$\cdot$...$\cdot$n = n! \\
\hline
$\circ$ & composition & (f$\circ$g)(x) = f(g(x)) \\
\hline
$\approx$ & approximately & $\pi\approx$ 3.14159 \\
\hline
dim & dimension & dim($\mathbb{R}^3$) = 3 \\
\hline
ker & kernel (null space) & ker(T) = {v : T(v) = 0} \\
\hline
Im & image (range) & Im(T) = {T(v) : v $\in$ V} \\
\hline
sup & supremum (least upper bound) & sup{1/n : n $\in\mathbb{N}$} = 1 \\
\hline
inf & infimum (greatest lower bound) & inf{1/n : n $\in\mathbb{N}$} = 0 \\
\hline
\end{tabular}
\end{center}

\bigskip\noindent\rule{\textwidth}{0.4pt}\bigskip


\begin{itemize}
  \item \textbf{Read this prerequisites chapter first} — make sure you're comfortable with every concept
  \item \textbf{For each unit:} Read the explanation  Study the example  Try to reproduce the proof without looking
  \item \textbf{Don't skip proofs} — university exams will ask "State and Prove"
  \item \textbf{Use pen and paper} — mathematics is NOT a spectator sport. Write as you read.
  \item \textbf{If stuck on a concept:} Re-read the prerequisites section for that topic, then return.
\end{itemize}

\begin{quote}
\textbf{Remember:} Every mathematician was once where you are. The difference between a student and a professor is simply years of consistent practice. You've got this! 
\end{quote}

\newpage

\chapter{Study Roadmap \& CSIR NET Guide}


\begin{quote}
\textbf{University:} Manipal University Jaipur (Online/Distance Education)
\textbf{Goal:} Clear semester exams with distinction + Qualify CSIR NET for Assistant Professorship
\textbf{Prepared for:} A learner with Class 12 Mathematics background stepping into postgraduate mathematics
\end{quote}

\bigskip\noindent\rule{\textwidth}{0.4pt}\bigskip

\subsection{ Your Path to Assistant Professorship}

\begin{verbatim}
  Class 12 Math ──▶ MSc Sem 1 ──▶ Sem 2,3,4 ──▶ CSIR NET ──▶ Assistant Professor
  (You are here)    (BUILD        (ADVANCE)      (QUALIFY)     (YOUR GOAL!)
                     FOUNDATIONS)
\end{verbatim}

\begin{quote}
\textbf{⚠️ IMPORTANT — Key Clarification:}
To become an Assistant Professor in India, you need to clear the \textbf{CSIR NET} (Council of Scientific and Industrial Research National Eligibility Test) in \textbf{Mathematical Sciences} — not UGC NET (which doesn't cover Mathematics). The CSIR NET is conducted by NTA (National Testing Agency).
\end{quote}

\bigskip\noindent\rule{\textwidth}{0.4pt}\bigskip

\subsection{ Semester 1 Subjects at a Glance}

\begin{center}
\begin{tabular}{|l|l|l|l|l|}
\hline
\textbf{#} & \textbf{Subject} & \textbf{Priority for CSIR NET} & \textbf{Difficulty Level} & \textbf{Study Order} \\
\hline
1 & [Advanced Linear Algebra](./01_Advanced_Linear_Algebra.md) & ⭐⭐⭐⭐⭐ (Very High) &  Medium & \textbf{Study 1st} \\
\hline
2 & [Mathematical Analysis](./02_Mathematical_Analysis.md) & ⭐⭐⭐⭐⭐ (Very High) &  Medium-Hard & \textbf{Study 2nd} \\
\hline
3 & [Topology-1](./03_Topology.md) & ⭐⭐⭐⭐ (High) &  Hard & \textbf{Study 3rd} \\
\hline
4 & [Advanced Complex Analysis](./04_Advanced_Complex_Analysis.md) & ⭐⭐⭐⭐ (High) &  Hard & \textbf{Study 4th} \\
\hline
5 & [Advanced Differential Equations](./05_Advanced_Differential_Equations.md) & ⭐⭐⭐⭐ (High) &  Medium-Hard & \textbf{Study 5th} \\
\hline
6 & [Dynamics of a Rigid Body](./06_Dynamics_of_Rigid_Body.md) & ⭐⭐⭐ (Moderate) &  Medium & \textbf{Study 6th} \\
\hline
\end{tabular}
\end{center}

\begin{quote}
\textbf{ TIP — Why this study order?}
Linear Algebra and Real Analysis form the \textit{language} of all other subjects. Topology builds on Analysis. Complex Analysis uses both. Differential Equations use Linear Algebra. Dynamics is the most applied and uses everything.
\end{quote}

\bigskip\noindent\rule{\textwidth}{0.4pt}\bigskip

\subsection{ CSIR NET Exam Pattern (Mathematical Sciences)}

\begin{verbatim}
┌──────────────────────────────────────────────────────────────┐
│                CSIR NET Mathematical Sciences                 │
├──────────────────┬───────────────────┬───────────────────────┤
│   Part A         │    Part B         │    Part C              │
│   General        │    Subject MCQs   │    Advanced MSQs       │
│   Aptitude       │                   │                        │
│   20 Qs · 30 M   │    25 Qs · 75 M   │    20 Qs · 95 M       │
├──────────────────┼───────────────────┼───────────────────────┤
│ Logical Reasoning│ Core Math Topics  │ Multi-Select Questions │
│ Quant Aptitude   │ Standard Problems │ Higher Difficulty      │
│ Data Interp.     │ Single Correct    │ Negative Marking       │
└──────────────────┴───────────────────┴───────────────────────┘
\end{verbatim}

\begin{center}
\begin{tabular}{|l|l|l|l|l|}
\hline
\textbf{Part} & \textbf{Questions} & \textbf{Max Marks} & \textbf{Negative Marking} & \textbf{Strategy} \\
\hline
\textbf{A} & 20 (attempt 15) & 30 & -0.5 per wrong & Must score 18+ \\
\hline
\textbf{B} & 25 (attempt 20) & 75 & -0.75 per wrong & Core strength — target 50+ \\
\hline
\textbf{C} & 20 (attempt 10) & 95 & -1.0 per wrong & Selective, high-reward \\
\hline
\end{tabular}
\end{center}

\textbf{Total: 200 marks | Duration: 3 hours | Qualifying: ~45-55% (varies by cutoff)}

\bigskip\noindent\rule{\textwidth}{0.4pt}\bigskip

\subsection{ Which Semester 1 Topics Map to CSIR NET?}

\begin{center}
\begin{tabular}{|l|l|l|}
\hline
\textbf{CSIR NET Topic} & \textbf{Your Sem 1 Subject} & \textbf{Units to Focus} \\
\hline
Linear Algebra (Part B+C) & Advanced Linear Algebra & All units — \textbf{highest weightage} \\
\hline
Real Analysis (Part B+C) & Mathematical Analysis & Units 1-5, 6-10 \\
\hline
Topology & Topology-1 & Units 1-4, 5-8 \\
\hline
Complex Analysis (Part B) & Advanced Complex Analysis & Units 1-4 (basics overlap) \\
\hline
ODE (Part B) & Advanced Differential Equations & Units 1-4, 7 \\
\hline
Classical Mechanics & Dynamics of a Rigid Body & Units 3-4 (Lagrangian overlap) \\
\hline
\end{tabular}
\end{center}

\bigskip\noindent\rule{\textwidth}{0.4pt}\bigskip

\subsection{ Suggested 16-Week Study Planner}

\begin{verbatim}
Week  1-3:   Advanced Linear Algebra (Units 1-5)
Week  4-5:   Advanced Linear Algebra (Units 6-10) + Review
Week  6-8:   Mathematical Analysis (Units 1-5)
Week  9-10:  Mathematical Analysis (Units 6-10) + Review
Week 11-12:  Topology-1 (Units 1-5)
Week 13:     Topology-1 (Units 6-10)
Week 14:     Advanced Complex Analysis (Units 1-5)
Week 15:     Adv. Complex Analysis (Units 6-10) + Adv. Diff. Eq. (Units 1-5)
Week 16:     Adv. Diff. Eq. (Units 6-10) + Dynamics (all) + Final Revision
\end{verbatim}

\bigskip\noindent\rule{\textwidth}{0.4pt}\bigskip

\subsection{ Exam Strategy}

\subsubsection{For Manipal Semester Exams}
\begin{itemize}[leftmargin=*]
  \item Focus on \textbf{all units equally} — university exams test breadth
  \item Practice \textbf{worked examples} from each unit — exam questions are variations
  \item Memorize \textbf{key theorems and their statements} — they ask "State and prove"
\end{itemize}

\subsubsection{For CSIR NET}
\begin{itemize}[leftmargin=*]
  \item Focus on \textbf{problem-solving speed} — it's MCQ-based
  \item Practice \textbf{previous year questions} extensively
  \item Build \textbf{conceptual clarity} over rote memorization
  \item Master \textbf{Linear Algebra + Real Analysis} first — they carry ~40% weightage combined
\end{itemize}

\bigskip\noindent\rule{\textwidth}{0.4pt}\bigskip

\subsection{ How to Use These Notes}

\begin{itemize}
  \item \textbf{Read sequentially within each subject} — each unit builds on the previous
  \item \textbf{Work through every example} — don't just read, compute
  \item \textbf{Attempt the practice problems} before looking at solutions
  \item \textbf{Use the mind maps} at the start of each unit for revision
  \item \textbf{Pay special attention to  CSIR NET markers} throughout the notes
  \item \textbf{Revisit the  Key Takeaway boxes} for quick revision before exams
\end{itemize}

\begin{quote}
\textbf{ CAUTION:}
\textbf{Do NOT skip the basics!} Even though these are "advanced" courses, every concept is built from fundamentals explained in these notes. Skipping ahead will leave gaps that compound painfully later.
\end{quote}

\bigskip\noindent\rule{\textwidth}{0.4pt}\bigskip

\textit{Let's begin your journey to becoming Professor Bhanu! Start with [Advanced Linear Algebra ](./01_Advanced_Linear_Algebra.md)}

\newpage

\chapter{Topic Hit-Rate Priority Analysis}


\begin{quote}
\textbf{How to read:} ⚡ = High ROI (study less, score more) |  = Medium ROI |  = Low ROI (needs heavy study for few marks)
\end{quote}

\bigskip\noindent\rule{\textwidth}{0.4pt}\bigskip

\subsection{OVERALL SUBJECT ROI FOR CSIR NET}

\begin{center}
\begin{tabular}{|l|l|l|l|l|}
\hline
\textbf{Priority} & \textbf{Subject} & \textbf{NET Questions} & \textbf{Study Hours Needed} & \textbf{ROI Rating} \\
\hline
1 & \textbf{Linear Algebra} & 4-5 per exam & ~80 hrs & ⚡⚡⚡ BEST ROI \\
\hline
2 & \textbf{Real Analysis} & 4-5 per exam & ~100 hrs & ⚡⚡⚡ BEST ROI \\
\hline
3 & \textbf{Topology} & 2-3 per exam & ~70 hrs & ⚡⚡ GOOD ROI \\
\hline
4 & \textbf{ODE} & 2-3 per exam & ~60 hrs & ⚡⚡ GOOD ROI \\
\hline
5 & \textbf{Complex Analysis} & 2-3 per exam & ~80 hrs & ⚡ MODERATE \\
\hline
6 & \textbf{Dynamics} & 0-1 per exam & ~50 hrs &  LOW (semester only) \\
\hline
\end{tabular}
\end{center}

\bigskip\noindent\rule{\textwidth}{0.4pt}\bigskip

\subsection{ ADVANCED LINEAR ALGEBRA — Unit-wise Hit Rate}

\begin{center}
\begin{tabular}{|l|l|l|l|l|l|}
\hline
\textbf{Unit} & \textbf{Topic} & \textbf{CSIR NET} & \textbf{Semester Exam} & \textbf{Hit Rate} & \textbf{Verdict} \\
\hline
7 & \textbf{Eigenvalues, Diagonalization} & ★★★★★ & ★★★★★ & ⚡⚡⚡ & \textbf{\#1 PRIORITY — Study this FIRST} \\
\hline
9 & \textbf{Cayley-Hamilton, Jordan Form} & ★★★★★ & ★★★★★ & ⚡⚡⚡ & \textbf{\#2 — Asked every single exam} \\
\hline
4-5 & \textbf{Inner Product, Gram-Schmidt} & ★★★★☆ & ★★★★☆ & ⚡⚡⚡ & \textbf{\#3 — Computation-heavy, easy marks} \\
\hline
1-2 & \textbf{Linear Transformations, Matrices} & ★★★★☆ & ★★★★★ & ⚡⚡ & Foundation — must know \\
\hline
10 & \textbf{Bilinear/Quadratic Forms} & ★★★☆☆ & ★★★★☆ & ⚡⚡ & Positive definite classification \\
\hline
3 & \textbf{Dual Spaces, Annihilators} & ★★☆☆☆ & ★★★☆☆ & ⚡ & Theory-heavy, fewer direct Qs \\
\hline
6 & \textbf{Inner Product Isomorphism} & ★★☆☆☆ & ★★★☆☆ & ⚡ & Riesz theorem — know statement \\
\hline
8 & \textbf{Invariant Subspaces} & ★★★☆☆ & ★★★☆☆ & ⚡⚡ & Connects to Jordan Form \\
\hline
\end{tabular}
\end{center}

\subsubsection{ LINEAR ALGEBRA: If you study ONLY 3 topics, study these:}
\begin{itemize}
  \item \textbf{Eigenvalues + Diagonalization} (guaranteed 2-3 questions)
  \item \textbf{Cayley-Hamilton + Jordan Form} (guaranteed 1-2 questions)
  \item \textbf{Gram-Schmidt + Orthogonal matrices} (guaranteed 1 question)
\end{itemize}

\bigskip\noindent\rule{\textwidth}{0.4pt}\bigskip

\subsection{ MATHEMATICAL ANALYSIS — Unit-wise Hit Rate}

\begin{center}
\begin{tabular}{|l|l|l|l|l|l|}
\hline
\textbf{Unit} & \textbf{Topic} & \textbf{CSIR NET} & \textbf{Semester Exam} & \textbf{Hit Rate} & \textbf{Verdict} \\
\hline
3-4 & \textbf{Uniform Convergence} & ★★★★★ & ★★★★★ & ⚡⚡⚡ & \textbf{\#1 — Tested EVERY exam} \\
\hline
9 & \textbf{Extrema, Lagrange Multipliers} & ★★★★★ & ★★★★☆ & ⚡⚡⚡ & \textbf{\#2 — Easy computation marks} \\
\hline
10 & \textbf{Jacobian} & ★★★★☆ & ★★★★☆ & ⚡⚡⚡ & \textbf{\#3 — Direct formula application} \\
\hline
6-7 & \textbf{Multivariable Calculus} & ★★★★☆ & ★★★★★ & ⚡⚡ & Differentiability, chain rule \\
\hline
8 & \textbf{Implicit/Inverse Function Thm} & ★★★★☆ & ★★★★☆ & ⚡⚡ & Statement + application \\
\hline
1-2 & \textbf{Riemann-Stieltjes Integral} & ★★★☆☆ & ★★★★★ & ⚡ & Semester important, less in NET \\
\hline
5 & \textbf{Stone-Weierstrass, Arzelà-Ascoli} & ★★★☆☆ & ★★★☆☆ &  & Know statements, rarely computed \\
\hline
\end{tabular}
\end{center}

\subsubsection{ ANALYSIS: If you study ONLY 3 topics:}
\begin{itemize}
  \item \textbf{Uniform vs Pointwise convergence + M-test} (2-3 questions guaranteed)
  \item \textbf{Lagrange multipliers + extrema} (1-2 questions)
  \item \textbf{Differentiability of multivariable functions} (1-2 questions)
\end{itemize}

\bigskip\noindent\rule{\textwidth}{0.4pt}\bigskip

\subsection{ TOPOLOGY — Unit-wise Hit Rate}

\begin{center}
\begin{tabular}{|l|l|l|l|l|l|}
\hline
\textbf{Unit} & \textbf{Topic} & \textbf{CSIR NET} & \textbf{Semester Exam} & \textbf{Hit Rate} & \textbf{Verdict} \\
\hline
7-8 & \textbf{Compactness + Heine-Borel} & ★★★★★ & ★★★★★ & ⚡⚡⚡ & \textbf{\#1 — "Is it compact?" = free marks} \\
\hline
6 & \textbf{Connectedness} & ★★★★★ & ★★★★★ & ⚡⚡⚡ & \textbf{\#2 — "Is it connected?" = free marks} \\
\hline
2 & \textbf{Open/Closed/Closure/Interior} & ★★★★☆ & ★★★★★ & ⚡⚡ & Foundation for everything \\
\hline
4 & \textbf{Basis, Product, Subspace Topology} & ★★★☆☆ & ★★★★☆ & ⚡⚡ & Constructive questions \\
\hline
9 & \textbf{Separation Axioms} & ★★★☆☆ & ★★★★☆ & ⚡ & Know T$_0$-T₄ hierarchy \\
\hline
5 & \textbf{Continuity, Homeomorphisms} & ★★★☆☆ & ★★★★☆ & ⚡⚡ & Topological definition \\
\hline
10 & \textbf{Separable, Second Countable} & ★★☆☆☆ & ★★★☆☆ &  & Rare in NET, semester appears \\
\hline
1,3 & \textbf{Basic Concepts, Boundary} & ★★☆☆☆ & ★★★★☆ & ⚡ & Easy but foundational \\
\hline
\end{tabular}
\end{center}

\subsubsection{ TOPOLOGY: If you study ONLY 3 topics:}
\begin{itemize}
  \item \textbf{Compactness} — "closed and bounded in $\mathbb{R}^n$" + Heine-Borel + Tychonoff
  \item \textbf{Connectedness} — intervals, path-connected, IVT
  \item \textbf{Standard counterexamples} — cofinite topology, Topologist's sine curve, etc.
\end{itemize}

\bigskip\noindent\rule{\textwidth}{0.4pt}\bigskip

\subsection{ COMPLEX ANALYSIS — Unit-wise Hit Rate}

\begin{center}
\begin{tabular}{|l|l|l|l|l|l|}
\hline
\textbf{Unit} & \textbf{Topic} & \textbf{CSIR NET} & \textbf{Semester Exam} & \textbf{Hit Rate} & \textbf{Verdict} \\
\hline
1 & \textbf{Entire Functions + Liouville} & ★★★★★ & ★★★★★ & ⚡⚡⚡ & \textbf{\#1 — Liouville asked EVERY exam} \\
\hline
4-5 & \textbf{Harmonic Functions + Maximum Principle} & ★★★★☆ & ★★★★☆ & ⚡⚡ & \textbf{\#2 — Harmonic conjugate} \\
\hline
9 & \textbf{Picard Theorems} & ★★★★☆ & ★★★★☆ & ⚡⚡⚡ & \textbf{\#3 — "How many values omitted?"} \\
\hline
3 & \textbf{Analytic Continuation} & ★★★☆☆ & ★★★★☆ & ⚡ & Identity theorem is key \\
\hline
6 & \textbf{Canonical Products, Jensen} & ★★☆☆☆ & ★★★★☆ &  & Semester-focused \\
\hline
7 & \textbf{Hadamard Three Circles} & ★★☆☆☆ & ★★★☆☆ &  & Statement only for NET \\
\hline
8 & \textbf{Hadamard Factorization} & ★★☆☆☆ & ★★★★☆ &  & Semester important \\
\hline
2 & \textbf{Gamma, Zeta Functions} & ★★☆☆☆ & ★★★☆☆ &  & Properties only \\
\hline
10 & \textbf{Univalent Functions} & ★☆☆☆☆ & ★★★☆☆ &  & Rarely in NET \\
\hline
\end{tabular}
\end{center}

\subsubsection{ COMPLEX ANALYSIS: If you study ONLY 3 topics:}
\begin{itemize}
  \item \textbf{Liouville's Theorem + applications} (entire + bounded = constant)
  \item \textbf{Picard's Theorem} (entire function omits at most 1 value)
  \item \textbf{Harmonic functions + conjugate finding} (Cauchy-Riemann)
\end{itemize}

\bigskip\noindent\rule{\textwidth}{0.4pt}\bigskip

\subsection{⚙️ DIFFERENTIAL EQUATIONS — Unit-wise Hit Rate}

\begin{center}
\begin{tabular}{|l|l|l|l|l|l|}
\hline
\textbf{Unit} & \textbf{Topic} & \textbf{CSIR NET} & \textbf{Semester Exam} & \textbf{Hit Rate} & \textbf{Verdict} \\
\hline
7 & \textbf{Critical Points + Lyapunov} & ★★★★★ & ★★★★★ & ⚡⚡⚡ & \textbf{\#1 — Classify stability = easy marks} \\
\hline
2 & \textbf{Picard-Lindelöf (Existence/Uniqueness)} & ★★★★★ & ★★★★★ & ⚡⚡⚡ & \textbf{\#2 — "Does unique soln exist?"} \\
\hline
4 & \textbf{Matrix Exponential} & ★★★★☆ & ★★★★☆ & ⚡⚡ & \textbf{\#3 — Direct computation} \\
\hline
6,10 & \textbf{Poincaré-Bendixson} & ★★★☆☆ & ★★★★☆ & ⚡⚡ & Statement + application \\
\hline
1 & \textbf{$\varepsilon$-Approximate Solutions} & ★☆☆☆☆ & ★★★★☆ &  & Semester only \\
\hline
3 & \textbf{Systems of DEs} & ★★★☆☆ & ★★★★☆ & ⚡⚡ & Foundation for matrix method \\
\hline
5 & \textbf{Nonlinear Systems} & ★★★☆☆ & ★★★☆☆ & ⚡ & Phase plane analysis \\
\hline
8 & \textbf{Paths of Linear Systems} & ★★☆☆☆ & ★★★☆☆ &  & Covered via Unit 7 \\
\hline
9 & \textbf{Parameter Dependence} & ★★☆☆☆ & ★★★☆☆ &  & Bifurcation basics \\
\hline
\end{tabular}
\end{center}

\subsubsection{ ODE: If you study ONLY 3 topics:}
\begin{itemize}
  \item \textbf{Critical point classification} (eigenvalues  node/saddle/spiral/center)
  \item \textbf{Existence & Uniqueness} (Lipschitz condition, Picard iteration)
  \item \textbf{Lyapunov function method} (V > 0, V̇ < 0  stable)
\end{itemize}

\bigskip\noindent\rule{\textwidth}{0.4pt}\bigskip

\subsection{ DYNAMICS — Unit-wise Hit Rate}

\begin{center}
\begin{tabular}{|l|l|l|l|l|l|}
\hline
\textbf{Unit} & \textbf{Topic} & \textbf{CSIR NET} & \textbf{Semester Exam} & \textbf{Hit Rate} & \textbf{Verdict} \\
\hline
1-2 & \textbf{Moments of Inertia + Theorems} & ★★☆☆☆ & ★★★★★ & ⚡⚡ & \textbf{Semester \#1 — direct formulas} \\
\hline
3 & \textbf{D'Alembert + Lagrangian} & ★★★☆☆ & ★★★★★ & ⚡⚡ & \textbf{Semester \#2 — NET overlap} \\
\hline
4 & \textbf{Angular Momentum + Euler Eqs} & ★★☆☆☆ & ★★★★★ & ⚡⚡ & \textbf{Semester \#3} \\
\hline
7 & \textbf{Rolling Motion} & ★☆☆☆☆ & ★★★★☆ & ⚡ & Computation-based \\
\hline
10 & \textbf{Conservation Laws} & ★★☆☆☆ & ★★★★☆ & ⚡⚡ & Energy/momentum problems \\
\hline
5,9 & \textbf{Impulsive Forces} & ★☆☆☆☆ & ★★★☆☆ &  & Semester only \\
\hline
6 & \textbf{Fixed Axis Rotation} & ★☆☆☆☆ & ★★★★☆ & ⚡ & Compound pendulum \\
\hline
8 & \textbf{2D Motion} & ★☆☆☆☆ & ★★★☆☆ &  & Application problems \\
\hline
\end{tabular}
\end{center}

\subsubsection{ DYNAMICS: For semester exam focus on:}
\begin{itemize}
  \item \textbf{Moment of inertia calculations + parallel/perpendicular axis theorems}
  \item \textbf{D'Alembert's principle + Lagrangian}
  \item \textbf{Rolling motion + conservation laws}
\end{itemize}

\bigskip\noindent\rule{\textwidth}{0.4pt}\bigskip

\subsection{ ULTIMATE PRIORITY MATRIX}

\subsubsection{MUST study (appear in EVERY exam):}
\begin{verbatim}
1. Eigenvalues, Diagonalization, Jordan Form     [Linear Algebra]
2. Uniform Convergence + M-test                   [Analysis]
3. Compactness + Heine-Borel                      [Topology]
4. Liouville's Theorem                            [Complex Analysis]
5. Critical Point Classification + Lyapunov       [ODE]
6. Cayley-Hamilton Theorem                        [Linear Algebra]
7. Picard-Lindelöf Existence/Uniqueness           [ODE]
8. Picard's Little Theorem                        [Complex Analysis]
9. Connectedness + IVT                            [Topology]
10. Lagrange Multipliers                          [Analysis]
\end{verbatim}

\subsubsection{SHOULD study (appear in 70% of exams):}
\begin{verbatim}
11. Gram-Schmidt Process                          [Linear Algebra]
12. Implicit/Inverse Function Theorem             [Analysis]
13. Matrix Exponential                            [ODE]
14. Separation Axioms (T₀-T₄)                    [Topology]
15. Harmonic Functions + Conjugate                [Complex Analysis]
16. Bilinear/Quadratic Forms                      [Linear Algebra]
17. Riemann-Stieltjes Integral                    [Analysis]
18. Poincaré-Bendixson Theorem                    [ODE]
\end{verbatim}

\subsubsection{COULD skip for NET (but study for semester):}
\begin{verbatim}
19. ε-Approximate Solutions                       [ODE]
20. Univalent Function Theory                     [Complex Analysis]
21. Impulsive Forces (2D)                         [Dynamics]
22. Dual Spaces (deep theory)                     [Linear Algebra]
23. Hadamard Three Circles (proof)                [Complex Analysis]
\end{verbatim}

\newpage

\chapter{Advanced Linear Algebra}


\begin{quote}
\textbf{CSIR NET Priority: ⭐⭐⭐⭐⭐ | This is the HIGHEST weightage subject in CSIR NET}
\end{quote}

\bigskip\noindent\rule{\textwidth}{0.4pt}\bigskip

\subsection{️ Subject Mind Map}

\begin{verbatim}
                    ADVANCED LINEAR ALGEBRA
                           │
         ┌─────────┬───────┴────────┬──────────┐
         │         │                │          │
    Transforms   Spaces          Forms     Canonical
         │         │                │          │
    ┌────┴────┐  ┌─┴──┐        ┌───┴───┐   ┌──┴──┐
    Linear  Dual Inner  Orth.  Bilinear  Jordan Cayley
    Trans.  Basis Prod.  Trans. Quadratic Form  Hamilton
\end{verbatim}

\bigskip\noindent\rule{\textwidth}{0.4pt}\bigskip

\section{Unit 1: Linear Transformations}

\begin{prereqbox}
 \textbf{Quick Refresher Before Starting:}
\textbf{Matrix:} A rectangular grid of numbers. A 2$\times$2 matrix looks like:
```
A = | 2  3 |    We write A = [[2,3],[1,4]]
| 1  4 |
```
\textbf{Matrix multiplication:} Row $\times$ Column. For AB, take each row of A, multiply element-by-element with each column of B, and add up.
```
| 2 3 | $\times$ | 1 | = | 2$\times$1 + 3$\times$0 | = | 2 |
| 1 4 |   | 0 |   | 1$\times$1 + 4$\times$0 |   | 1 |
```
\textbf{Vector:} An ordered list of numbers, like v = (3, -1, 2) in $\mathbb{R}^3$.
\textbf{Linear combination:} c$_1$v$_1$ + c$_2$v$_2$ means scaling and adding vectors.
\textbf{System Ax = b:} Finding x means solving multiple equations simultaneously.
\end{prereqbox}



\subsection{Prerequisites from Class 12}
You know matrices and solving equations. Now we go deeper — a \textbf{linear transformation} is the \textit{abstract idea} behind matrix multiplication.

\subsection{What is a Linear Transformation?}

\begin{tcolorbox}[colback=lightblue!30,colframe=defgreen,title=\textbf{Definition},breakable,boxrule=1.5pt]
\textbf{Definition:} Let V and W be vector spaces over a field F. A function T: V  W is a \textbf{linear transformation} if for all vectors u, v $\in$ V and all scalars c $\in$ F:
\end{tcolorbox}

\begin{itemize}
  \item \textbf{T(u + v) = T(u) + T(v)} (preserves addition)
  \item \textbf{T(cv) = cT(v)} (preserves scalar multiplication)
\end{itemize}

\textbf{In plain English:} A linear transformation is a function between vector spaces that "plays nice" with addition and scaling. No bending, no twisting — just stretching, rotating, reflecting, or projecting.

\subsubsection{ Memory Aid}
Think of a linear transformation as a "well-behaved machine": if you double the input, the output doubles. If you add two inputs, the output is the sum of individual outputs.

\subsubsection{Example 1: Rotation in $\mathbb{R}^2$}
The rotation by angle $\theta$ counterclockwise is T: $\mathbb{R}^2\mathbb{R}^2$ defined by:

\begin{verbatim}
T(x, y) = (x cos θ − y sin θ, x sin θ + y cos θ)
\end{verbatim}

\textbf{Verification of linearity:}
Let u = (x$_1$, y$_1$), v = (x$_2$, y$_2$)

T(u + v) = T(x$_1$+x$_2$, y$_1$+y$_2$)
= ((x$_1$+x$_2$)cos$\theta$ − (y$_1$+y$_2$)sin$\theta$, (x$_1$+x$_2$)sin$\theta$ + (y$_1$+y$_2$)cos$\theta$)
= (x$_1$cos$\theta$ − y$_1$sin$\theta$, x$_1$sin$\theta$ + y$_1$cos$\theta$) + (x$_2$cos$\theta$ − y$_2$sin$\theta$, x$_2$sin$\theta$ + y$_2$cos$\theta$)
= T(u) + T(v) 

Similarly, T(cu) = cT(u) 

\subsubsection{Example 2: Differentiation}
D: P$_n$($\mathbb{R}$)  P$_n$₋$_1$($\mathbb{R}$), where D(p) = p' (the derivative)

\begin{verbatim}
D(3x² + 2x + 1) = 6x + 2
\end{verbatim}

This IS linear because: d/dx(f + g) = f' + g' and d/dx(cf) = cf'

\subsubsection{Example 3 (Non-example): T(x) = x$^2$ is NOT linear}
T(2) = 4, but 2$\cdot$T(1) = 2$\cdot$1 = 2 $\neq$ 4. So T(cu) $\neq$ cT(u). 

\subsection{Key Concepts}

\subsubsection{Kernel (Null Space)}
\textbf{ker(T) = {v $\in$ V : T(v) = 0}} — all vectors that T "kills" (maps to zero)

\subsubsection{Image (Range)}
\textbf{Im(T) = {T(v) : v $\in$ V}} — all possible outputs of T

\subsubsection{ CSIR NET Critical Theorem: Rank-Nullity Theorem}

\textbf{dim(V) = dim(ker T) + dim(Im T)}

That is: \textbf{dim(domain) = nullity + rank}

\begin{tcolorbox}[colback=examplebg!30,colframe=exampleborder,title=\textbf{Worked Example},breakable,boxrule=1pt]
\textbf{Example:} T: $\mathbb{R}^4\mathbb{R}^3$ with rank(T) = 2. Then nullity = 4 − 2 = 2.
\end{tcolorbox}

\subsubsection{Worked Example (CSIR NET Style)}
\textbf{Q:} Let T: $\mathbb{R}^3\mathbb{R}^3$ be defined by T(x,y,z) = (x+y, y+z, 0). Find rank and nullity.

\textbf{Solution:}
\begin{itemize}[leftmargin=*]
  \item Matrix of T: A = [[1,1,0],[0,1,1],[0,0,0]]
  \item Row reduce: already in echelon form. Rank = 2 (two non-zero rows)
  \item By Rank-Nullity: Nullity = 3 − 2 = 1
  \item ker(T): x+y=0, y+z=0  y=−x, z=x  ker = span{(1,−1,1)}
\end{itemize}

\bigskip\noindent\rule{\textwidth}{0.4pt}\bigskip

\section{Unit 2: Linear Transformations and Vector Spaces}

\begin{prereqbox}
 \textbf{Quick Refresher Before Starting:}
\textbf{Basis:} A set of vectors that (1) spans the space and (2) is linearly independent.
Example: {(1,0), (0,1)} is the standard basis for $\mathbb{R}^2$. Every vector (a,b) = a(1,0) + b(0,1).
\textbf{Dimension:} Number of vectors in a basis. dim($\mathbb{R}^3$) = 3.
\textbf{Linear transformation T:} A function between vector spaces that preserves addition and scaling:
T(u + v) = T(u) + T(v) and T(cv) = cT(v).
Example: T(x,y) = (2x + y, x - y) is linear. T(x,y) = (x$^2$, y) is NOT (because T(2v) $\neq$ 2T(v)).
\end{prereqbox}


\subsection{Matrix Representation of Linear Transformations}

Every linear transformation between finite-dimensional spaces can be represented by a matrix (once you choose bases).

\textbf{Recipe:}
\begin{itemize}
  \item Choose ordered basis B = {v$_1$,...,v$_n$} for V and B' = {w$_1$,...,w$_m$} for W
  \item Compute T(v$_1$), T(v$_2$), ..., T(v$_n$)
  \item Express each T(v$_j$) as a linear combination of w$_1$,...,w$_m$
  \item The coefficients form the columns of the matrix [T]
\end{itemize}

\subsubsection{Example}
Let T: $\mathbb{R}^2\mathbb{R}^2$ be T(x,y) = (2x+y, x−y). With standard basis e$_1$=(1,0), e$_2$=(0,1):

\begin{verbatim}
T(e₁) = (2,1) = 2e₁ + 1e₂
T(e₂) = (1,−1) = 1e₁ + (−1)e₂

Matrix [T] = | 2   1 |
             | 1  −1 |
\end{verbatim}

\subsection{Change of Basis}

If B and B' are two bases for V, and P is the change-of-basis matrix from B to B', then:

\textbf{[T]_B' = P$^{-}^1$ [T]_B P}

This is called a \textbf{similarity transformation}. Two matrices A and B are \textbf{similar} if B = P$^{-}^1$AP for some invertible P.

\begin{quote}
\textbf{ CSIR NET:} Similar matrices have the same eigenvalues, trace, determinant, rank, and characteristic polynomial!
\end{quote}

\subsection{Isomorphisms}

A linear transformation T: V  W is an \textbf{isomorphism} if it is bijective (one-to-one and onto).

\begin{tcolorbox}[colback=warningbg!30,colframe=warningborder,title=\textbf{Important -- Exam Alert},breakable,boxrule=1.5pt]
\textbf{Key Result:} Two finite-dimensional vector spaces over F are isomorphic \textbf{if and only if} they have the same dimension.
\end{tcolorbox}

So $\mathbb{R}^3\cong$ P$_2$($\mathbb{R}$) $\cong$ M$_1$ₓ$_3$($\mathbb{R}$) — they're all "the same" as vector spaces (all dimension 3).

\bigskip\noindent\rule{\textwidth}{0.4pt}\bigskip

\section{Unit 3: Dual Basis, Annihilators, and Transformations}

\begin{prereqbox}
 \textbf{Quick Refresher Before Starting:}
\textbf{Linear functional:} A linear map from V to $\mathbb{R}$ (or $\mathbb{C}$). It takes a vector, outputs a single number.
Example: f(x,y,z) = 2x + 3y - z is a linear functional on $\mathbb{R}^3$.
\textbf{Dual space V\textit{:} The set of ALL linear functionals on V. Surprisingly, dim(V}) = dim(V).
\textbf{If V has basis {e$_1$, e$_2$, e$_3$}}, then V* has a dual basis {f$_1$, f$_2$, f$_3$} where f$_i$(e$_j$) = 1 if i=j, 0 otherwise.
Think of dual basis as "coordinate extractors" — f$_1$ extracts the first coordinate.
\end{prereqbox}


\subsection{Dual Space}

\begin{tcolorbox}[colback=lightblue!30,colframe=defgreen,title=\textbf{Definition},breakable,boxrule=1.5pt]
\textbf{Definition:} The \textbf{dual space} V* of V is the vector space of all linear functionals on V.
\end{tcolorbox}

A \textbf{linear functional} is a linear transformation f: V  F (from V to the scalar field).

\subsubsection{Example}
On $\mathbb{R}^3$, the function f(x,y,z) = 2x − 3y + z is a linear functional.

\subsubsection{Dual Basis}
If B = {v$_1$,...,v$_n$} is a basis for V, the \textbf{dual basis} B* = {f$_1$,...,f$_n$} is defined by:

\begin{verbatim}
fᵢ(vⱼ) = δᵢⱼ = { 1 if i=j}
                  { 0 if i≠j}
\end{verbatim}

\begin{tcolorbox}[colback=warningbg!30,colframe=warningborder,title=\textbf{Important -- Exam Alert},breakable,boxrule=1.5pt]
\textbf{Key fact:} dim(V*) = dim(V)
\end{tcolorbox}

\subsubsection{Example}
For $\mathbb{R}^2$, with standard basis {e$_1$=(1,0), e$_2$=(0,1)}:
\begin{itemize}[leftmargin=*]
  \item f$_1$(x,y) = x (picks out the first coordinate)
  \item f$_2$(x,y) = y (picks out the second coordinate)
\end{itemize}

\subsection{Annihilators}

For a subspace W $\subseteq$ V, the \textbf{annihilator} W$^0$ is:

\textbf{W$^0$ = {f $\in$ V* : f(w) = 0 for all w $\in$ W}}

\subsubsection{ CSIR NET Key Formula:}
\textbf{dim(W) + dim(W$^0$) = dim(V)}

\subsubsection{Example}
V = $\mathbb{R}^3$, W = span{(1,0,0), (0,1,0)} (the xy-plane).
W$^0$ = {f(x,y,z) = cz : c $\in\mathbb{R}$} — all functionals that vanish on the xy-plane.
dim(W) = 2, dim(W$^0$) = 1, sum = 3 = dim(V) 

\subsection{Transpose (Dual) of a Transformation}

If T: V  W is linear, the \textbf{transpose} T\textit{: W}  V* is defined by:

\textbf{(T\textit{f)(v) = f(T(v))} for all f $\in$ W}, v $\in$ V

The matrix of T* is the transpose of the matrix of T (with respect to dual bases).

\bigskip\noindent\rule{\textwidth}{0.4pt}\bigskip

\section{Unit 4: Inner Product Spaces}

\begin{prereqbox}
 \textbf{Quick Refresher Before Starting:}
\textbf{Dot product} (from Class 12): For u = (u$_1$,u$_2$), v = (v$_1$,v$_2$):
u $\cdot$ v = u$_1$v$_1$ + u$_2$v$_2$ (multiply corresponding entries and add)
Example: (3,4) $\cdot$ (1,2) = 3$\times$1 + 4$\times$2 = 11
\textbf{Key properties of dot product:}
- u $\cdot$ v = |u||v|cos$\theta$ where $\theta$ is the angle between them
- u $\perp$ v (perpendicular) $\Longleftrightarrow$ u $\cdot$ v = 0
- |u| = $\sqrt{u $\cdot$ u}$ = $\sqrt{u$_1^2$ + u$_2^2$}$ (Pythagoras!)
\textbf{Inner product} is a generalization of dot product. It can be defined on ANY vector space, not just $\mathbb{R}^n$.
For polynomials: ⟨p,q⟩ = $\int_0^1$ p(x)q(x)dx is a valid inner product!
\end{prereqbox}


\subsection{What is an Inner Product?}

An \textbf{inner product} on a real vector space V is a function ⟨$\cdot$,$\cdot$⟩: V $\times$ V  $\mathbb{R}$ satisfying:

\begin{itemize}
  \item \textbf{Symmetry:} ⟨u,v⟩ = ⟨v,u⟩
  \item \textbf{Linearity in first argument:} ⟨au+bv, w⟩ = a⟨u,w⟩ + b⟨v,w⟩
  \item \textbf{Positive definiteness:} ⟨v,v⟩ $\geq$ 0, and ⟨v,v⟩ = 0 iff v = 0
\end{itemize}

\textbf{In plain English:} An inner product is a way to measure "angles" and "lengths" in abstract vector spaces.

\subsubsection{The Standard Inner Product on $\mathbb{R}^n$}
⟨(x$_1$,...,x$_n$), (y$_1$,...,y$_n$)⟩ = x$_1$y$_1$ + x$_2$y$_2$ + ... + x$_n$y$_n$ (the dot product you know!)

\subsubsection{Norm (Length)}
\begin{center}
\begin{tabular}{|l|l|l|}
\hline
\textbf{} & \textbf{v} & \textbf{} \\
\hline
\end{tabular}
\end{center}

\subsubsection{ CSIR NET: Cauchy-Schwarz Inequality}
\textbf{|⟨u,v⟩| $\leq$ ||u|| $\cdot$ ||v||}

Equality holds iff u and v are linearly dependent (parallel).

\subsubsection{ CSIR NET: Triangle Inequality}
\textbf{||u + v|| $\leq$ ||u|| + ||v||}

\subsubsection{Worked Example}
In $\mathbb{R}^3$ with standard inner product, let u = (1,2,3), v = (4,−1,2):

\begin{verbatim}
⟨u,v⟩ = 1(4) + 2(−1) + 3(2) = 4 − 2 + 6 = 8
||u|| = √(1+4+9) = √14
||v|| = √(16+1+4) = √21

Check C-S: |8| ≤ √14 · √21 = √294 ≈ 17.15 ✓
\end{verbatim}

\subsection{Orthogonality}

Two vectors u, v are \textbf{orthogonal} if ⟨u,v⟩ = 0.

A set {v$_1$,...,v$_k$} is \textbf{orthogonal} if ⟨v$_i$,v$_j$⟩ = 0 for all i $\neq$ j.

It is \textbf{orthonormal} if additionally ||v$_i$|| = 1 for all i.

\bigskip\noindent\rule{\textwidth}{0.4pt}\bigskip

\section{Unit 5: Orthonormal Systems and Orthogonal Transformations}

\begin{prereqbox}
 \textbf{Quick Refresher Before Starting:}
\textbf{Orthogonal} = perpendicular (inner product = 0).
\textbf{Orthonormal} = orthogonal AND each vector has length 1.
Example: {(1,0), (0,1)} is orthonormal. {(1,1), (-1,1)} is orthogonal but NOT orthonormal (length = $\sqrt{2}$).
To make it orthonormal: divide each by its length  {(1/$\sqrt{2}$, 1/$\sqrt{2}$), (-1/$\sqrt{2}$, 1/$\sqrt{2}$)}.
\textbf{Gram-Schmidt Process} (you'll learn this here): Takes ANY basis and produces an orthonormal basis.
It works by: take first vector, normalize it. For each next vector, subtract its projections onto previous orthonormal vectors, then normalize.
\end{prereqbox}


\subsection{Gram-Schmidt Process  (CSIR NET FAVORITE!)}

Given a basis {v$_1$,...,v$_n$}, produce an orthonormal basis {e$_1$,...,e$_n$}:

\begin{verbatim}
Step 1: u₁ = v₁,  e₁ = u₁/||u₁||

Step 2: u₂ = v₂ − ⟨v₂,e₁⟩e₁,  e₂ = u₂/||u₂||

Step 3: u₃ = v₃ − ⟨v₃,e₁⟩e₁ − ⟨v₃,e₂⟩e₂,  e₃ = u₃/||u₃||

General: uₖ = vₖ − Σᵢ₌₁ᵏ⁻¹ ⟨vₖ,eᵢ⟩eᵢ,  eₖ = uₖ/||uₖ||
\end{verbatim}

\subsubsection{Worked Example}
Apply Gram-Schmidt to {v$_1$=(1,1,0), v$_2$=(1,0,1), v$_3$=(0,1,1)} in $\mathbb{R}^3$:

\textbf{Step 1:}
u$_1$ = (1,1,0), ||u$_1$|| = $\sqrt{2}$, e$_1$ = (1/$\sqrt{2}$, 1/$\sqrt{2}$, 0)

\textbf{Step 2:}
⟨v$_2$,e$_1$⟩ = 1/$\sqrt{2}$ + 0 + 0 = 1/$\sqrt{2}$
u$_2$ = (1,0,1) − (1/$\sqrt{2}$)(1/$\sqrt{2}$, 1/$\sqrt{2}$, 0) = (1,0,1) − ($\frac{1}{2}$, $\frac{1}{2}$, 0) = ($\frac{1}{2}$, −$\frac{1}{2}$, 1)
\begin{center}
\begin{tabular}{|l|l|l|}
\hline
\textbf{} & \textbf{u$_2$} & \textbf{} \\
\hline
\end{tabular}
\end{center}
e$_2$ = (1/$\sqrt{6}$, −1/$\sqrt{6}$, 2/$\sqrt{6}$)

\textbf{Step 3:}
⟨v$_3$,e$_1$⟩ = 1/$\sqrt{2}$, ⟨v$_3$,e$_2$⟩ = −1/$\sqrt{6}$ + 2/$\sqrt{6}$ = 1/$\sqrt{6}$
u$_3$ = (0,1,1) − (1/$\sqrt{2}$)e$_1$ − (1/$\sqrt{6}$)e$_2$ = (−$\frac{2}{3}$, $\frac{2}{3}$, $\frac{2}{3}$)... [normalize]
e$_3$ = (−1/$\sqrt{3}$, 1/$\sqrt{3}$, 1/$\sqrt{3}$)

\subsection{Orthogonal Matrices & Transformations}

\begin{tcolorbox}[colback=lightblue!30,colframe=defgreen,title=\textbf{Definition},breakable,boxrule=1.5pt]
\textbf{Definition:} A matrix Q is \textbf{orthogonal} if Q^T Q = I (equivalently, Q^T = Q$^{-}^1$).
\end{tcolorbox}

\textbf{Properties of orthogonal matrices:}
\begin{itemize}[leftmargin=*]
  \item det(Q) = $\pm$1
  \item Columns form an orthonormal set
  \item Preserves lengths: ||Qv|| = ||v||
  \item Preserves angles: ⟨Qu, Qv⟩ = ⟨u,v⟩
\end{itemize}

\subsubsection{ CSIR NET Type Question}
\textbf{Q:} Which of the following is an orthogonal matrix?

\begin{verbatim}
A = |cos θ  −sin θ|    B = |1  1|    C = |1  0|
    |sin θ   cos θ|        |0  1|        |0  2|
\end{verbatim}

\textbf{Answer:} A only. Check: A^T A = I . B and C fail this test.

\bigskip\noindent\rule{\textwidth}{0.4pt}\bigskip

\section{Unit 6: Inner Product Space Isomorphism}

\begin{prereqbox}
 \textbf{Quick Refresher Before Starting:}
\textbf{Isomorphism} = a perfect structural match between two spaces.
If T: V  W is a bijective linear map, then V and W are "the same" structurally.
\textbf{Key theorem:} V $\cong$ W (isomorphic) $\Longleftrightarrow$ dim(V) = dim(W). So ALL n-dimensional real vector spaces are "the same" as $\mathbb{R}^n$!
This unit shows: inner product spaces of the same dimension are isomorphic in a way that preserves lengths and angles.
\end{prereqbox}


\subsection{Riesz Representation Theorem}

\begin{tcolorbox}[colback=lightblue,colframe=thmblue,title=\textbf{Theorem},breakable,boxrule=1.5pt]
\textbf{Theorem:} For every linear functional f $\in$ V*, there exists a unique vector v$_0\in$ V such that:
\end{tcolorbox}

\textbf{f(v) = ⟨v, v$_0$⟩ for all v $\in$ V}

\textbf{In plain English:} Every linear functional on an inner product space is just "taking the inner product with some fixed vector."

This gives an \textbf{isomorphism} V $\cong$ V* (the space equals its dual, via inner product).

\subsubsection{Example}
On $\mathbb{R}^3$, the functional f(x,y,z) = 2x − 3y + z is represented by v$_0$ = (2,−3,1) because:
f(v) = ⟨v, (2,−3,1)⟩ for all v.

\subsection{Orthogonal Complements}

For a subspace W $\subseteq$ V: \textbf{W$\perp$ = {v $\in$ V : ⟨v,w⟩ = 0 for all w $\in$ W}}

\subsubsection{Key Results:}
\begin{itemize}[leftmargin=*]
  \item V = W $\oplus$ W$\perp$ (direct sum decomposition)
  \item dim(W) + dim(W$\perp$) = dim(V)
  \item (W$\perp$)$\perp$ = W
\end{itemize}

\subsubsection{Projection Theorem}
Every v $\in$ V can be uniquely written as v = w + w$\perp$ where w $\in$ W, w$\perp\in$ W$\perp$.

The \textbf{orthogonal projection} onto W is: proj_W(v) = $\Sigma_i$ ⟨v,e$_i$⟩e$_i$ (where {e$_i$} is an orthonormal basis for W)

\bigskip\noindent\rule{\textwidth}{0.4pt}\bigskip

\section{Unit 7: Algebra of Hom(V,V)}

\begin{prereqbox}
 \textbf{Quick Refresher Before Starting:}
\textbf{Polynomial in a matrix:} If p(x) = x$^2$ - 3x + 2, then p(A) = A$^2$ - 3A + 2I.
Example: A = [[1,0],[0,2]]. A$^2$ = [[1,0],[0,4]]. p(A) = [[1,0],[0,4]] - [[3,0],[0,6]] + [[2,0],[0,2]] = [[0,0],[0,0]].
So A satisfies the polynomial x$^2$ - 3x + 2 = (x-1)(x-2)!
\textbf{Hom(V,V)} = set of all linear transformations from V to itself. This set is itself a vector space AND an algebra (you can multiply = compose transformations).
\textbf{Minimal polynomial:} The smallest-degree polynomial that A satisfies.
\end{prereqbox}


\subsection{The Space of All Linear Operators}

\textbf{Hom(V,V)} = {T: V  V | T is linear} is itself a vector space of dimension n$^2$.

It is also an \textbf{algebra} (you can compose/multiply transformations).

\subsection{Key Concepts}

\subsubsection{Eigenvalues and Eigenvectors  (HIGHEST PRIORITY FOR CSIR NET)}

\begin{tcolorbox}[colback=lightblue!30,colframe=defgreen,title=\textbf{Definition},breakable,boxrule=1.5pt]
\textbf{Definition:} $\lambda$ is an eigenvalue of T if T(v) = $\lambda$v for some v $\neq$ 0. Such v is an eigenvector.
\end{tcolorbox}

\textbf{Characteristic polynomial:} p($\lambda$) = det(A − $\lambda$I)

\textbf{To find eigenvalues:} Solve det(A − $\lambda$I) = 0

\subsubsection{Worked Example}
\begin{verbatim}
A = |2  1|
    |1  2|

det(A − λI) = (2−λ)² − 1 = λ² − 4λ + 3 = (λ−1)(λ−3) = 0

Eigenvalues: λ₁ = 1, λ₂ = 3

For λ₁ = 1: (A−I)v = 0 → |1  1| |x| = |0| → x+y=0 → v₁ = (1,−1)
                            |1  1| |y|   |0|

For λ₂ = 3: (A−3I)v = 0 → |−1  1| |x| = |0| → −x+y=0 → v₂ = (1,1)
                             | 1 −1| |y|   |0|
\end{verbatim}

\subsubsection{Minimal Polynomial}
The \textbf{minimal polynomial} m($\lambda$) is the monic polynomial of least degree such that m(A) = 0.

\begin{tcolorbox}[colback=warningbg!30,colframe=warningborder,title=\textbf{Important -- Exam Alert},breakable,boxrule=1.5pt]
\textbf{Key facts:}
- m($\lambda$) divides the characteristic polynomial p($\lambda$)
- They have the same roots (but possibly different multiplicities)
- A is diagonalizable iff m($\lambda$) has no repeated roots
\end{tcolorbox}

\bigskip\noindent\rule{\textwidth}{0.4pt}\bigskip

\section{Unit 8: Diagonalization and Invariant Subspaces}

\begin{prereqbox}
 \textbf{Quick Refresher Before Starting:}
\textbf{Eigenvalue and Eigenvector:} If Av = $\lambda$v (A stretches v by factor $\lambda$), then:
- $\lambda$ is an eigenvalue, v is an eigenvector
- To find eigenvalues: solve det(A - $\lambda$I) = 0
\textbf{Worked example:} A = [[2,1],[0,3]].
det(A - $\lambda$I) = det([[2-$\lambda$, 1],[0, 3-$\lambda$]]) = (2-$\lambda$)(3-$\lambda$) = 0
So $\lambda_1$ = 2, $\lambda_2$ = 3.
For $\lambda$=2: (A-2I)v = 0  [[0,1],[0,1]]v = 0  v = (1,0).
For $\lambda$=3: (A-3I)v = 0  [[-1,1],[0,0]]v = 0  v = (1,1).
\textbf{Diagonalization:} A = PDP$^{-}^1$ where D = diagonal matrix of eigenvalues, P = matrix of eigenvectors. Only works if there are enough independent eigenvectors.
\end{prereqbox}


\subsection{Diagonalization }

A matrix A is \textbf{diagonalizable} if there exists an invertible P such that P$^{-}^1$AP = D (diagonal).

\textbf{When is A diagonalizable?}
\begin{itemize}[leftmargin=*]
  \item A has n linearly independent eigenvectors
  \item Equivalently, the minimal polynomial has no repeated roots
  \item For each eigenvalue: geometric multiplicity = algebraic multiplicity
\end{itemize}

\subsubsection{Example (Diagonalizable)}
\begin{verbatim}
A = |4  1|, eigenvalues: λ=3, λ=5
    |0  5|  (note: not symmetric, but still diagonalizable since distinct eigenvalues)

P = |1  1|, P⁻¹AP = |3  0|
    |−1 1|           |0  5|
\end{verbatim}

\subsubsection{Example (NOT Diagonalizable)}
\begin{verbatim}
A = |1  1|, eigenvalue: λ=1 (repeated)
    |0  1|  Only one eigenvector: (1,0). Not diagonalizable!
\end{verbatim}

\subsection{Invariant Subspaces}

A subspace W $\subseteq$ V is \textbf{T-invariant} if T(W) $\subseteq$ W (T maps W into itself).

\begin{tcolorbox}[colback=examplebg!30,colframe=exampleborder,title=\textbf{Worked Example},breakable,boxrule=1pt]
\textbf{Examples of invariant subspaces:}
- {0} and V (trivially)
- ker(T) and Im(T)
- Eigenspaces E$\lambda$ = ker(T − $\lambda$I)
\end{tcolorbox}

\bigskip\noindent\rule{\textwidth}{0.4pt}\bigskip

\section{Unit 9: Cayley-Hamilton, Canonical Forms, Jordan Form }

\begin{prereqbox}
 \textbf{Quick Refresher Before Starting:}
\textbf{Characteristic polynomial:} p($\lambda$) = det(A - $\lambda$I). Its roots are the eigenvalues.
\textbf{Cayley-Hamilton Theorem} (you'll prove here): Every matrix satisfies its own characteristic polynomial. p(A) = 0.
Example: If char poly is $\lambda^2$ - 5$\lambda$ + 6 = 0, then A$^2$ - 5A + 6I = 0 (the zero matrix!).
\textbf{Jordan form:} When a matrix can't be diagonalized, Jordan form is the "next best thing."
A Jordan block J$_2$($\lambda$) = [[$\lambda$,1],[0,$\lambda$]] — it's almost diagonal, with 1's on the superdiagonal.
\end{prereqbox}


\subsection{Cayley-Hamilton Theorem}

\textbf{Every square matrix satisfies its own characteristic polynomial.}

If p($\lambda$) = det(A − $\lambda$I), then \textbf{p(A) = 0} (the zero matrix).

\subsubsection{Example}
\begin{verbatim}
A = |1  2|
    |3  4|

p(λ) = λ² − 5λ − 2

Check: A² − 5A − 2I = |7  10| − |5  10| − |2  0| = |0  0| ✓
                        |15 22|   |15 20|   |0  2|   |0  0|
\end{verbatim}

\subsubsection{ CSIR NET Application}
Use Cayley-Hamilton to compute A$^{-}^1$:
From A$^2$ − 5A − 2I = 0: A$^2$ = 5A + 2I  A(A − 5I) = 2I  A$^{-}^1$ = (A − 5I)/2

\subsection{Jordan Canonical Form}

When a matrix is NOT diagonalizable, we use the \textbf{Jordan form} — the "next best thing."

A \textbf{Jordan block} J$_k$($\lambda$) is a k$\times$k matrix:
\begin{verbatim}
Jₖ(λ) = |λ  1  0  ··· 0|
         |0  λ  1  ··· 0|
         |0  0  λ  ··· 0|
         |⋮        ⋱   1|
         |0  0  0  ··· λ|
\end{verbatim}

\textbf{Jordan Normal Form Theorem:} Every square matrix over $\mathbb{C}$ is similar to a block diagonal matrix of Jordan blocks.

\subsubsection{Example}
\begin{verbatim}
A = |5  4  2  1|     Jordan form:  |1  0  0  0|
    |0  1  −1 −1|                   |0  1  1  0|
    |−1 −1  3  0|                   |0  0  1  0|
    |1  1  −1  2|                   |0  0  0  5|

Here: J = J₁(1) ⊕ J₂(1) ⊕ J₁(5)
\end{verbatim}

\bigskip\noindent\rule{\textwidth}{0.4pt}\bigskip

\section{Unit 10: Forms on Vector Spaces}

\begin{prereqbox}
 \textbf{Quick Refresher Before Starting:}
\textbf{Quadratic form} = a homogeneous degree-2 polynomial in several variables.
Examples: Q(x,y) = 3x$^2$ + 2xy - y$^2$ or Q(x,y,z) = x$^2$ + 4y$^2$ + z$^2$ - 2xz
Every quadratic form can be written as Q(x) = xᵀAx where A is symmetric.
Example: Q = 3x$^2$ + 2xy - y$^2$ corresponds to A = [[3,1],[1,-1]] (off-diagonal = half the xy coefficient).
\textbf{Classification:}
- Positive definite: Q > 0 for all x $\neq$ 0 (like x$^2$ + y$^2$ — a bowl shape)
- Negative definite: Q < 0 for all x $\neq$ 0
- Indefinite: Q takes both positive and negative values (like x$^2$ - y$^2$ — a saddle)
\end{prereqbox}


\subsection{Bilinear Forms}

A \textbf{bilinear form} on V is a function B: V $\times$ V  F that is linear in each argument:
\begin{itemize}[leftmargin=*]
  \item B(au+bv, w) = aB(u,w) + bB(v,w)
  \item B(u, av+bw) = aB(u,v) + bB(u,w)
\end{itemize}

\textbf{Matrix representation:} B(x,y) = x^T A y

\subsubsection{Types:}
\begin{itemize}[leftmargin=*]
  \item \textbf{Symmetric:} B(u,v) = B(v,u)  A^T = A
  \item \textbf{Skew-symmetric:} B(u,v) = −B(v,u)  A^T = −A
  \item \textbf{Non-degenerate:} B(u,v) = 0 for all v implies u = 0
\end{itemize}

\subsection{Quadratic Forms }

A \textbf{quadratic form} Q: V  F is Q(v) = B(v,v) for some symmetric bilinear form B.

\subsubsection{Example}
Q(x,y,z) = 2x$^2$ + 3y$^2$ − z$^2$ + 4xy − 2xz

Matrix: A = | 2   2  −1|
            | 2   3   0|
            |−1   0  −1|

\subsubsection{Sylvester's Law of Inertia }
Every real quadratic form can be reduced to: Q = x$_1^2$ + ... + xₚ$^2$ − xₚ₊$_1^2$ − ... − xᵣ$^2$

The pair (p, r−p) is the \textbf{signature} — it's an invariant (doesn't change with basis).

\subsubsection{Classification:}
\begin{itemize}[leftmargin=*]
  \item \textbf{Positive definite:} All eigenvalues > 0 (signature (n,0))
  \item \textbf{Negative definite:} All eigenvalues < 0
  \item \textbf{Indefinite:} Mixed signs
  \item \textbf{Positive semi-definite:} All eigenvalues $\geq$ 0
\end{itemize}

\bigskip\noindent\rule{\textwidth}{0.4pt}\bigskip


\textbf{Q1.} Let T: $\mathbb{R}^3\mathbb{R}^3$ be T(x,y,z) = (y+z, x+z, x+y). Find eigenvalues of T.

\textbf{Solution:} Matrix A = |0 1 1|, det(A−$\lambda$I) = −$\lambda^3$+3$\lambda$+2 = −($\lambda$−2)($\lambda$+1)$^2$ = 0
                          |1 0 1|
                          |1 1 0|
Eigenvalues: $\lambda$ = 2 (once), $\lambda$ = −1 (twice)

\textbf{Q2.} If A is a 5$\times$5 matrix with characteristic polynomial ($\lambda$−2)$^3$($\lambda$+1)$^2$ and minimal polynomial ($\lambda$−2)$^2$($\lambda$+1), what are the possible Jordan forms?

\textbf{Solution:} Jordan blocks for $\lambda$=2: one J$_2$(2) + one J$_1$(2). For $\lambda$=−1: two J$_1$(−1).
Jordan form = diag(J$_2$(2), J$_1$(2), J$_1$(−1), J$_1$(−1))

\textbf{Q3.} True or False: Every real symmetric matrix is diagonalizable.
\textbf{Answer:} TRUE. (Spectral Theorem — orthogonally diagonalizable!)

\textbf{Q4.} Let V = P$_3$($\mathbb{R}$). Define T(p) = p'. What is the Jordan form of T?
\textbf{Solution:} T is nilpotent: T$^4$ = 0. Matrix is 4$\times$4 with $\lambda$=0 only. Jordan form = J₄(0).

\bigskip\noindent\rule{\textwidth}{0.4pt}\bigskip

\begin{quote}
\textbf{ Recommended Reading:} Hoffman & Kunze "Linear Algebra", Friedberg et al. "Linear Algebra"
\textbf{Next Subject:} [Mathematical Analysis ](./02_Mathematical_Analysis.md)
\end{quote}

\newpage

\chapter{Linear Algebra --- Proofs \& Problems}


\begin{quote}
\textbf{This supplements 01_Advanced_Linear_Algebra.md with complete proofs and extra examples needed for >95%}
\end{quote}

\bigskip\noindent\rule{\textwidth}{0.4pt}\bigskip


\subsection{Proof 1: Rank-Nullity Theorem ⚡⚡⚡}

\begin{tcolorbox}[colback=lightblue,colframe=thmblue,title=\textbf{Theorem},breakable,boxrule=1.5pt]
\textbf{Theorem:} If T: V  W is linear and V is finite-dimensional, then dim(V) = dim(ker T) + dim(Im T).
\end{tcolorbox}

\begin{proof}

Let dim(ker T) = k. Choose basis {v$_1$,...,v$_k$} for ker T.
Extend to basis {v$_1$,...,v$_k$, v$_k$₊$_1$,...,v$_n$} for V.
\end{proof}

\textbf{Claim:} {T(v$_k$₊$_1$),...,T(v$_n$)} is a basis for Im T.

\textbf{Spanning:} Any w $\in$ Im T has form w = T(c$_1$v$_1$+...+c$_n$v$_n$) = c$_1$T(v$_1$)+...+c$_n$T(v$_n$) = 0+...+0+c$_k$₊$_1$T(v$_k$₊$_1$)+...+c$_n$T(v$_n$). 

\textbf{Independence:} Suppose c$_k$₊$_1$T(v$_k$₊$_1$)+...+c$_n$T(v$_n$) = 0.
Then T(c$_k$₊$_1$v$_k$₊$_1$+...+c$_n$v$_n$) = 0, so c$_k$₊$_1$v$_k$₊$_1$+...+c$_n$v$_n\in$ ker T.
So c$_k$₊$_1$v$_k$₊$_1$+...+c$_n$v$_n$ = a$_1$v$_1$+...+a$_k$v$_k$ for some a$_i$.
Since {v$_1$,...,v$_n$} is a basis, all coefficients = 0. In particular c$_k$₊$_1$=...=c$_n$=0. 

Therefore dim(Im T) = n − k, giving n = k + (n−k). $\blacksquare$

\bigskip\noindent\rule{\textwidth}{0.4pt}\bigskip

\subsection{Proof 2: Cayley-Hamilton Theorem ⚡⚡⚡}

\begin{tcolorbox}[colback=lightblue,colframe=thmblue,title=\textbf{Theorem},breakable,boxrule=1.5pt]
\textbf{Theorem:} Every matrix satisfies its characteristic polynomial.
\end{tcolorbox}

\begin{proof}
\textbf{Proof (for diagonalizable case — exam-friendly version):}
Let A = PDP$^{-}^1$ where D = diag($\lambda_1$,...,$\lambda_n$).
Characteristic polynomial: p($\lambda$) = $\Pi$($\lambda_i$ − $\lambda$).
\end{proof}

p(A) = P$\cdot$p(D)$\cdot$P$^{-}^1$ = P$\cdot$diag(p($\lambda_1$),...,p($\lambda_n$))$\cdot$P$^{-}^1$

But p($\lambda_i$) = 0 for each eigenvalue $\lambda_i$ (by definition of characteristic polynomial).

So p(A) = P$\cdot$0$\cdot$P$^{-}^1$ = 0. $\blacksquare$

\begin{proof}
\textbf{Proof (general case via adjugate):}
Let B($\lambda$) = adj($\lambda$I − A) be the adjugate. Then:
($\lambda$I − A)$\cdot$B($\lambda$) = det($\lambda$I − A)$\cdot$I = p($\lambda$)$\cdot$I
\end{proof}

B($\lambda$) is a matrix polynomial: B($\lambda$) = B$_n$₋$_1\lambda^n^{-}^1$ + ... + B$_1\lambda$ + B$_0$

Expanding ($\lambda$I − A)(B$_n$₋$_1\lambda^n^{-}^1$ + ... + B$_0$) and comparing coefficients of $\lambda^n$, $\lambda^n^{-}^1$, ..., $\lambda^0$:

B$_n$₋$_1$ = I
B$_n$₋$_2$ − AB$_n$₋$_1$ = c$_n$₋$_1$I
...
−AB$_0$ = c$_0$I

Multiply these by A$^n$, A$^n^{-}^1$, ..., A$^0$ respectively and add:
0 = A$^n$ + c$_n$₋$_1$A$^n^{-}^1$ + ... + c$_0$I = p(A). $\blacksquare$

\bigskip\noindent\rule{\textwidth}{0.4pt}\bigskip

\subsection{Proof 3: Spectral Theorem (Real Symmetric Matrices) ⚡⚡⚡}

\begin{tcolorbox}[colback=lightblue,colframe=thmblue,title=\textbf{Theorem},breakable,boxrule=1.5pt]
\textbf{Theorem:} Every real symmetric matrix is orthogonally diagonalizable (A = QDQᵀ where Q is orthogonal).
\end{tcolorbox}

\begin{proof}
\textbf{Proof sketch (exam-level):}
\end{proof}

\textbf{Step 1:} Eigenvalues of symmetric matrices are real.
Let Av = $\lambda$v, v $\neq$ 0. Then v̄ᵀAv = $\lambda$||v||$^2$ and v̄ᵀAv = (Av̄)ᵀv = (Āv̄)ᵀv = $\lambda$̄||v||$^2$.
So $\lambda$ = $\lambda$̄, hence $\lambda\in\mathbb{R}$. 

\textbf{Step 2:} Eigenvectors for distinct eigenvalues are orthogonal.
Av$_1$ = $\lambda_1$v$_1$, Av$_2$ = $\lambda_2$v$_2$ with $\lambda_1\neq\lambda_2$.
$\lambda_1$⟨v$_1$,v$_2$⟩ = ⟨Av$_1$,v$_2$⟩ = ⟨v$_1$,Aᵀv$_2$⟩ = ⟨v$_1$,Av$_2$⟩ = $\lambda_2$⟨v$_1$,v$_2$⟩.
So ($\lambda_1$−$\lambda_2$)⟨v$_1$,v$_2$⟩ = 0, giving ⟨v$_1$,v$_2$⟩ = 0. 

\textbf{Step 3:} Induction on dimension to get full orthogonal eigenbasis. $\blacksquare$

\bigskip\noindent\rule{\textwidth}{0.4pt}\bigskip

\subsection{Proof 4: Cauchy-Schwarz Inequality ⚡⚡⚡}

\begin{tcolorbox}[colback=lightblue,colframe=thmblue,title=\textbf{Theorem},breakable,boxrule=1.5pt]
\textbf{Theorem:} |⟨u,v⟩| $\leq$ ||u||$\cdot$||v|| with equality iff u, v linearly dependent.
\end{tcolorbox}

\begin{proof}

If v = 0, both sides = 0. 
If v $\neq$ 0, for any scalar t $\in\mathbb{R}$:
0 $\leq$ ||u − tv||$^2$ = ⟨u−tv, u−tv⟩ = ||u||$^2$ − 2t⟨u,v⟩ + t$^2$||v||$^2$
\end{proof}

This quadratic in t is $\geq$ 0, so discriminant $\leq$ 0:
4⟨u,v⟩$^2$ − 4||u||$^2$||v||$^2\leq$ 0
⟨u,v⟩$^2\leq$ ||u||$^2$||v||$^2$

Taking square roots: |⟨u,v⟩| $\leq$ ||u||$\cdot$||v||. $\blacksquare$

\bigskip\noindent\rule{\textwidth}{0.4pt}\bigskip

\subsection{Proof 5: Sylvester's Law of Inertia ⚡⚡}

\begin{tcolorbox}[colback=lightblue,colframe=thmblue,title=\textbf{Theorem},breakable,boxrule=1.5pt]
\textbf{Theorem:} The signature (p, q) of a real quadratic form is invariant under change of basis.
\end{tcolorbox}

\begin{proof}
\textbf{Proof sketch:}
Suppose Q = x$_1^2$+...+xₚ$^2$−xₚ₊$_1^2$−...−xᵣ$^2$ in one basis and Q = y$_1^2$+...+yₛ$^2$−yₛ₊$_1^2$−...−yᵣ$^2$ in another.
\end{proof}

Assume s > p. Let V⁺ = span{y$_1$,...,yₛ} and V$^{-}$ = span{xₚ₊$_1$,...,x$_n$}.
dim(V⁺) + dim(V$^{-}$) = s + (n−p) > n, so V⁺ $\cap$ V$^{-}\neq$ {0}.

Take nonzero v $\in$ V⁺ $\cap$ V$^{-}$. Then Q(v) > 0 (from V⁺) and Q(v) $\leq$ 0 (from V$^{-}$). Contradiction! So s = p. $\blacksquare$

\bigskip\noindent\rule{\textwidth}{0.4pt}\bigskip


\subsection{Eigenvalue Problems — Variations}

\textbf{Ex 1:} A = [[3,1],[0,3]]. Char poly: (3−$\lambda$)$^2$ = 0. $\lambda$ = 3 (double). Only one eigenvector (1,0). NOT diagonalizable. Jordan form: J$_2$(3).

\textbf{Ex 2:} A = [[0,−1],[1,0]]. Char poly: $\lambda^2$+1 = 0. $\lambda$ = $\pm$i (no real eigenvalues). Over $\mathbb{R}$: not diagonalizable. Over $\mathbb{C}$: diagonalizable!

\textbf{Ex 3:} A = [[1,2,0],[0,1,0],[0,0,2]]. Char poly: (1−$\lambda$)$^2$(2−$\lambda$). $\lambda$=1 (double), $\lambda$=2. For $\lambda$=1: rank(A−I) = 1, nullity = 2. Two independent eigenvectors  diagonalizable!

\textbf{Ex 4:} A$^3$ = I (rotation by 120°). Eigenvalues must satisfy $\lambda^3$ = 1. So $\lambda$ = 1, e^(2$\pi$i/3), e^(4$\pi$i/3).

\textbf{Ex 5:} If A$^2$ = A (idempotent), eigenvalues satisfy $\lambda^2$ = $\lambda$, so $\lambda$ = 0 or 1 only.

\textbf{Ex 6 (NET 2023):} A is 4$\times$4 with eigenvalues 1,1,2,2 and rank(A−I) = rank(A−2I) = 2. Is A diagonalizable?
For $\lambda$=1: geometric mult = 4−rank(A−I) = 4−2 = 2 = algebraic mult. 
For $\lambda$=2: geometric mult = 4−rank(A−2I) = 4−2 = 2 = algebraic mult. 
YES, diagonalizable!

\subsection{Gram-Schmidt — Full Worked Examples}

\textbf{Ex 1:} Orthogonalize {(1,1,1), (1,2,3), (1,4,9)} in $\mathbb{R}^3$:

v$_1$ = (1,1,1), ||v$_1$|| = $\sqrt{3}$, e$_1$ = (1,1,1)/$\sqrt{3}$

⟨v$_2$,e$_1$⟩ = (1+2+3)/$\sqrt{3}$ = 6/$\sqrt{3}$ = 2$\sqrt{3}$
u$_2$ = (1,2,3) − 2$\sqrt{3}$$\cdot$(1,1,1)/$\sqrt{3}$ = (1,2,3) − (2,2,2) = (−1,0,1)
\begin{center}
\begin{tabular}{|l|l|l|}
\hline
\textbf{} & \textbf{u$_2$} & \textbf{} \\
\hline
\end{tabular}
\end{center}

⟨v$_3$,e$_1$⟩ = (1+4+9)/$\sqrt{3}$ = 14/$\sqrt{3}$
⟨v$_3$,e$_2$⟩ = (−1+0+9)/$\sqrt{2}$ = 8/$\sqrt{2}$
u$_3$ = (1,4,9) − (14/$\sqrt{3}$)e$_1$ − (8/$\sqrt{2}$)e$_2$
= (1,4,9) − ($\frac{14}{3}$)(1,1,1) − 4(−1,0,1)
= (1−$\frac{14}{3}$+4, 4−$\frac{14}{3}$, 9−$\frac{14}{3}$−4) = ($\frac{1}{3}$, −$\frac{2}{3}$, $\frac{1}{3}$)
e$_3$ = (1,−2,1)/$\sqrt{6}$

\textbf{Verify:} ⟨e$_1$,e$_2$⟩ = (−1+0+1)/$\sqrt{6}$ = 0 , ⟨e$_1$,e$_3$⟩ = (1−2+1)/$\sqrt{18}$ = 0 

\subsection{Jordan Form — Complete Examples}

\textbf{Ex 1:} A with char poly ($\lambda$−2)$^3$, minimal poly ($\lambda$−2)$^2$.
Possible Jordan forms: J$_2$(2) $\oplus$ J$_1$(2) (this is the ONLY option matching both).

\textbf{Ex 2:} A with char poly ($\lambda$−1)$^2$($\lambda$−3)$^2$, minimal poly ($\lambda$−1)($\lambda$−3)$^2$.
For $\lambda$=1: block sizes must divide 1  two J$_1$(1)
For $\lambda$=3: largest block = 2  one J$_2$(3)
But we need total size 4: J$_1$(1) $\oplus$ J$_1$(1) $\oplus$ J$_2$(3) 

\textbf{Ex 3:} Nilpotent 4$\times$4 with rank = 2, rank(A$^2$) = 1, A$^3$ = 0.
Jordan form: J$_3$(0) $\oplus$ J$_1$(0) (verify: rank of J$_3$(0) = 2 , rank of J$_3$(0)$^2$ = 1 )

\bigskip\noindent\rule{\textwidth}{0.4pt}\bigskip


\textbf{P1.} True/False: If AB = I then BA = I (for square matrices).
\textbf{Answer:} TRUE (for square matrices). If A is n$\times$n and AB = I, then A is invertible and B = A$^{-}^1$, so BA = I.

\textbf{P2.} If A is 3$\times$3 with eigenvalues 1, 2, 3, find det(A), tr(A), det(A$^{-}^1$).
\textbf{Answer:} det(A) = 1$\cdot$2$\cdot$3 = 6. tr(A) = 1+2+3 = 6. det(A$^{-}^1$) = $\frac{1}{6}$.

\textbf{P3.} Rank of [[1,2,3],[4,5,6],[7,8,9]]?
\textbf{Answer:} R$_2$  R$_2$−4R$_1$, R$_3$  R$_3$−7R$_1$: [[1,2,3],[0,−3,−6],[0,−6,−12]]. R$_3$  R$_3$−2R$_2$: [[1,2,3],[0,−3,−6],[0,0,0]]. Rank = \textbf{2}.

\textbf{P4.} A is skew-symmetric (Aᵀ = −A). Show eigenvalues are pure imaginary or zero.
\begin{proof}
Av = $\lambda$v. Then $\lambda$̄||v||$^2$ = v̄ᵀAv = (Av̄)ᵀv = (−Āv̄)ᵀv = −$\lambda$̄||v||$^2$. Wait: v̄ᵀAv = vᵀĀᵀv̄... Let me redo over $\mathbb{C}$. ⟨Av,v⟩ = ⟨v,Aᵀv⟩ = ⟨v,−Av⟩ = −⟨v,Av⟩ = −⟨Av,v⟩̄. Wait: for real A, ⟨Av,v⟩ = vᵀAv $\in\mathbb{C}$ generally. Let's use: if Av = $\lambda$v, then v\textit{Av = $\lambda$||v||$^2$ and (v}Av)\textit{ = v}Aᵀv = −v\textit{Av. So v}Av is pure imaginary, hence $\lambda$ is pure imaginary.
\end{proof}

\textbf{P5.} Dimension of the space of n$\times$n symmetric matrices?
\textbf{Answer:} n(n+1)/2. For n=3: 6. For n=4: 10.

\textbf{P6.} If A$^2$ − 5A + 6I = 0, what are possible eigenvalues?
\textbf{Answer:} $\lambda^2$ − 5$\lambda$ + 6 = 0  ($\lambda$−2)($\lambda$−3) = 0. Eigenvalues $\in$ {2, 3}.

\textbf{P7.} T: P$_2$  P$_2$ defined by T(p) = p' + p. Find matrix w.r.t. {1, x, x$^2$}.
\textbf{Answer:} T(1)=1, T(x)=1+x, T(x$^2$)=2x+x$^2$. Matrix: [[1,1,0],[0,1,2],[0,0,1]].

\textbf{P8.} Show that the trace is similarity-invariant.
\begin{proof}
tr(P$^{-}^1$AP) = tr(APP$^{-}^1$) = tr(A) (using tr(XY) = tr(YX)). $\blacksquare$
\end{proof}

\textbf{P9.} A is orthogonal, det(A) = −1. Show $\lambda$ = −1 is an eigenvalue.
\begin{proof}
Eigenvalues of orthogonal matrix have |$\lambda$|=1. Complex eigenvalues come in conjugate pairs (A is real). For odd n, one real eigenvalue must exist. Product of eigenvalues = det(A) = −1. If the real eigenvalue were 1, the remaining pairs multiply to −1, but conjugate pairs multiply to positive numbers. Contradiction. So −1 must be an eigenvalue. For even n, similar argument with parity.
\end{proof}

\textbf{P10.} Find the quadratic form corresponding to A = [[2,−1],[−1,3]].
\textbf{Answer:} Q(x,y) = 2x$^2$ + 3y$^2$ − 2xy. Eigenvalues: (5$\pm$$\sqrt{5}$)/2, both > 0  positive definite.

\textbf{P11.} A is nilpotent (Aᵏ=0). Show det(I+A) = 1.
\begin{proof}
Eigenvalues of A are all 0. Eigenvalues of I+A are all 1. det(I+A) = product = 1. $\blacksquare$
\end{proof}

\textbf{P12.} If A and B are similar, they have the same minimal polynomial.
\begin{proof}
B = P$^{-}^1$AP. If m(A) = 0, then m(B) = m(P$^{-}^1$AP) = P$^{-}^1$m(A)P = 0. So minimal poly of B divides that of A. By symmetry (A = PBP$^{-}^1$), minimal poly of A divides that of B. Hence equal. $\blacksquare$
\end{proof}

\textbf{P13-P15.} [Left as exercises — attempt before checking textbook]
P13: Find all 2$\times$2 matrices that commute with [[1,1],[0,1]].
P14: Prove that the set of diagonalizable matrices is dense in M_n($\mathbb{C}$).
P15: If T$^2$ = T, show V = ker(T) $\oplus$ Im(T).

\newpage

\chapter{Mathematical Analysis}


\begin{quote}
\textbf{CSIR NET Priority: ⭐⭐⭐⭐⭐ | Real Analysis is THE most tested area in CSIR NET}
\end{quote}

\bigskip\noindent\rule{\textwidth}{0.4pt}\bigskip

\subsection{️ Subject Mind Map}

\begin{verbatim}
                    MATHEMATICAL ANALYSIS
                           │
         ┌─────────┬───────┴────────┬──────────┐
         │         │                │          │
    Integration  Convergence    Multivariable  Applications
         │         │                │          │
    ┌────┴────┐  ┌─┴──┐        ┌───┴───┐   ┌──┴──┐
    Riemann   Uniform  Point-  Implicit  Extrema Jacobian
    Stieltjes Convg.   wise    Function
\end{verbatim}

\bigskip\noindent\rule{\textwidth}{0.4pt}\bigskip

\section{Unit 1: Riemann-Stieltjes Integral — Part I}

\begin{prereqbox}
 \textbf{Quick Refresher Before Starting:}
\textbf{Definite integral} (Class 12): $\int_0^1$ x$^2$ dx = [x$^3$/3]$_0^1$ = $\frac{1}{3}$. It gives the area under the curve y = x$^2$.
\textbf{Riemann sum:} Divide [a,b] into small intervals, approximate area by rectangles, sum them up.
As rectangles get thinner, the sum  the integral.
\textbf{Riemann-Stieltjes integral} (new concept): Instead of $\int$f(x)dx, we compute $\int$f(x)d$\alpha$(x) where $\alpha$ is another function.
Think of it as: $\alpha$ controls "how we weight" different parts of the interval.
When $\alpha$(x) = x, this reduces to the ordinary integral.
\end{prereqbox}



\subsection{Why Beyond Riemann?}

You know the Riemann integral $\int$f(x)dx from Class 12. The \textbf{Riemann-Stieltjes integral} generalizes this by replacing dx with d$\alpha$(x), where $\alpha$ is a monotonically increasing function.

\textbf{Notation:} $\int$ₐᵇ f d$\alpha$ or $\int$ₐᵇ f(x) d$\alpha$(x)

\textbf{Intuition:} Instead of measuring "width" of intervals by (x$_i$ − x$_i$₋$_1$), we measure by ($\alpha$(x$_i$) − $\alpha$(x$_i$₋$_1$)). This "reweights" the integral.

\textbf{Special case:} When $\alpha$(x) = x, we get the ordinary Riemann integral.

\subsection{Formal Definition}

\textbf{Partition:} P = {a = x$_0$ < x$_1$ < ... < x$_n$ = b}

\textbf{Upper and Lower Sums:}
\begin{verbatim}
U(P, f, α) = Σᵢ Mᵢ · Δαᵢ    where Mᵢ = sup f on [xᵢ₋₁, xᵢ], Δαᵢ = α(xᵢ) − α(xᵢ₋₁)
L(P, f, α) = Σᵢ mᵢ · Δαᵢ    where mᵢ = inf f on [xᵢ₋₁, xᵢ]
\end{verbatim}

f is \textbf{Riemann-Stieltjes integrable} w.r.t. $\alpha$ if:

\textbf{sup L(P, f, $\alpha$) = inf U(P, f, $\alpha$)}

\subsubsection{Example 1: Step Function Integrator}
Let $\alpha$(x) = ⌊x⌋ (floor function) on [0, 3]. Then:

$\int_0^3$ f(x) d$\alpha$ = f(1) + f(2) + f(3) — it becomes a \textbf{sum}!

This shows how R-S integral unifies sums and integrals.

\subsubsection{Example 2}
Compute $\int_0^1$ x d$\alpha$ where $\alpha$(x) = x$^2$.

\textbf{Solution:} Here d$\alpha$ = 2x dx ($\alpha$ is differentiable), so:
$\int_0^1$ x $\cdot$ 2x dx = $\int_0^1$ 2x$^2$ dx = [2x$^3$/3]$_0^1$ = $\frac{2}{3}$

\subsubsection{ Key Theorem: Integration by Parts}
\textbf{$\int$ₐᵇ f d$\alpha$ + $\int$ₐᵇ $\alpha$ df = f(b)$\alpha$(b) − f(a)$\alpha$(a)}

\bigskip\noindent\rule{\textwidth}{0.4pt}\bigskip

\section{Unit 2: Riemann-Stieltjes Integral — Part II}

\begin{prereqbox}
 \textbf{Quick Refresher Before Starting:}
\textbf{Integration by parts} (Class 12): $\int$u dv = uv - $\int$v du
Example: $\int$x$\cdot$e$^x$ dx. Let u = x, dv = e$^x$dx. Then du = dx, v = e$^x$.
Result: xe$^x$ - $\int$e$^x$dx = xe$^x$ - e$^x$ + C = e$^x$(x-1) + C.
This unit extends integration by parts to Stieltjes integrals: $\int$f d$\alpha$ + $\int\alpha$ df = f(b)$\alpha$(b) - f(a)$\alpha$(a).
\end{prereqbox}


\subsection{Existence Conditions}

\subsubsection{Theorem 1}
If f is continuous on [a,b] and $\alpha$ is monotonically increasing, then f $\in$ R($\alpha$) (f is R-S integrable w.r.t. $\alpha$).

\subsubsection{Theorem 2}
If f is monotonic and $\alpha$ is continuous and monotonically increasing on [a,b], then f $\in$ R($\alpha$).

\subsubsection{Theorem 3 (Reduction to Riemann)}
If f is continuous and $\alpha$ has a continuous derivative $\alpha$', then:

\textbf{$\int$ₐᵇ f d$\alpha$ = $\int$ₐᵇ f(x)$\alpha$'(x) dx}

\subsection{Properties of the R-S Integral}

\begin{itemize}
  \item \textbf{Linearity:} $\int$(af + bg)d$\alpha$ = a$\int$f d$\alpha$ + b$\int$g d$\alpha$
  \item \textbf{Linearity in integrator:} $\int$f d(c$\alpha$ + $\beta$) = c$\int$f d$\alpha$ + $\int$f d$\beta$
  \item \textbf{Additivity:} $\int$ₐᵇ f d$\alpha$ = $\int$ₐᶜ f d$\alpha$ + $\int$ᶜᵇ f d$\alpha$ (for a < c < b)
  \item \textbf{|$\int$f d$\alpha$| $\leq\int$|f| d$\alpha$} (if both integrals exist)
\end{itemize}

\subsubsection{ CSIR NET Example}
\textbf{Q:} Evaluate $\int_0^2$ f(x) d$\alpha$ where f(x) = x$^2$ and $\alpha$(x) = x + ⌊x⌋.

\textbf{Solution:} $\alpha$(x) = x + ⌊x⌋ has a jump of 1 at x = 1.
Split: $\int_0^2$ = $\int_0^1^{-}$ (continuous part) + jump contribution + $\int_1$⁺$^2$ (continuous part)

On [0,1): $\alpha$(x) = x, d$\alpha$ = dx  $\int_0^1$ x$^2$ dx = $\frac{1}{3}$
Jump at x=1: f(1)$\cdot$(jump of $\alpha$ at 1) = 1$\cdot$1 = 1
On [1,2]: $\alpha$(x) = x+1, d$\alpha$ = dx  $\int_1^2$ x$^2$ dx = $\frac{7}{3}$

\textbf{Total = $\frac{1}{3}$ + 1 + $\frac{7}{3}$ = $\frac{1}{3}$ + $\frac{3}{3}$ + $\frac{7}{3}$ = $\frac{11}{3}$}

\bigskip\noindent\rule{\textwidth}{0.4pt}\bigskip

\section{Unit 3: Uniform Convergence — Part I}

\begin{prereqbox}
 \textbf{Quick Refresher Before Starting:}
\textbf{Sequence:} a$_1$, a$_2$, a$_3$, ... (a list of numbers). Example: a$_n$ = 1/n gives 1, $\frac{1}{2}$, $\frac{1}{3}$, ...
\textbf{Convergence:} a$_n$  L means the terms get arbitrarily close to L. Example: 1/n  0.
\textbf{Series:} $\Sigma$a$_n$ = a$_1$ + a$_2$ + a$_3$ + ... Example: $\Sigma$($\frac{1}{2}$)$^n$ = 1 + $\frac{1}{2}$ + $\frac{1}{4}$ + ... = 2.
\textbf{Pointwise convergence of functions:} f$_n$(x)  f(x) means: for EACH fixed x, the sequence of numbers f$_n$(x) converges to f(x).
Example: f$_n$(x) = x$^n$ on [0,1]. For x < 1: x$^n$  0. For x = 1: 1$^n$  1. So the limit is: f(x) = 0 if x < 1, f(1) = 1.
Notice: each f$_n$ is continuous but the limit f is NOT! This motivates uniform convergence.
\end{prereqbox}


\subsection{Pointwise vs Uniform Convergence}

\subsubsection{Pointwise Convergence}
{f$_n$} converges \textbf{pointwise} to f on S if: for each x $\in$ S and each $\varepsilon$ > 0, $\exists$N(x,$\varepsilon$) such that |f$_n$(x) − f(x)| < $\varepsilon$ for all n $\geq$ N.

Note: N depends on BOTH x and $\varepsilon$.

\subsubsection{Uniform Convergence }
{f$_n$} converges \textbf{uniformly} to f on S if: for each $\varepsilon$ > 0, $\exists$N($\varepsilon$) such that |f$_n$(x) − f(x)| < $\varepsilon$ for ALL x $\in$ S and all n $\geq$ N.

Note: N depends ONLY on $\varepsilon$, NOT on x. The same N works everywhere!

\subsubsection{Critical Example: Why This Matters}
f$_n$(x) = x$^n$ on [0,1].

\textbf{Pointwise limit:}
\begin{itemize}[leftmargin=*]
  \item For x $\in$ [0,1): lim x$^n$ = 0
  \item For x = 1: lim 1$^n$ = 1
\end{itemize}

So f(x) = 0 for x $\in$ [0,1) and f(1) = 1. Each f$_n$ is continuous, but f is NOT!

\textbf{This convergence is NOT uniform} because continuity is not preserved.

\subsubsection{Uniform Convergence Preserves Continuity}
\begin{tcolorbox}[colback=lightblue,colframe=thmblue,title=\textbf{Theorem},breakable,boxrule=1.5pt]
\textbf{Theorem:} If f$_n$  f uniformly on S and each f$_n$ is continuous at c $\in$ S, then f is continuous at c.
\end{tcolorbox}

\subsection{Tests for Uniform Convergence}

\subsubsection{Weierstrass M-Test  (CSIR NET FAVORITE)}
If |f$_n$(x)| $\leq$ M$_n$ for all x $\in$ S and $\Sigma$ M$_n$ converges, then $\Sigma$ f$_n$ converges uniformly on S.

\subsubsection{Example}
$\Sigma$ sin(nx)/n$^2$ on all of $\mathbb{R}$. Here |sin(nx)/n$^2$| $\leq$ 1/n$^2$ = M$_n$ and $\Sigma$ 1/n$^2$ converges (p-series, p=2 > 1).
By M-test, the series converges uniformly. 

\bigskip\noindent\rule{\textwidth}{0.4pt}\bigskip

\section{Unit 4: Uniform Convergence — Part II}

\begin{prereqbox}
 \textbf{Quick Refresher Before Starting:}
\textbf{Pointwise vs Uniform convergence} (the key distinction):
\textbf{Pointwise:} For each x, f$_n$(x) eventually gets close to f(x). Different x's may converge at different rates.
\textbf{Uniform:} ALL x's converge at the SAME rate. Formally: sup|f$_n$(x) - f(x)|  0.
\textbf{Why it matters:} Uniform convergence preserves continuity, integrability, and (with extra conditions) differentiability. Pointwise does NOT.
\textbf{Visual:} Imagine f$_n$ as a rubber band. Pointwise = each point of the band settles down. Uniform = the ENTIRE band settles down simultaneously.
\end{prereqbox}


\subsection{Interchange Theorems }

\subsubsection{Theorem 1: Limit and Integral}
If f$_n$  f uniformly on [a,b] and each f$_n$ is integrable, then:

\textbf{$\int$ₐᵇ lim f$_n$ = lim $\int$ₐᵇ f$_n$} (can swap limit and integral)

\subsubsection{Theorem 2: Limit and Derivative}
If {f$_n$} converges at some point, each f'$_n$ exists, and {f'$_n$} converges uniformly, then:

\textbf{(lim f$_n$)' = lim f'$_n$} (can swap limit and derivative)

\subsubsection{ CSIR NET Application}
\textbf{Q:} Does $\Sigma_n$₌$_1$^$\infty$ x$^n$/n! converge uniformly on [−R, R] for any R > 0?

\textbf{Solution:} |x$^n$/n!| $\leq$ R$^n$/n! on [−R,R]. Since $\Sigma$ R$^n$/n! = eᴿ < $\infty$, by M-test, YES — uniform convergence on any bounded interval. The limit is e$^x$, and we can differentiate term by term.

\subsection{Abel's and Dirichlet's Tests}

\textbf{Abel's Test:} If $\Sigma$ a$_n$ converges and {b$_n$} is monotone bounded, then $\Sigma$ a$_n$b$_n$ converges.

\textbf{Dirichlet's Test:} If partial sums of $\Sigma$ a$_n$ are bounded and b$_n$  0 monotonically, then $\Sigma$ a$_n$b$_n$ converges.

\bigskip\noindent\rule{\textwidth}{0.4pt}\bigskip

\section{Unit 5: Approximation and Convergence Theorems}

\begin{prereqbox}
 \textbf{Quick Refresher Before Starting:}
\textbf{Polynomial:} p(x) = a$_n$x$^n$ + ... + a$_1$x + a$_0$. Continuous everywhere, easy to work with.
\textbf{Weierstrass Approximation Theorem} (key result here): ANY continuous function on [a,b] can be uniformly approximated by polynomials.
This is remarkable — even weird, jagged-looking continuous functions can be approximated by smooth polynomials!
Stone-Weierstrass generalizes this to other function algebras.
\end{prereqbox}


\subsection{Stone-Weierstrass Theorem }

\textbf{Weierstrass Approximation Theorem:} Every continuous function on [a,b] can be uniformly approximated by polynomials.

\textbf{In plain English:} No matter how complicated a continuous function is on a closed interval, polynomials can get arbitrarily close to it everywhere simultaneously.

\textbf{Stone's Generalization:} Extends this to subalgebras of C(X) on compact spaces.

\subsubsection{Example}
f(x) = |x| on [−1,1] is continuous. So there exist polynomials p$_n$(x) such that p$_n$  |x| uniformly on [−1,1], even though |x| is not differentiable at 0.

\subsection{Equicontinuity and Arzelà-Ascoli Theorem }

\textbf{Equicontinuous family:} {f$_n$} is equicontinuous at x$_0$ if for every $\varepsilon$ > 0, there exists $\delta$ > 0 (independent of n) such that |f$_n$(x) − f$_n$(x$_0$)| < $\varepsilon$ whenever |x − x$_0$| < $\delta$.

\textbf{Arzelà-Ascoli Theorem:} A sequence {f$_n$} in C[a,b] has a uniformly convergent subsequence if and only if {f$_n$} is uniformly bounded and equicontinuous.

\bigskip\noindent\rule{\textwidth}{0.4pt}\bigskip

\section{Unit 6: Calculus of Functions of Several Variables — Part I}

\begin{prereqbox}
 \textbf{Quick Refresher Before Starting:}
\textbf{Partial derivatives} (extending Class 12 calculus to several variables):
If f(x,y) = x$^2$y + 3y$^2$, then:
$\partial$f/$\partial$x = 2xy (differentiate treating y as a constant)
$\partial$f/$\partial$y = x$^2$ + 6y (differentiate treating x as a constant)
\textbf{Gradient:} $\nabla$f = ($\partial$f/$\partial$x, $\partial$f/$\partial$y) points in the direction of steepest increase.
Example: f = x$^2$ + y$^2$. $\nabla$f = (2x, 2y), pointing radially outward (away from minimum at origin).
\end{prereqbox}


\subsection{From Single Variable to Multivariable}

Now we extend differentiation to functions f: $\mathbb{R}^n\mathbb{R}$ᵐ.

\subsubsection{Partial Derivatives}
For f: $\mathbb{R}^n\mathbb{R}$, the \textbf{partial derivative} with respect to x$_i$:

$\partial$f/$\partial$x$_i$ = lim_{h0} [f(x$_1$,...,x$_i$+h,...,x$_n$) − f(x$_1$,...,x$_i$,...,x$_n$)] / h

\subsubsection{Example}
f(x,y) = x$^2$y + sin(xy)

$\partial$f/$\partial$x = 2xy + y$\cdot$cos(xy)
$\partial$f/$\partial$y = x$^2$ + x$\cdot$cos(xy)

\subsubsection{Total Derivative (Fréchet Derivative) }

f: $\mathbb{R}^n\mathbb{R}$ᵐ is \textbf{differentiable} at a if there exists a linear map L: $\mathbb{R}^n\mathbb{R}$ᵐ such that:

lim_{h0} ||f(a+h) − f(a) − L(h)|| / ||h|| = 0

L is called the \textbf{total derivative} Df(a). Its matrix is the \textbf{Jacobian matrix}.

\begin{tcolorbox}[colback=warningbg!30,colframe=warningborder,title=\textbf{Important -- Exam Alert},breakable,boxrule=1.5pt]
\textbf{Key:} Existence of partial derivatives does NOT guarantee differentiability! But if partial derivatives exist and are continuous, then f IS differentiable.
\end{tcolorbox}

\bigskip\noindent\rule{\textwidth}{0.4pt}\bigskip

\section{Unit 7: Calculus of Functions of Several Variables — Part II}

\begin{prereqbox}
 \textbf{Quick Refresher Before Starting:}
\textbf{Chain rule} (Class 12): d/dx[f(g(x))] = f'(g(x))$\cdot$g'(x).
Example: d/dx[sin(x$^2$)] = cos(x$^2$)$\cdot$2x.
\textbf{Multivariable chain rule} (new): If z = f(x,y) where x = x(t), y = y(t):
dz/dt = ($\partial$f/$\partial$x)(dx/dt) + ($\partial$f/$\partial$y)(dy/dt)
\textbf{Total derivative:} The matrix of all partial derivatives — generalizes f'(x) to multiple dimensions.
\end{prereqbox}


\subsection{Higher-Order Derivatives and Mixed Partials}

\subsubsection{Clairaut's Theorem (Equality of Mixed Partials) }
If f_xy and f_yx are both continuous at (a,b), then:

\textbf{$\partial^2$f/$\partial$x$\partial$y = $\partial^2$f/$\partial$y$\partial$x}

\subsubsection{Example}
f(x,y) = x$^3$y$^2$ + e^(xy)

f_x = 3x$^2$y$^2$ + ye^(xy)
f_xy = 6x$^2$y + e^(xy) + xye^(xy)

f_y = 2x$^3$y + xe^(xy)
f_yx = 6x$^2$y + e^(xy) + xye^(xy)  (Equal!)

\subsection{Chain Rule for Multivariable Functions}

If f: $\mathbb{R}^n\mathbb{R}$ᵐ and g: $\mathbb{R}$ᵐ  $\mathbb{R}$ᵖ are differentiable, then:

\textbf{D(g$\circ$f)(a) = Dg(f(a)) $\cdot$ Df(a)} (matrix multiplication of Jacobians)

\subsubsection{Example}
x = r cos $\theta$, y = r sin $\theta$ (polar coordinates)

If f(x,y) is given, then:
$\partial$f/$\partial$r = ($\partial$f/$\partial$x)cos $\theta$ + ($\partial$f/$\partial$y)sin $\theta$
$\partial$f/$\partial\theta$ = −($\partial$f/$\partial$x)r sin $\theta$ + ($\partial$f/$\partial$y)r cos $\theta$

\subsection{Taylor's Theorem in Several Variables}

f(a+h) = f(a) + Df(a)$\cdot$h + ½ h^T $\cdot$ D$^2$f(a) $\cdot$ h + ... (higher order terms)

where D$^2$f is the \textbf{Hessian matrix} of second partial derivatives.

\bigskip\noindent\rule{\textwidth}{0.4pt}\bigskip

\section{Unit 8: Implicit and Explicit Functions}

\begin{prereqbox}
 \textbf{Quick Refresher Before Starting:}
\textbf{Implicit function:} An equation F(x,y) = 0 that defines y as a function of x without explicitly solving for y.
Example: x$^2$ + y$^2$ = 1 (circle). We can't write y = single formula for the whole circle, but near any point (except ($\pm$1,0)), we can locally solve for y.
\textbf{Implicit differentiation} (Class 12): Differentiate both sides of F(x,y) = 0:
Fₓ + Fᵧ$\cdot$y' = 0, so y' = -Fₓ/Fᵧ
Example: x$^2$ + y$^2$ = 1  2x + 2y$\cdot$y' = 0  y' = -x/y.
\end{prereqbox}


\subsection{Implicit Function Theorem  (CSIR NET HIGH PRIORITY)}

\begin{tcolorbox}[colback=lightblue,colframe=thmblue,title=\textbf{Theorem},breakable,boxrule=1.5pt]
\textbf{Theorem:} Let F: $\mathbb{R}^n$⁺ᵐ  $\mathbb{R}$ᵐ be continuously differentiable. If F(a,b) = 0 and the partial derivative matrix $\partial$F/$\partial$y (m$\times$m block) is invertible at (a,b), then:
\end{tcolorbox}

Near (a,b), the equation F(x,y) = 0 implicitly defines y as a function of x: y = g(x).

Moreover: \textbf{Dg(a) = −[$\partial$F/$\partial$y]$^{-}^1\cdot$ [$\partial$F/$\partial$x]} evaluated at (a,b).

\subsubsection{Example}
F(x,y) = x$^2$ + y$^2$ − 1 = 0 (unit circle)

$\partial$F/$\partial$y = 2y. This is non-zero when y $\neq$ 0.

So near any point with y $\neq$ 0, we can solve for y = g(x).
dy/dx = −F_x/F_y = −2x/(2y) = −x/y  (matches implicit differentiation)

\subsection{Inverse Function Theorem }

If f: $\mathbb{R}^n\mathbb{R}^n$ is C$^1$ and Df(a) is invertible, then f is a local diffeomorphism near a.

Moreover: \textbf{D(f$^{-}^1$)(f(a)) = [Df(a)]$^{-}^1$}

\bigskip\noindent\rule{\textwidth}{0.4pt}\bigskip

\section{Unit 9: Extrema of Functions of Several Variables }

\begin{prereqbox}
 \textbf{Quick Refresher Before Starting:}
\textbf{Finding max/min} (Class 12 for one variable): Set f'(x) = 0, check f''(x).
\textbf{Two variables:} Set BOTH $\partial$f/$\partial$x = 0 AND $\partial$f/$\partial$y = 0 simultaneously. These give critical points.
\textbf{Second derivative test (2D):} Compute D = fₓₓfᵧᵧ - (fₓᵧ)$^2$.
- D > 0, fₓₓ > 0  local MINIMUM
- D > 0, fₓₓ < 0  local MAXIMUM
- D < 0  SADDLE POINT (neither max nor min)
- D = 0  test inconclusive
\textbf{Lagrange multipliers:} For extrema of f subject to constraint g = 0: solve $\nabla$f = $\lambda\nabla$g and g = 0.
\end{prereqbox}


\subsection{Local Extrema}

For f: $\mathbb{R}^n\mathbb{R}$, at a critical point (where $\nabla$f = 0), use the \textbf{Hessian matrix:}

\begin{verbatim}
H = |f_xx  f_xy|
    |f_yx  f_yy|
\end{verbatim}

\textbf{Second Derivative Test (2 variables):}
Let D = f_xx$\cdot$f_yy − (f_xy)$^2$ = det(H) at the critical point.

\begin{center}
\begin{tabular}{|l|l|}
\hline
\textbf{Condition} & \textbf{Conclusion} \\
\hline
D > 0, f_xx > 0 & Local minimum \\
\hline
D > 0, f_xx < 0 & Local maximum \\
\hline
D < 0 & Saddle point \\
\hline
D = 0 & Test inconclusive \\
\hline
\end{tabular}
\end{center}

\subsubsection{Worked Example}
f(x,y) = x$^3$ + y$^3$ − 3xy

$\nabla$f = (3x$^2$−3y, 3y$^2$−3x) = (0,0)  x$^2$=y, y$^2$=x  Critical points: (0,0) and (1,1)

At (0,0): H = |0  −3|, D = 0−9 = −9 < 0  \textbf{Saddle point}
              |−3  0|

At (1,1): H = |6  −3|, D = 36−9 = 27 > 0, f_xx = 6 > 0  \textbf{Local minimum}
              |−3  6|

\subsection{Lagrange Multipliers  (CSIR NET FAVORITE)}

To optimize f(x) subject to constraint g(x) = 0:

\textbf{$\nabla$f = $\lambda\nabla$g} (at the optimum)

\subsubsection{Example}
Maximize f(x,y) = xy subject to x$^2$ + y$^2$ = 1.

$\nabla$f = (y,x), $\nabla$g = (2x,2y)
y = 2$\lambda$x, x = 2$\lambda$y  y/(2x) = x/(2y)  y$^2$ = x$^2$

With constraint: 2x$^2$ = 1  x = $\pm$1/$\sqrt{2}$, y = $\pm$1/$\sqrt{2}$
Maximum value: f = $\frac{1}{2}$

\bigskip\noindent\rule{\textwidth}{0.4pt}\bigskip

\section{Unit 10: Jacobian and Its Properties }

\begin{prereqbox}
 \textbf{Quick Refresher Before Starting:}
\textbf{Determinant} (Class 12): For 2$\times$2 matrix [[a,b],[c,d]], det = ad - bc.
\textbf{Jacobian} = determinant of the matrix of partial derivatives.
For transformation (u,v)  (x,y): J = |$\partial$x/$\partial$u  $\partial$x/$\partial$v|
|$\partial$y/$\partial$u  $\partial$y/$\partial$v|
\textbf{Why it matters:} When changing variables in integrals:
$\int\int$f(x,y)dxdy = $\int\int$f(x(u,v), y(u,v))|J|dudv
\textbf{Example:} Polar coordinates x = rcos$\theta$, y = rsin$\theta$  J = r.
So $\int\int$f dxdy = $\int\int$f$\cdot$r drd$\theta$ (that's where the 'r' comes from!).
\end{prereqbox}


\subsection{Definition}

For f: $\mathbb{R}^n\mathbb{R}^n$, the \textbf{Jacobian} is:

J = det(Df) = det[$\partial$f$_i$/$\partial$x$_j$]

\subsection{Geometric Meaning}
\begin{center}
\begin{tabular}{|l|}
\hline
\textbf{J} \\
\hline
\end{tabular}
\end{center}

\subsubsection{Example: Polar Coordinates}
x = r cos $\theta$, y = r sin $\theta$

\begin{verbatim}
J = |∂x/∂r  ∂x/∂θ| = |cos θ  −r sin θ| = r cos²θ + r sin²θ = r
    |∂y/∂r  ∂y/∂θ|   |sin θ   r cos θ|
\end{verbatim}

So area element: dx dy = |J| dr d$\theta$ = r dr d$\theta$ (familiar from calculus!)

\subsubsection{Change of Variables Formula}

\textbf{$\int\int$_D f(x,y) dx dy = $\int\int$_D' f(x(u,v), y(u,v)) |J| du dv}

\subsubsection{Properties}
\begin{itemize}
  \item If f and g are differentiable: J(g$\circ$f) = J(g) $\cdot$ J(f) (chain rule)
  \item J(f$^{-}^1$) = 1/J(f) (at corresponding points)
  \item J = 0 means the mapping is singular (not locally invertible)
\end{itemize}

\bigskip\noindent\rule{\textwidth}{0.4pt}\bigskip


\textbf{Q1.} The sequence f$_n$(x) = x/(1+nx$^2$) converges uniformly on $\mathbb{R}$ to what function?

\textbf{Solution:} Pointwise limit: for each x, f$_n$(x)  0. Is it uniform?
\begin{center}
\begin{tabular}{|l|l|l|}
\hline
\textbf{f$_n$(x)} & \textbf{=} & \textbf{x} \\
\hline
\end{tabular}
\end{center}

\textbf{Q2.} Let f(x,y) = (x$^2$−y$^2$)/(x$^2$+y$^2$) for (x,y) $\neq$ (0,0). Do iterated limits exist at origin?

\textbf{Solution:} lim_{x0} lim_{y0} f = lim_{x0} 1 = 1
lim_{y0} lim_{x0} f = lim_{y0} (−1) = −1
Iterated limits differ  limit does NOT exist at origin.

\textbf{Q3.} Is f(x,y) = $\sqrt{|xy|}$ differentiable at (0,0)?

\textbf{Solution:} f_x(0,0) = 0, f_y(0,0) = 0 (check from definition). But along y=x:
\begin{center}
\begin{tabular}{|l|l|l|l|l|l|l|}
\hline
\textbf{f(h,h) − 0 − 0} & \textbf{/} & \textbf{} & \textbf{(h,h)} & \textbf{} & \textbf{=} & \textbf{h} \\
\hline
\end{tabular}
\end{center}

\bigskip\noindent\rule{\textwidth}{0.4pt}\bigskip

\begin{quote}
\textbf{ Recommended Reading:} Rudin "Principles of Mathematical Analysis", Apostol "Mathematical Analysis"
\textbf{Next Subject:} [Topology ](./03_Topology.md)
\end{quote}

\newpage

\chapter{Analysis --- Proofs \& Problems}


\begin{quote}
\textbf{Supplement to 02_Mathematical_Analysis.md}
\end{quote}

\bigskip\noindent\rule{\textwidth}{0.4pt}\bigskip


\subsection{Proof 1: Weierstrass M-Test ⚡⚡⚡}

\begin{tcolorbox}[colback=lightblue,colframe=thmblue,title=\textbf{Theorem},breakable,boxrule=1.5pt]
\textbf{Theorem:} If |f$_n$(x)| $\leq$ M$_n$ for all x $\in$ S and $\Sigma$M$_n$ converges, then $\Sigma$f$_n$ converges uniformly on S.
\end{tcolorbox}

\begin{proof}

Let s$_n$(x) = $\Sigma_k$₌$_1^n$ f$_k$(x) and S = $\Sigma$M$_n$ < $\infty$.
\end{proof}

For m > n: |s$_m$(x) − s$_n$(x)| = |$\Sigma_k$₌$_n$₊$_1$ᵐ f$_k$(x)| $\leq\Sigma_k$₌$_n$₊$_1$ᵐ |f$_k$(x)| $\leq\Sigma_k$₌$_n$₊$_1$ᵐ M$_k$

Since $\Sigma$M$_n$ converges, the tail $\Sigma_k$₌$_n$₊$_1$^$\infty$ M$_k$  0 as n  $\infty$.

Given $\varepsilon$ > 0, choose N so that $\Sigma_k$₌$_n$₊$_1$^$\infty$ M$_k$ < $\varepsilon$ for all n $\geq$ N.
Then |s$_m$(x) − s$_n$(x)| < $\varepsilon$ for ALL x $\in$ S and all m > n $\geq$ N.

By Cauchy criterion for uniform convergence, $\Sigma$f$_n$ converges uniformly. $\blacksquare$

\bigskip\noindent\rule{\textwidth}{0.4pt}\bigskip

\subsection{Proof 2: Uniform Convergence Preserves Continuity ⚡⚡⚡}

\begin{tcolorbox}[colback=lightblue,colframe=thmblue,title=\textbf{Theorem},breakable,boxrule=1.5pt]
\textbf{Theorem:} If f$_n$  f uniformly on S and each f$_n$ is continuous at c, then f is continuous at c.
\end{tcolorbox}

\begin{proof}

|f(x) − f(c)| $\leq$ |f(x) − f$_n$(x)| + |f$_n$(x) − f$_n$(c)| + |f$_n$(c) − f(c)|
\end{proof}

Given $\varepsilon$ > 0:
\begin{itemize}[leftmargin=*]
  \item By uniform convergence, $\exists$N: |f$_n$(x) − f(x)| < $\varepsilon$/3 for ALL x, for n $\geq$ N (choose such n)
  \item By continuity of f$_n$ at c, $\exists\delta$: |f$_n$(x) − f$_n$(c)| < $\varepsilon$/3 when |x−c| < $\delta$
\end{itemize}

Then |f(x) − f(c)| < $\varepsilon$/3 + $\varepsilon$/3 + $\varepsilon$/3 = $\varepsilon$ when |x−c| < $\delta$. $\blacksquare$

\bigskip\noindent\rule{\textwidth}{0.4pt}\bigskip

\subsection{Proof 3: Uniform Convergence and Integration ⚡⚡}

\begin{tcolorbox}[colback=lightblue,colframe=thmblue,title=\textbf{Theorem},breakable,boxrule=1.5pt]
\textbf{Theorem:} If f$_n$  f uniformly on [a,b] and each f$_n$ is integrable, then f is integrable and $\int$f$_n\int$f.
\end{tcolorbox}

\begin{proof}

|$\int$ₐᵇ f$_n$ − $\int$ₐᵇ f| = |$\int$ₐᵇ (f$_n$ − f)| $\leq\int$ₐᵇ |f$_n$ − f| $\leq$ (b−a) $\cdot$ sup|f$_n$ − f|  0. $\blacksquare$
\end{proof}

\bigskip\noindent\rule{\textwidth}{0.4pt}\bigskip

\subsection{Proof 4: Implicit Function Theorem (2D case) ⚡⚡}

\begin{tcolorbox}[colback=lightblue,colframe=thmblue,title=\textbf{Theorem},breakable,boxrule=1.5pt]
\textbf{Theorem:} If F(a,b) = 0, F is C$^1$, and Fᵧ(a,b) $\neq$ 0, then near (a,b), y = g(x) with g'(a) = −Fₓ(a,b)/Fᵧ(a,b).
\end{tcolorbox}

\begin{proof}
\textbf{Proof idea:} Apply Banach fixed point theorem to the map T(y) = y − F(x,y)/Fᵧ(a,b) for fixed x near a. Since Fᵧ(a,b) $\neq$ 0, T is a contraction near (a,b). The derivative formula follows by implicit differentiation of F(x,g(x)) = 0. $\blacksquare$
\end{proof}

\bigskip\noindent\rule{\textwidth}{0.4pt}\bigskip

\subsection{Proof 5: Second Derivative Test (2D) ⚡⚡⚡}

\begin{tcolorbox}[colback=lightblue,colframe=thmblue,title=\textbf{Theorem},breakable,boxrule=1.5pt]
\textbf{Theorem:} At critical point (a,b) with D = fₓₓfᵧᵧ − fₓᵧ$^2$ > 0, fₓₓ > 0 $\Longrightarrow$ local min.
\end{tcolorbox}

\begin{proof}
By Taylor: f(a+h,b+k) − f(a,b) = ½(fₓₓh$^2$ + 2fₓᵧhk + fᵧᵧk$^2$) + higher order
= ½fₓₓ[(h + fₓᵧk/fₓₓ)$^2$ + (D/fₓₓ$^2$)k$^2$]
\end{proof}

When D > 0 and fₓₓ > 0, the bracketed expression is sum of squares  positive. $\blacksquare$

\bigskip\noindent\rule{\textwidth}{0.4pt}\bigskip


\subsection{Uniform Convergence — 6 Examples}

\textbf{Ex 1:} f$_n$(x) = x/n on [0,1]. sup|f$_n$| = 1/n  0. UNIFORM. 
\textbf{Ex 2:} f$_n$(x) = x$^n$ on [0,1]. sup|f$_n$ − f| = sup on [0,1) of x$^n$ = 1 (at x1$^{-}$). NOT uniform.
\textbf{Ex 3:} f$_n$(x) = x$^n$ on [0, $\frac{1}{2}$]. sup = ($\frac{1}{2}$)$^n$  0. UNIFORM on [0, $\frac{1}{2}$]. 
\textbf{Ex 4:} f$_n$(x) = sin(x/n) on $\mathbb{R}$. |sin(x/n)| $\leq$ |x/n|  need bound independent of x. On $\mathbb{R}$: NOT uniform. On [−M,M]: uniform.
\textbf{Ex 5:} $\Sigma$ x$^n$/n! on [−R,R]. By M-test with M$_n$ = R$^n$/n!. UNIFORM on bounded intervals.
\textbf{Ex 6:} f$_n$(x) = nx$\cdot$e^(−nx$^2$). Max at x = 1/$\sqrt{2n}$, max value = $\sqrt{n/(2e}$)  $\infty$. NOT uniform on [0,1].

\subsection{Lagrange Multipliers — 5 Examples}

\textbf{Ex 1:} Max/min f = x+y on x$^2$+y$^2$ = 1.
$\nabla$f = (1,1) = $\lambda$(2x,2y). So x=y=1/(2$\lambda$). Constraint: 2/(4$\lambda^2$)=1, $\lambda$=$\pm$1/$\sqrt{2}$.
Max = $\sqrt{2}$ at (1/$\sqrt{2}$, 1/$\sqrt{2}$). Min = −$\sqrt{2}$ at (−1/$\sqrt{2}$, −1/$\sqrt{2}$).

\textbf{Ex 2:} Min f = x$^2$+y$^2$+z$^2$ subject to x+y+z = 3.
$\nabla$f = $\lambda\nabla$g: (2x,2y,2z) = $\lambda$(1,1,1). So x=y=z. From constraint: x=y=z=1. Min = 3.

\textbf{Ex 3:} Max f = xyz subject to x+y+z = 12, x,y,z > 0.
By AM-GM or Lagrange: x=y=z=4. Max = 64.

\textbf{Ex 4:} Closest point on plane 2x+y−z = 5 to origin.
Min x$^2$+y$^2$+z$^2$ subject to 2x+y−z=5. Get (2$\lambda$,$\lambda$,−$\lambda$) with 4$\lambda$+$\lambda$+$\lambda$=5, $\lambda$=$\frac{5}{6}$.
Point: ($\frac{5}{3}$, $\frac{5}{6}$, −$\frac{5}{6}$). Distance = 5/$\sqrt{6}$.

\textbf{Ex 5 (CSIR NET style):} Max of x$^2$y$^3$ on x$^2$+y$^2$ = 1, x,y $\geq$ 0.
Use substitution x = cos$\theta$, y = sin$\theta$. Maximize cos$^2\theta\cdot$sin$^3\theta$. Take derivative, set = 0.
Or Lagrange: 2xy$^3$ = 2$\lambda$x and 3x$^2$y$^2$ = 2$\lambda$y. From first: $\lambda$ = y$^3$ (if x$\neq$0). From second: $\lambda$ = 3x$^2$y/2.
So y$^3$ = 3x$^2$y/2, y$^2$ = 3x$^2$/2, with x$^2$+y$^2$ = 1: x$^2$+3x$^2$/2 = 1, x$^2$ = $\frac{2}{5}$, y$^2$ = $\frac{3}{5}$.
Max = ($\frac{2}{5}$)($\frac{3}{5}$)^($\frac{3}{2}$) = ($\frac{2}{5}$)$\cdot$3$\sqrt{3}$/(5$\sqrt{5}$) = 6$\sqrt{3}$/(25$\sqrt{5}$).

\bigskip\noindent\rule{\textwidth}{0.4pt}\bigskip


\textbf{P1.} Show f$_n$(x) = x$^2$/(x$^2$+n)  0 pointwise on $\mathbb{R}$. Is convergence uniform?
\textbf{Solution:} |f$_n$(x)| = x$^2$/(x$^2$+n) < x$^2$/n. But sup over $\mathbb{R}$ = lim_{x$\infty$} x$^2$/(x$^2$+n) = 1 for all n. NOT uniform.

\textbf{P2.} Compute $\int_0^1$ x$^3$ d(x$^2$).
\textbf{Solution:} $\alpha$(x) = x$^2$, $\alpha$'(x) = 2x. $\int_0^1$ x$^3\cdot$2x dx = 2$\int_0^1$ x$^4$ dx = $\frac{2}{5}$.

\textbf{P3.} f(x,y) = x$^2$y/(x$^4$+y$^2$). Does lim_{(x,y)(0,0)} f exist?
\textbf{Solution:} Along y = x$^2$: f = x$^4$/(2x$^4$) = $\frac{1}{2}$. Along y = 0: f = 0. Different paths  limit does NOT exist.

\textbf{P4.} Is f(x,y) = (x$^2$+y$^2$)sin(1/(x$^2$+y$^2$)) differentiable at (0,0)?
\textbf{Solution:} fₓ(0,0) = 0 (check). |f(h,k)−0−0|/$\sqrt{h$^2$+k$^2$}$ = $\sqrt{h$^2$+k$^2$}$|sin(1/(h$^2$+k$^2$))| $\leq$ $\sqrt{h$^2$+k$^2$}$  0. YES, differentiable.

\textbf{P5.} Jacobian of spherical coordinates (r,$\theta$,$\varphi$)  (x,y,z).
\textbf{Answer:} J = r$^2$sin$\varphi$. (Standard result — memorize this!)

\textbf{P6.} Find extrema of f(x,y) = x$^2$+y$^2$−2x−4y+5.
\textbf{Solution:} fₓ = 2x−2 = 0, fᵧ = 2y−4 = 0. Critical point (1,2). D = fₓₓfᵧᵧ−fₓᵧ$^2$ = 4−0 = 4 > 0, fₓₓ = 2 > 0. Local minimum = f(1,2) = 1+4−2−8+5 = 0.

\textbf{P7.} Prove: if $\Sigma$f$_n$ converges uniformly and each f$_n$ is continuous, then $\Sigma$f$_n$ is continuous.
\textbf{Hint:} Apply "uniform limit preserves continuity" to partial sums.

\textbf{P8.} Stone-Weierstrass: Can every continuous function on [0,1] be uniformly approximated by polynomials with integer coefficients?
\textbf{Answer:} NO. The set of such polynomials is countable at each point, but C[0,1] is uncountable.

\newpage

\chapter{Topology-1}


\begin{quote}
\textbf{CSIR NET Priority: ⭐⭐⭐⭐ | Topology questions appear regularly in Part B and C}
\end{quote}

\bigskip\noindent\rule{\textwidth}{0.4pt}\bigskip

\subsection{️ Subject Mind Map}

\begin{verbatim}
                         TOPOLOGY
                           │
         ┌─────────┬───────┴────────┬──────────┐
         │         │                │          │
    Foundations  Structure       Properties   Axioms
         │         │                │          │
    ┌────┴────┐  ┌─┴──┐        ┌───┴───┐   ┌──┴──┐
    Open/Closed  Basis         Connected  T0,T1,T2
    Boundary   Subbasis       Compact    Regular
    Interior   Nbhd bases     Path-conn  Normal
\end{verbatim}

\bigskip\noindent\rule{\textwidth}{0.4pt}\bigskip

\section{Unit 1: Basic Concepts}

\begin{prereqbox}
 \textbf{Quick Refresher Before Starting:}
\textbf{Open interval} (a,b) = {x $\in\mathbb{R}$ : a < x < b}. Does NOT include endpoints.
\textbf{Closed interval} [a,b] = {x $\in\mathbb{R}$ : a $\leq$ x $\leq$ b}. INCLUDES endpoints.
\textbf{Why (0,1) is "open":} Every point has wiggle room. Pick any x $\in$ (0,1), say x = 0.3. Then (0.2, 0.4) $\subset$ (0,1) is a small interval around x entirely inside (0,1). 
\textbf{Why [0,1] is NOT open:} The point 0 has no wiggle room to the left — any interval around 0 escapes [0,1].
\textbf{Topology} abstracts this: we declare which sets are "open" (have wiggle room) via axioms, WITHOUT needing distances or coordinates.
\end{prereqbox}



\subsection{What is Topology?}

\textbf{In plain English:} Topology studies properties of spaces that don't change when you stretch, bend, or deform them (but NOT tear or glue). It's "rubber sheet geometry."

A coffee mug and a donut are "the same" topologically — each has exactly one hole!

\subsection{Definition of a Topological Space}

A \textbf{topology} $\tau$ on a set X is a collection of subsets of X satisfying:

\begin{itemize}
  \item \textbf{$\emptyset\in\tau$ and X $\in\tau$} (empty set and whole space are open)
  \item \textbf{Finite intersections:} If U$_1$,...,U$_n\in\tau$, then U$_1\cap$ ... $\cap$ U$_n\in\tau$
  \item \textbf{Arbitrary unions:} If {U$\alpha$} $\subseteq\tau$, then $\cup\alpha$ U$\alpha\in\tau$
\end{itemize}

The pair (X, $\tau$) is called a \textbf{topological space}. Members of $\tau$ are called \textbf{open sets}.

\subsubsection{ Memory Aid: "FInite Intersections, Arbitrary Unions"  Think "FIAU"}

\subsubsection{Example 1: The Real Line $\mathbb{R}$}
Standard topology: $\tau$ = all sets that are unions of open intervals (a,b).
This is what you're used to from calculus!

\subsubsection{Example 2: Discrete Topology}
$\tau$ = P(X) (power set — ALL subsets are open). This is the "finest" topology.

\subsubsection{Example 3: Indiscrete (Trivial) Topology}
$\tau$ = {$\emptyset$, X} (only the empty set and X are open). This is the "coarsest" topology.

\subsubsection{Example 4: Cofinite Topology}
On X, $\tau$ = {$\emptyset$} $\cup$ {U $\subseteq$ X : X \ U is finite}. Open sets are those whose complement is finite.

\textbf{Verify:} $\emptyset\in\tau$ , X $\in\tau$ (X \ X = $\emptyset$ is finite) 
Finite intersection of cofinite sets is cofinite  (complement of intersection is union of finite sets = finite)
Arbitrary union of cofinite sets is cofinite  (complement of union is intersection, subset of a finite set) 

\bigskip\noindent\rule{\textwidth}{0.4pt}\bigskip

\section{Unit 2: Open and Closed Sets}

\begin{prereqbox}
 \textbf{Quick Refresher Before Starting:}
\textbf{Set operations you need:}
- A $\cup$ B (union) = elements in A OR B
- A $\cap$ B (intersection) = elements in A AND B
- Aᶜ (complement) = everything NOT in A
- \textbf{De Morgan:} (A $\cup$ B)ᶜ = Aᶜ $\cap$ Bᶜ and (A $\cap$ B)ᶜ = Aᶜ $\cup$ Bᶜ
\textbf{Key rule in topology:} A set is CLOSED if its complement is open.
So [0,1] is closed in $\mathbb{R}$ because $\mathbb{R}$\[0,1] = (-$\infty$,0) $\cup$ (1,$\infty$) is open (union of open intervals).
\end{prereqbox}


\subsection{Closed Sets}
A set F $\subseteq$ X is \textbf{closed} if X \ F is open (i.e., X \ F $\in\tau$).

\begin{quote}
\textbf{⚠️ COMMON MISCONCEPTION:} "Open" and "closed" are NOT opposites! A set can be:
- Both open and closed (called \textbf{clopen}): $\emptyset$ and X always are
- Neither open nor closed: [0,1) in $\mathbb{R}$ with standard topology
- Open but not closed: (0,1)
- Closed but not open: [0,1]
\end{quote}

\subsection{Interior, Closure, Limit Points}

\subsubsection{Interior}
\textbf{Int(A) = A°} = largest open set contained in A = $\cup${U $\in\tau$ : U $\subseteq$ A}

\subsubsection{Closure}
\textbf{Cl(A) = Ā} = smallest closed set containing A = $\cap${F : F is closed, A $\subseteq$ F}

\subsubsection{Limit Point (Accumulation Point)}
x is a \textbf{limit point} of A if every open set containing x intersects A in a point other than x itself.

\textbf{A' = set of all limit points of A} (derived set)

\subsubsection{Key Formula: \textbf{Ā = A $\cup$ A'}}

\subsubsection{Example}
In $\mathbb{R}$ (standard topology), A = (0,1) $\cup$ {2}:
\begin{itemize}[leftmargin=*]
  \item Int(A) = (0,1) — the isolated point 2 has no open set around it inside A
  \item A' = [0,1] — limit points include 0 and 1 but NOT 2 (it's isolated)
  \item Ā = [0,1] $\cup$ {2}
\end{itemize}

\subsection{Dense Sets}
A is \textbf{dense} in X if Ā = X. Example: $\mathbb{Q}$ is dense in $\mathbb{R}$.

\subsubsection{ CSIR NET Frequent Question}
\textbf{Q:} In the cofinite topology on $\mathbb{R}$, what is the closure of any infinite set A?
\textbf{A:} Ā = $\mathbb{R}$. Because any closed set containing an infinite set must be all of $\mathbb{R}$ (its complement would need to be finite, but $\mathbb{R}$ \ A could be anything).

\bigskip\noindent\rule{\textwidth}{0.4pt}\bigskip

\section{Unit 3: Boundary of a Set}

\begin{prereqbox}
 \textbf{Quick Refresher Before Starting:}
\textbf{Boundary intuition:} A boundary point is "on the edge" — every neighborhood contains points both INSIDE and OUTSIDE the set.
Example: For (0,1) in $\mathbb{R}$, the boundary is {0, 1}. Around 0, any interval (-$\varepsilon$, $\varepsilon$) contains points both in (0,1) (like $\varepsilon$/2) and outside (like -$\varepsilon$/2).
Example: For $\mathbb{Q}$ (rationals) in $\mathbb{R}$, EVERY real number is a boundary point! Because every interval contains both rationals and irrationals.
\end{prereqbox}


\subsection{Definition}
\textbf{Bd(A) = $\partial$A = Ā $\cap$ (X\A)̄} = Ā \ Int(A)

\textbf{In English:} The boundary consists of points that are "on the edge" — every neighborhood intersects both A and its complement.

\subsubsection{Example}
In $\mathbb{R}$: A = (0,1]
\begin{itemize}[leftmargin=*]
  \item Ā = [0,1], Int(A) = (0,1)
  \item $\partial$A = [0,1] \ (0,1) = {0, 1}
\end{itemize}

\subsubsection{Key Relationships}
\begin{verbatim}
X = Int(A) ∪ Bd(A) ∪ Int(X\A)   (disjoint union!)

Equivalently:
- x ∈ Int(A) ⟺ some neighborhood of x lies entirely in A
- x ∈ Int(X\A) ⟺ some neighborhood of x lies entirely outside A  
- x ∈ Bd(A) ⟺ every neighborhood of x meets both A and X\A
\end{verbatim}

\subsubsection{Properties}
\begin{itemize}
  \item $\partial$A is always closed
  \item A is open $\Longleftrightarrow$ A $\cap\partial$A = $\emptyset$ (A doesn't contain any of its boundary)
  \item A is closed $\Longleftrightarrow\partial$A $\subseteq$ A (A contains all its boundary)
  \item $\partial$A = $\partial$(X\A)
\end{itemize}

\bigskip\noindent\rule{\textwidth}{0.4pt}\bigskip

\section{Unit 4: Constructing Topologies}

\begin{prereqbox}
 \textbf{Quick Refresher Before Starting:}
\textbf{Basis for a topology} works like a basis for a vector space — a smaller collection that generates all open sets.
Example: Open intervals (a,b) form a basis for the usual topology on $\mathbb{R}$. Every open set is a union of open intervals.
\textbf{Product topology on X $\times$ Y:} Open sets are unions of "boxes" U $\times$ V where U is open in X and V is open in Y.
Think of it as: the plane $\mathbb{R}^2$ gets its topology from open rectangles (a,b) $\times$ (c,d).
\end{prereqbox}


\subsection{Basis for a Topology }

A \textbf{basis} B for a topology on X is a collection of subsets such that:
\begin{itemize}
  \item For each x $\in$ X, $\exists$B $\in$ B with x $\in$ B
  \item If x $\in$ B$_1\cap$ B$_2$ (with B$_1$, B$_2\in$ B), $\exists$B$_3\in$ B with x $\in$ B$_3\subseteq$ B$_1\cap$ B$_2$
\end{itemize}

The topology \textbf{generated} by B: U is open iff for every x $\in$ U, $\exists$B $\in$ B with x $\in$ B $\subseteq$ U.

\subsubsection{Examples of Bases}
\begin{center}
\begin{tabular}{|l|l|}
\hline
\textbf{Space} & \textbf{Basis} \\
\hline
$\mathbb{R}$ (standard) & {(a,b) : a < b} — open intervals \\
\hline
$\mathbb{R}^2$ (standard) & {open disks} or {open rectangles} \\
\hline
$\mathbb{R}$ (lower limit) & {[a,b) : a < b} — half-open intervals \\
\hline
\end{tabular}
\end{center}

\subsection{Subbasis}
A \textbf{subbasis} S is any collection of subsets whose union is X. The topology generated by S has as basis all finite intersections of members of S.

\subsection{Product Topology }
On X $\times$ Y, the \textbf{product topology} has basis: {U $\times$ V : U open in X, V open in Y}

\subsubsection{Example}
$\mathbb{R}^2$ = $\mathbb{R}\times\mathbb{R}$ with product of standard topologies gives the standard topology on $\mathbb{R}^2$.
Basis: open rectangles (a,b) $\times$ (c,d).

\subsection{Subspace Topology}
If (X,$\tau$) is a topological space and Y $\subseteq$ X, the \textbf{subspace topology} on Y is:
$\tau$_Y = {U $\cap$ Y : U $\in\tau$}

\subsubsection{Example}
[0,1] with subspace topology from $\mathbb{R}$. The set [0, $\frac{1}{2}$) is open in [0,1] because [0,$\frac{1}{2}$) = (-1, $\frac{1}{2}$) $\cap$ [0,1].

\bigskip\noindent\rule{\textwidth}{0.4pt}\bigskip

\section{Unit 5: Topological Maps, Neighbourhood Bases, Connectedness}

\begin{prereqbox}
 \textbf{Quick Refresher Before Starting:}
\textbf{Continuous function} (Class 12): No jumps, no breaks. f is continuous at a if lim(xa) f(x) = f(a).
\textbf{Topological definition:} f: X  Y is continuous if the preimage of every open set is open.
f$^{-}^1$(V) = {x $\in$ X : f(x) $\in$ V}. This must be open in X whenever V is open in Y.
\textbf{Homeomorphism:} A continuous bijection whose inverse is also continuous. If X $\cong$ Y (homeomorphic), they are "topologically the same."
Example: (0,1) $\cong\mathbb{R}$ via f(x) = tan($\pi$(x - $\frac{1}{2}$)). They look different but are topologically identical!
\end{prereqbox}


\subsection{Continuous Functions (Topological Definition) }

f: X  Y is \textbf{continuous} iff for every open set V in Y, f$^{-}^1$(V) is open in X.

\textbf{This generalizes the $\varepsilon$-$\delta$ definition from calculus!}

\subsection{Homeomorphism}
f: X  Y is a \textbf{homeomorphism} if f is bijective, continuous, and f$^{-}^1$ is continuous.

Homeomorphic spaces are "topologically identical." We write X $\cong$ Y.

\subsubsection{Example}
f: (0,1)  $\mathbb{R}$ given by f(x) = tan($\pi$(x − $\frac{1}{2}$)) is a homeomorphism.
So (0,1) $\cong\mathbb{R}$ — the open unit interval and the entire real line are topologically the same!

\subsection{Neighbourhood Bases}
A \textbf{neighbourhood base} at x is a collection Bₓ of neighborhoods of x such that every neighborhood of x contains some member of Bₓ.

\subsubsection{First Countable Space}
X is \textbf{first countable} if every point has a countable neighbourhood base.

All metric spaces are first countable (use balls B(x, 1/n)).

\bigskip\noindent\rule{\textwidth}{0.4pt}\bigskip

\section{Unit 6: Connected Spaces }

\begin{prereqbox}
 \textbf{Quick Refresher Before Starting:}
\textbf{Connected} = "one piece." A space X is connected if you CANNOT write X = A $\cup$ B where A,B are both open, both nonempty, and disjoint.
\textbf{Example:} $\mathbb{R}$ is connected. (0,1) $\cup$ (2,3) is NOT connected (two separate pieces).
\textbf{Example:} $\mathbb{Q}$ is NOT connected! For any irrational r, $\mathbb{Q}$ = (-$\infty$,r)$\cap\mathbb{Q}\cup$ (r,$\infty$)$\cap\mathbb{Q}$ splits it.
\textbf{Path connected:} Any two points can be joined by a continuous curve. Path connected $\Longrightarrow$ connected (but not always vice versa).
\end{prereqbox}


\subsection{Definition of Connectedness}

X is \textbf{connected} if it cannot be written as X = U $\cup$ V where U, V are non-empty, disjoint, open sets.

Equivalently: X is connected iff the only clopen subsets are $\emptyset$ and X.

\subsubsection{ Memory Aid: A connected space is "one piece" — you can't split it into two open chunks.}

\subsubsection{Examples}
\begin{itemize}[leftmargin=*]
  \item $\mathbb{R}$ is connected 
  \item $\mathbb{Q}$ is NOT connected (e.g., split at $\sqrt{2}$: U = $\mathbb{Q}\cap$ (−$\infty$,$\sqrt{2}$), V = $\mathbb{Q}\cap$ ($\sqrt{2}$,$\infty$))
  \item [0,1] is connected 
  \item {0} $\cup$ {1} is NOT connected (discrete)
\end{itemize}

\subsection{Key Theorems }

\subsubsection{Intermediate Value Theorem (Topological Version)}
If f: X  Y is continuous and X is connected, then f(X) is connected.

\textbf{Corollary (IVT from calculus):} If f: [a,b]  $\mathbb{R}$ is continuous with f(a) < 0 < f(b), then $\exists$c with f(c) = 0.

\subsubsection{Connected Subsets of $\mathbb{R}$}
A subset of $\mathbb{R}$ is connected iff it is an interval (possibly infinite, open, closed, or half-open).

\subsubsection{Properties}
\begin{itemize}
  \item Continuous image of connected space is connected
  \item If A is connected and A $\subseteq$ B $\subseteq$ Ā, then B is connected
  \item Union of connected sets with common point is connected
  \item Product of connected spaces is connected
\end{itemize}

\subsection{Path Connectedness}
X is \textbf{path connected} if for any x,y $\in$ X, there exists a continuous f: [0,1]  X with f(0) = x, f(1) = y.

\textbf{Path connected $\Longrightarrow$ Connected} (but NOT the converse!)

\subsubsection{ CSIR NET: Topologist's Sine Curve}
The set A = {(x, sin(1/x)) : x > 0} $\cup$ {(0,y) : −1 $\leq$ y $\leq$ 1} is connected but NOT path connected.

\bigskip\noindent\rule{\textwidth}{0.4pt}\bigskip

\section{Unit 7: Compactness }

\begin{prereqbox}
 \textbf{Quick Refresher Before Starting:}
\textbf{Open cover:} A collection of open sets whose union contains X. Example: {(-n,n) : n = 1,2,...} covers $\mathbb{R}$.
\textbf{Finite subcover:} Choosing finitely many sets from the cover that still cover X.
\textbf{Compact:} Every open cover has a finite subcover. Intuitively: compact sets behave like finite sets.
\textbf{In $\mathbb{R}^n$ (Heine-Borel):} Compact $\Longleftrightarrow$ closed AND bounded.
Example: [0,1] is compact  (closed + bounded). (0,1) is NOT compact  (not closed). $\mathbb{R}$ is NOT compact  (not bounded).
\end{prereqbox}


\subsection{Open Cover Definition}

An \textbf{open cover} of X is a collection {U$\alpha$} of open sets with X = $\cup\alpha$ U$\alpha$.

X is \textbf{compact} if every open cover has a \textbf{finite subcover}.

\subsubsection{ Memory Aid: "You can always find a finite team to do the whole job."}

\subsubsection{Example: [0,1] is compact}
Take any open cover. By the Heine-Borel theorem (proven using completeness), we can extract a finite subcover.

\subsubsection{Example: (0,1) is NOT compact}
Cover: {(1/n, 1) : n $\geq$ 1}. This covers (0,1) but no finite subcollection does (always misses points near 0).

\subsection{Heine-Borel Theorem  (CRITICAL for CSIR NET)}

A subset of $\mathbb{R}^n$ is compact \textbf{if and only if} it is \textbf{closed and bounded}.

\subsubsection{Quick Classification}
\begin{center}
\begin{tabular}{|l|l|l|}
\hline
\textbf{Set} & \textbf{Compact?} & \textbf{Why?} \\
\hline
[0,1] & Yes & Closed + bounded \\
\hline
(0,1) & No & Not closed \\
\hline
[0,$\infty$) & No & Not bounded \\
\hline
{1/n : n $\in\mathbb{N}$} & No & Not closed (0 is a limit point not in the set) \\
\hline
{0} $\cup$ {1/n : n $\in\mathbb{N}$} & Yes & Closed + bounded \\
\hline
S$^n$ (n-sphere) & Yes & Closed + bounded in $\mathbb{R}^n$⁺$^1$ \\
\hline
\end{tabular}
\end{center}

\bigskip\noindent\rule{\textwidth}{0.4pt}\bigskip

\section{Unit 8: Compact Spaces — Properties }

\begin{prereqbox}
 \textbf{Quick Refresher Before Starting:}
\textbf{Why compactness matters} — it guarantees things that can fail in non-compact spaces:
- Continuous functions on compact sets are bounded and attain their max/min (Extreme Value Theorem)
- Continuous functions on compact sets are uniformly continuous
- In Hausdorff spaces: compact sets are closed
- A continuous bijection from compact to Hausdorff is automatically a homeomorphism!
\end{prereqbox}


\subsection{Key Theorems}

\begin{itemize}
  \item \textbf{Closed subset of compact space is compact}
  \item \textbf{Compact subset of Hausdorff space is closed}
  \item \textbf{Continuous image of compact space is compact}
  \item \textbf{Continuous bijection from compact to Hausdorff is homeomorphism}
  \item \textbf{Extreme Value Theorem:} Continuous real-valued function on compact space attains max and min
\end{itemize}

\subsubsection{ CSIR NET Application}
\textbf{Q:} f: [0,1]  $\mathbb{R}$ is continuous. Must f be bounded?
\textbf{A:} Yes! [0,1] is compact, so f([0,1]) is compact in $\mathbb{R}$, hence closed and bounded.

\subsection{Tychonoff's Theorem}
Product of compact spaces is compact (in the product topology). This is one of the most powerful theorems in topology.

\subsection{Sequential Compactness}
In metric spaces: X is compact $\Longleftrightarrow$ every sequence has a convergent subsequence.

\subsection{Bolzano-Weierstrass Property}
Every infinite subset of a compact space has a limit point.

\bigskip\noindent\rule{\textwidth}{0.4pt}\bigskip

\section{Unit 9: Separation Axioms}

\begin{prereqbox}
 \textbf{Quick Refresher Before Starting:}
\textbf{Separation axioms} describe how well a topology can "separate" points and sets:
- \textbf{T$_1$:} For any two distinct points, each has a neighborhood not containing the other.
- \textbf{T$_2$ (Hausdorff):} Any two distinct points have DISJOINT neighborhoods. Most "nice" spaces are Hausdorff.
- \textbf{T$_3$ (Regular):} Points and closed sets can be separated by open sets.
- \textbf{T₄ (Normal):} Any two disjoint closed sets can be separated by open sets.
\textbf{Hierarchy:} T₄ $\Longrightarrow$ T$_3\Longrightarrow$ T$_2\Longrightarrow$ T$_1$. Every metric space is T₄.
\end{prereqbox}


\subsection{The Hierarchy }

\begin{verbatim}
T0 ──▶ T1 ──▶ T2 (Hausdorff) ──▶ T3 (Regular) ──▶ T4 (Normal)
                                    (each implies the previous)
\end{verbatim}

\begin{center}
\begin{tabular}{|l|l|l|}
\hline
\textbf{Axiom} & \textbf{Name} & \textbf{Condition} \\
\hline
T$_0$ & Kolmogorov & For any two distinct points, at least one has an open set not containing the other \\
\hline
T$_1$ & Fréchet & For any two distinct points, each has an open set not containing the other \\
\hline
T$_2$ & Hausdorff & Any two distinct points have disjoint open neighborhoods \\
\hline
T$_3$ & Regular & T$_1$ + point and closed set can be separated by open sets \\
\hline
T$_3$½ & Tychonoff & T$_1$ + point and closed set can be separated by continuous function \\
\hline
T₄ & Normal & T$_1$ + two disjoint closed sets can be separated by open sets \\
\hline
\end{tabular}
\end{center}

\subsubsection{ Memory Aid}
\textbf{T$_2$ (Hausdorff):} "Two points can be \textit{housed} apart" — put each in its own open house (neighborhood).

\subsubsection{Examples}
\begin{center}
\begin{tabular}{|l|l|}
\hline
\textbf{Space} & \textbf{Highest Separation Axiom} \\
\hline
$\mathbb{R}$ (standard) & T₄ (Normal) — all metric spaces are normal \\
\hline
$\mathbb{R}$ (cofinite) & T$_1$ but NOT T$_2$ \\
\hline
$\mathbb{R}$ (indiscrete) & T$_0$ only (fails even T$_1$) \\
\hline
Any metric space & T₄ \\
\hline
\end{tabular}
\end{center}

\subsubsection{ Urysohn's Lemma}
If X is normal and A, B are disjoint closed sets, then there exists a continuous function f: X  [0,1] with f(A) = {0} and f(B) = {1}.

\bigskip\noindent\rule{\textwidth}{0.4pt}\bigskip

\section{Unit 10: Separable and Second Countable Spaces}

\begin{prereqbox}
 \textbf{Quick Refresher Before Starting:}
\textbf{Countable set:} Can be listed as a sequence: a$_1$, a$_2$, a$_3$, ... Examples: $\mathbb{N}$, $\mathbb{Z}$, $\mathbb{Q}$ are countable. $\mathbb{R}$ is UNcountable.
\textbf{Dense set:} A $\subset$ X is dense if every open set contains a point of A. Example: $\mathbb{Q}$ is dense in $\mathbb{R}$ (every interval contains a rational).
\textbf{Separable:} X has a countable dense subset. $\mathbb{R}$ is separable ($\mathbb{Q}$ is countable and dense).
\textbf{Second countable:} X has a countable basis. Example: $\mathbb{R}$ has basis {(a,b) : a,b $\in\mathbb{Q}$} which is countable.
\end{prereqbox}


\subsection{Separable Spaces}
X is \textbf{separable} if it has a countable dense subset.

\subsubsection{Examples}
\begin{itemize}[leftmargin=*]
  \item $\mathbb{R}$ is separable ($\mathbb{Q}$ is countable and dense)
  \item ℓ$^2$ (separable Hilbert space) is separable
  \item ℓ^$\infty$ with sup norm is NOT separable
\end{itemize}

\subsection{Second Countable Spaces}
X is \textbf{second countable} if it has a countable basis for its topology.

\subsubsection{Examples}
\begin{itemize}[leftmargin=*]
  \item $\mathbb{R}$ is second countable (basis: intervals with rational endpoints)
  \item $\mathbb{R}^n$ is second countable
  \item Any uncountable set with discrete topology is NOT second countable
\end{itemize}

\subsubsection{Key Relationships }
\begin{verbatim}
Second Countable ──▶ First Countable
       │                    
       ▼                    
   Separable               
       │                    
       ▼                    
   Lindelöf
\end{verbatim}

\textbf{In metric spaces:} Separable $\Longleftrightarrow$ Second Countable $\Longleftrightarrow$ Lindelöf

\subsubsection{Lindelöf's Theorem}
Every second countable space is Lindelöf (every open cover has a countable subcover).

\bigskip\noindent\rule{\textwidth}{0.4pt}\bigskip


\textbf{Q1.} In the cofinite topology on an infinite set X, is X compact?
\textbf{Answer:} YES. Any open cover {U$\alpha$} — pick one U$\alpha_0$. Its complement X\U$\alpha_0$ is finite = {x$_1$,...,x$_n$}. For each x$_i$, pick U$\alpha_i$ containing x$_i$. Then {U$\alpha_0$, U$\alpha_1$,...,U$\alpha_n$} is a finite subcover.

\textbf{Q2.} Is $\mathbb{R}$ with the lower limit topology (basis = {[a,b)}) connected?
\textbf{Answer:} NO. [0,$\infty$) is both open (union of [0,n)) and closed (complement = $\cup_n$[−n,0) is open). So $\mathbb{R}$ splits as [0,$\infty$) $\cup$ (−$\infty$,0).

\textbf{Q3.} Product of [0,1] with itself uncountably many times — is it compact?
\textbf{Answer:} YES — by Tychonoff's theorem!

\textbf{Q4.} Give an example of a space that is T$_1$ but not T$_2$.
\textbf{Answer:} $\mathbb{R}$ with cofinite topology. Singletons are closed (T$_1$), but you can't separate two points by disjoint open sets (any two non-empty open sets intersect, since their complements are finite).

\bigskip\noindent\rule{\textwidth}{0.4pt}\bigskip

\begin{quote}
\textbf{ Recommended Reading:} Munkres "Topology", Simmons "Introduction to Topology and Modern Analysis"
\textbf{Next Subject:} [Advanced Complex Analysis ](./04_Advanced_Complex_Analysis.md)
\end{quote}

\newpage

\chapter{Topology --- Proofs \& Counterexamples}


\begin{quote}
\textbf{Supplement to 03_Topology.md — Counterexamples are GOLD for CSIR NET MCQs}
\end{quote}

\bigskip\noindent\rule{\textwidth}{0.4pt}\bigskip


\subsection{Proof 1: Heine-Borel Theorem ⚡⚡⚡}

\begin{tcolorbox}[colback=lightblue,colframe=thmblue,title=\textbf{Theorem},breakable,boxrule=1.5pt]
\textbf{Theorem:} A subset of $\mathbb{R}^n$ is compact iff it is closed and bounded.
\end{tcolorbox}

\begin{proof}
\textbf{Proof ($\Longleftarrow$ direction for $\mathbb{R}$, exam version):}
Let K $\subseteq$ [a,b] (closed and bounded). Let {U$\alpha$} be an open cover.
\end{proof}

Let S = {x $\in$ [a,b] : [a,x] can be covered by finitely many U$\alpha$}.
S is nonempty (a $\in$ S) and bounded above by b. Let c = sup S.

\textbf{Claim: c $\in$ S.} Since c $\in$ some U$\alpha_0$ (which is open), $\exists\delta$: (c−$\delta$,c+$\delta$) $\subseteq$ U$\alpha_0$. Since c = sup S, $\exists$x $\in$ S with x > c−$\delta$. So [a,x] has finite subcover. Adding U$\alpha_0$ covers [a,c]. So c $\in$ S. 

\textbf{Claim: c = b.} If c < b, then [a,c+$\varepsilon$] can also be finitely covered (extend using U$\alpha_0$). This contradicts c = sup S. So c = b $\in$ S. 

Therefore [a,b] is compact. Since K is closed subset of compact [a,b], K is compact. $\blacksquare$

\bigskip\noindent\rule{\textwidth}{0.4pt}\bigskip

\subsection{Proof 2: Continuous Image of Compact is Compact ⚡⚡⚡}

\begin{proof}

Let f: X  Y be continuous, X compact. Let {V$\alpha$} be open cover of f(X).
Then {f$^{-}^1$(V$\alpha$)} is open cover of X (by continuity). By compactness of X, $\exists$ finite subcover f$^{-}^1$(V$\alpha_1$),...,f$^{-}^1$(V$\alpha_n$). Then V$\alpha_1$,...,V$\alpha_n$ covers f(X). $\blacksquare$
\end{proof}

\bigskip\noindent\rule{\textwidth}{0.4pt}\bigskip

\subsection{Proof 3: Continuous Image of Connected is Connected ⚡⚡⚡}

\begin{proof}
\textbf{Proof (by contradiction):}
Let f: X  Y be continuous, X connected. Suppose f(X) = A $\cup$ B with A,B open in f(X), disjoint, nonempty.
Then X = f$^{-}^1$(A) $\cup$ f$^{-}^1$(B). These are open (continuity), disjoint, nonempty. This contradicts X connected. $\blacksquare$
\end{proof}

\bigskip\noindent\rule{\textwidth}{0.4pt}\bigskip

\subsection{Proof 4: Compact + Hausdorff $\Longrightarrow$ Normal ⚡⚡}

\begin{proof}
\textbf{Proof sketch:}
Step 1: In Hausdorff space, compact sets are closed.
Step 2: Given disjoint closed sets A, B in compact Hausdorff X — both A, B are compact.
Step 3: For each a $\in$ A, b $\in$ B: separate by open sets (Hausdorff). Use compactness of B to get finite subcover  separate a from all of B.
Step 4: Use compactness of A to get finite subcover  separate all of A from all of B. $\blacksquare$
\end{proof}

\bigskip\noindent\rule{\textwidth}{0.4pt}\bigskip

\subsection{Proof 5: Connected Subsets of $\mathbb{R}$ are Intervals ⚡⚡}

\begin{proof}

Let S $\subseteq\mathbb{R}$ be connected. Suppose a, b $\in$ S with a < c < b but c $\notin$ S.
Then S = (S $\cap$ (−$\infty$,c)) $\cup$ (S $\cap$ (c,$\infty$)) — two nonempty disjoint open (in S) sets. Contradicts connectedness. $\blacksquare$
\end{proof}

\bigskip\noindent\rule{\textwidth}{0.4pt}\bigskip


\subsection{The Big List}

\begin{center}
\begin{tabular}{|l|l|l|}
\hline
\textbf{Statement} & \textbf{True/False} & \textbf{Counterexample} \\
\hline
Connected $\Longrightarrow$ Path connected & \textbf{FALSE} & Topologist's sine curve: {(x,sin(1/x)):x>0} $\cup$ {0}$\times$[−1,1] \\
\hline
Compact $\Longrightarrow$ Closed & \textbf{FALSE} (in general) & {a,b} with indiscrete topology — compact but not closed \\
\hline
Compact $\Longrightarrow$ Closed & \textbf{TRUE in Hausdorff} & — \\
\hline
Closed + Bounded $\Longrightarrow$ Compact & \textbf{FALSE} (general metric) & In (0,1) with usual metric: (0,1) itself is closed (as full space) and bounded \\
\hline
Closed + Bounded $\Longrightarrow$ Compact & \textbf{TRUE in $\mathbb{R}^n$} & Heine-Borel \\
\hline
Product of compact = compact & \textbf{TRUE} & Tychonoff's theorem \\
\hline
Continuous bijection = homeomorphism & \textbf{FALSE} & f: [0,2$\pi$)  S$^1$ by f(t)=e^(it). Continuous bijection but f$^{-}^1$ not continuous \\
\hline
Continuous bijection = homeomorphism & \textbf{TRUE} & If domain is compact, codomain Hausdorff \\
\hline
Hausdorff $\Longrightarrow$ Normal & \textbf{FALSE} & $\mathbb{R}$ with K-topology (Sorgenfrey) \\
\hline
Metric $\Longrightarrow$ Normal & \textbf{TRUE} & Every metric space is normal \\
\hline
Separable $\Longrightarrow$ Second countable & \textbf{FALSE} (general) & Sorgenfrey line \\
\hline
Separable $\Longleftrightarrow$ Second countable & \textbf{TRUE in metric spaces} & — \\
\hline
$\mathbb{Q}$ is connected & \textbf{FALSE} & Split at $\sqrt{2}$ \\
\hline
$\mathbb{R}$ \ $\mathbb{Q}$ is connected & \textbf{TRUE} & Cannot split irrationals into two open sets \\
\hline
Subspace of Hausdorff is Hausdorff & \textbf{TRUE} & Subspace inherits separation \\
\hline
Subspace of compact is compact & \textbf{FALSE} & (0,1) $\subset$ [0,1] — not compact \\
\hline
Closed subspace of compact = compact & \textbf{TRUE} & — \\
\hline
Finite $\times$ compact = compact & \textbf{TRUE} & — \\
\hline
\end{tabular}
\end{center}

\subsection{5 KEY Spaces to Know Cold}

\textbf{1. Cofinite Topology on $\mathbb{R}$:} T$_1$, NOT T$_2$, connected, compact, every infinite set is dense.

\textbf{2. Discrete Topology on $\mathbb{N}$:} T₄, NOT compact (cover by singletons), NOT connected, second countable, separable.

\textbf{3. Sorgenfrey Line (Lower limit topology):} T₄, separable, NOT second countable, NOT connected ([0,$\infty$) is clopen), Lindelöf but NOT compact.

\textbf{4. Topologist's Sine Curve:} Connected, NOT path connected, NOT locally connected.

\textbf{5. Long Line:} Connected, locally compact, NOT compact, NOT metrizable.

\bigskip\noindent\rule{\textwidth}{0.4pt}\bigskip


\textbf{P1.} Is [0,1] $\cap\mathbb{Q}$ compact in $\mathbb{R}$? 
\textbf{Answer:} NO. Not closed in $\mathbb{R}$ (limit point √$\frac{2}{2}$ $\notin$ set). Heine-Borel fails.

\textbf{P2.} Is {(x,y) : x$^2$+y$^2\leq$ 1} \ {(0,0)} compact?
\textbf{Answer:} NO. Closed? No — (0,0) is a limit point not in the set. Not closed  not compact.

\textbf{P3.} Is $\mathbb{R}^2$ \ {line} connected?
\textbf{Answer:} NO. $\mathbb{R}^2$ minus a line has two components.

\textbf{P4.} Is $\mathbb{R}^2$ \ {point} connected?
\textbf{Answer:} YES. Can connect any two points by a path avoiding one point.

\textbf{P5.} Is $\mathbb{R}^3$ \ {line} connected?
\textbf{Answer:} YES. Can go "around" a line in 3D.

\textbf{P6.} In the cofinite topology on $\mathbb{Z}$, is {even integers} open?
\textbf{Answer:} NO. Its complement (odd integers) is infinite, not finite. So not cofinite, so not open.

\textbf{P7.} True/False: Intersection of two compact sets is compact (in Hausdorff space).
\textbf{Answer:} TRUE. Compact sets are closed in Hausdorff. Intersection of closed sets is closed. Closed subset of compact is compact.

\textbf{P8.} Is the product topology on $\mathbb{R}$^$\omega$ (countable product) metrizable?
\textbf{Answer:} YES — use the metric d(x,y) = $\Sigma$ min(|x$_n$−y$_n$|,1)/2$^n$.

\textbf{P9.} Every second countable space is separable. Prove.
\begin{proof}
Let {B$_n$} be countable basis. Pick x$_n\in$ B$_n$ for each n. Then {x$_n$} is dense: for any open U $\neq\emptyset$, U contains some B$_n$, which contains x$_n$. $\blacksquare$
\end{proof}

\textbf{P10.} Is GL(n,$\mathbb{R}$) (invertible matrices) connected?
\textbf{Answer:} NO. GL(n,$\mathbb{R}$) = {det>0} $\cup$ {det<0}. Two components. But GL⁺(n,$\mathbb{R}$) = {det>0} IS connected.

\newpage

\chapter{Advanced Complex Analysis}


\begin{quote}
\textbf{CSIR NET Priority: ⭐⭐⭐⭐ | Complex Analysis basics are heavily tested; advanced topics appear in Part C}
\end{quote}

\bigskip\noindent\rule{\textwidth}{0.4pt}\bigskip

\subsection{️ Subject Mind Map}

\begin{verbatim}
                   ADVANCED COMPLEX ANALYSIS
                           │
       ┌───────────┬───────┴────────┬───────────┐
       │           │                │           │
   Entire      Analytic        Harmonic     Advanced
   Functions   Continuation    Functions    Theorems
       │           │                │           │
   ┌───┴───┐   ┌──┴──┐        ┌───┴───┐   ┌───┴───┐
   Order   Canonical  Dirichlet  Green  Picard  Bloch
   Growth  Products   Problem    Func.  Theorems Schottky
\end{verbatim}

\bigskip\noindent\rule{\textwidth}{0.4pt}\bigskip

\section{Unit 1: Integral Functions (Entire Functions)}

\begin{prereqbox}
 \textbf{Quick Refresher Before Starting:}
\textbf{Complex number} z = a + bi where i$^2$ = -1. |z| = $\sqrt{a$^2$+b$^2$}$. Polar form: z = re^(i$\theta$).
\textbf{Analytic function:} f(z) is analytic if it's complex-differentiable. Must satisfy Cauchy-Riemann equations:
If f = u + iv, then $\partial$u/$\partial$x = $\partial$v/$\partial$y and $\partial$u/$\partial$y = -$\partial$v/$\partial$x.
\textbf{Entire function:} Analytic EVERYWHERE in $\mathbb{C}$. Examples: eᶻ, sin z, polynomials. NOT entire: 1/z (pole at 0).
\textbf{Taylor series:} f(z) = $\Sigma$ a$_n$z$^n$. Converges in a disk. Entire = infinite radius of convergence.
\end{prereqbox}



\subsection{Prerequisites Refresher}

An \textbf{analytic function} (holomorphic function) is one that is complex-differentiable in a region. Key facts from basic complex analysis:

\begin{itemize}[leftmargin=*]
  \item \textbf{Cauchy-Riemann equations:} If f(z) = u(x,y) + iv(x,y), then f is analytic iff $\partial$u/$\partial$x = $\partial$v/$\partial$y and $\partial$u/$\partial$y = −$\partial$v/$\partial$x
  \item \textbf{Cauchy's theorem:} ∮_C f(z)dz = 0 for analytic f on simply connected domain
  \item \textbf{Power series:} f(z) = $\Sigma$ a$_n$z$^n$ converges in a disk |z| < R
\end{itemize}

\subsection{What is an Entire Function?}

An \textbf{entire function} is a function that is analytic on ALL of $\mathbb{C}$ (the entire complex plane).

\subsubsection{Examples}
\begin{center}
\begin{tabular}{|l|l|l|}
\hline
\textbf{Function} & \textbf{Entire?} & \textbf{Why} \\
\hline
eᶻ & Yes & Power series converges everywhere \\
\hline
sin z, cos z & Yes & Power series converge everywhere \\
\hline
Polynomials p(z) & Yes & Analytic everywhere \\
\hline
1/z & No & Singularity at z = 0 \\
\hline
tan z & No & Poles at z = $\pi$/2 + n$\pi$ \\
\hline
log z & No & Branch point at z = 0 \\
\hline
\end{tabular}
\end{center}

\subsection{Liouville's Theorem }

\textbf{If f is entire and bounded, then f is constant.}

This is remarkably powerful! It immediately proves the Fundamental Theorem of Algebra.

\subsubsection{Proof Sketch}
By Cauchy's integral formula, |f'(z$_0$)| $\leq$ M/R for any R > 0. As R  $\infty$, f'(z$_0$) = 0 everywhere, so f is constant.

\subsubsection{ CSIR NET Application}
\textbf{Q:} f is entire with |f(z)| $\leq$ 3|z|$^2$ + 7 for all z. What can you say about f?

\textbf{Solution:} f(z)/z$^3$  0 as |z|  $\infty$. By an extended Liouville theorem, f must be a polynomial of degree $\leq$ 2. So f(z) = az$^2$ + bz + c.

\subsection{Growth and Order of Entire Functions}

The \textbf{order} of an entire function f is:

\textbf{$\rho$ = lim sup_{r$\infty$} [log log M(r)] / [log r]}

where M(r) = max_{|z|=r} |f(z)|.

\begin{center}
\begin{tabular}{|l|l|}
\hline
\textbf{Function} & \textbf{Order $\rho$} \\
\hline
Polynomial & 0 \\
\hline
eᶻ & 1 \\
\hline
e^(z$^2$) & 2 \\
\hline
sin z & 1 \\
\hline
e^(eᶻ) & $\infty$ \\
\hline
\end{tabular}
\end{center}

The \textbf{type} $\sigma$ of an entire function of order $\rho$ is:

\textbf{$\sigma$ = lim sup_{r$\infty$} log M(r) / r^$\rho$}

\begin{itemize}[leftmargin=*]
  \item $\sigma$ = 0: minimal type
  \item 0 < $\sigma$ < $\infty$: normal (mean) type
  \item $\sigma$ = $\infty$: maximal type
\end{itemize}

\subsubsection{Example}
f(z) = eᶻ: M(r) = eʳ, so log M(r) = r, log log M(r) = log r.
$\rho$ = lim (log r)/(log r) = 1. Type $\sigma$ = lim r/r$^1$ = 1. So eᶻ has order 1, type 1.

\bigskip\noindent\rule{\textwidth}{0.4pt}\bigskip

\section{Unit 2: Special Functions and Fundamental Theorems}

\begin{prereqbox}
 \textbf{Quick Refresher Before Starting:}
\textbf{Gamma function:} $\Gamma$(n) = (n-1)! for positive integers. $\Gamma$($\frac{1}{2}$) = √$\pi$. Defined by $\Gamma$(s) = $\int_0$^$\infty$ t^(s-1)e^(-t)dt.
\textbf{Functional equation:} $\Gamma$(s+1) = s$\cdot\Gamma$(s). This extends factorial to all complex numbers!
\textbf{Riemann zeta function:} $\zeta$(s) = $\Sigma$ 1/nˢ = 1 + $\frac{1}{2}$ˢ + $\frac{1}{3}$ˢ + ... Converges for Re(s) > 1.
\end{prereqbox}


\subsection{The Gamma Function}

\textbf{$\Gamma$(z) = $\int_0$^$\infty$ t^(z-1) e^(-t) dt} for Re(z) > 0

Key properties:
\begin{itemize}[leftmargin=*]
  \item $\Gamma$(n+1) = n! for positive integers
  \item $\Gamma$(z+1) = z$\Gamma$(z) (functional equation)
  \item $\Gamma$($\frac{1}{2}$) = √$\pi$
  \item Has poles at z = 0, −1, −2, ... (simple poles)
  \item Never zero: $\Gamma$(z) $\neq$ 0 for all z
\end{itemize}

\subsubsection{Euler's Reflection Formula}
\textbf{$\Gamma$(z)$\Gamma$(1−z) = $\pi$ / sin($\pi$z)}

\subsection{The Zeta Function}

\textbf{$\zeta$(s) = $\Sigma$_{n=1}^$\infty$ 1/nˢ} for Re(s) > 1

\begin{itemize}[leftmargin=*]
  \item Extends to a meromorphic function on all of $\mathbb{C}$
  \item Only pole: simple pole at s = 1 with residue 1
  \item Riemann Hypothesis: all non-trivial zeros have Re(s) = $\frac{1}{2}$ (unsolved!)
\end{itemize}

\subsection{Mittag-Leffler Theorem}

Given prescribed poles and principal parts, there exists a meromorphic function with exactly those singularities. This is the "partial fractions" theorem for meromorphic functions.

\subsubsection{Example}
To construct a function with simple poles at z = n (n $\in\mathbb{Z}$) with residue 1:
The function $\pi$ cot($\pi$z) has exactly these properties!

\bigskip\noindent\rule{\textwidth}{0.4pt}\bigskip

\section{Unit 3: Analytic Continuation}

\begin{prereqbox}
 \textbf{Quick Refresher Before Starting:}
\textbf{Power series} (Class 12): f(x) = $\Sigma$ a$_n$x$^n$ converges inside a disk of radius R (radius of convergence).
Outside R: diverges. On the boundary: depends on the specific series.
\textbf{Analytic continuation:} If f is defined by a power series in one disk, can we "extend" it to a larger region?
Example: f(z) = $\Sigma$ z$^n$ = 1/(1-z) converges only for |z| < 1, but 1/(1-z) makes sense everywhere except z = 1.
So 1/(1-z) is the analytic continuation of $\Sigma$ z$^n$ beyond the unit disk.
\end{prereqbox}


\subsection{The Big Idea}

If f is analytic on a region D$_1$ and g is analytic on D$_2$, with D$_1\cap$ D$_2\neq\emptyset$ and f = g on D$_1\cap$ D$_2$, then g is the \textbf{analytic continuation} of f to D$_2$.

\subsubsection{Identity Theorem }
If two analytic functions agree on a set with a limit point in their common domain, they agree everywhere on that domain.

\textbf{In plain English:} An analytic function is completely determined by its values on ANY set with a limit point.

\subsubsection{Example: Extending the Geometric Series}
f(z) = $\Sigma$_{n=0}^$\infty$ z$^n$ = 1/(1−z) is defined for |z| < 1.

But g(z) = 1/(1−z) is analytic on $\mathbb{C}$ \ {1}.

So g is the analytic continuation of f to $\mathbb{C}$ \ {1}.

\subsubsection{Example: The Gamma Function}
$\Gamma$(z) is initially defined for Re(z) > 0 by the integral. Using $\Gamma$(z) = $\Gamma$(z+1)/z, we analytically continue to Re(z) > −1 (except z = 0), then to Re(z) > −2, etc.

\subsection{Monodromy Theorem}
If f can be analytically continued along every path in a simply connected domain D, then the continuation defines a single-valued analytic function on D.

\subsection{Natural Boundary}
Some functions CANNOT be continued beyond their circle of convergence.

\begin{tcolorbox}[colback=examplebg!30,colframe=exampleborder,title=\textbf{Worked Example},breakable,boxrule=1pt]
\textbf{Example:} f(z) = $\Sigma$ z^(2$^n$) = z + z$^2$ + z$^4$ + z⁸ + ... has |z| = 1 as a \textbf{natural boundary} — every point on the unit circle is a singularity!
\end{tcolorbox}

\bigskip\noindent\rule{\textwidth}{0.4pt}\bigskip

\section{Unit 4: Advanced Topics and Harmonic Functions}

\begin{prereqbox}
 \textbf{Quick Refresher Before Starting:}
\textbf{Harmonic function:} u(x,y) satisfying Laplace's equation: $\partial^2$u/$\partial$x$^2$ + $\partial^2$u/$\partial$y$^2$ = 0 (written $\nabla^2$u = 0).
\textbf{Connection to analytic functions:} If f = u + iv is analytic, then BOTH u and v are harmonic!
\textbf{Harmonic conjugate:} Given u, find v such that u + iv is analytic. Use Cauchy-Riemann: vₓ = -uᵧ, vᵧ = uₓ.
Example: u = x$^2$ - y$^2$  uₓ = 2x, uᵧ = -2y  vₓ = 2y, vᵧ = 2x  v = 2xy. So f = z$^2$.
\end{prereqbox}


\subsection{Harmonic Functions}

A real-valued function u(x,y) is \textbf{harmonic} if:

\textbf{$\nabla^2$u = $\partial^2$u/$\partial$x$^2$ + $\partial^2$u/$\partial$y$^2$ = 0} (Laplace's equation)

\subsubsection{Connection to Analytic Functions }
If f(z) = u + iv is analytic, then BOTH u and v are harmonic. Moreover, v is the \textbf{harmonic conjugate} of u.

\subsubsection{Finding the Harmonic Conjugate}
Given u, find v using Cauchy-Riemann equations:
$\partial$v/$\partial$y = $\partial$u/$\partial$x and $\partial$v/$\partial$x = −$\partial$u/$\partial$y

\subsubsection{Example}
u(x,y) = x$^2$ − y$^2$ (harmonic: u_xx + u_yy = 2 + (−2) = 0 )

$\partial$v/$\partial$x = −$\partial$u/$\partial$y = 2y  v = 2xy + $\varphi$(y)
$\partial$v/$\partial$y = $\partial$u/$\partial$x = 2x  2x + $\varphi$'(y) = 2x  $\varphi$'(y) = 0  $\varphi$(y) = C

So v = 2xy + C, and f(z) = (x$^2$−y$^2$) + i(2xy) = z$^2$ + iC 

\subsection{Mean Value Property }
If u is harmonic in a disk and continuous on its closure:

\textbf{u(center) = ($\frac{1}{2}$$\pi$) $\int_0^2\pi$ u(center + re^(i$\theta$)) d$\theta$}

The value at the center equals the average over any circle around it!

\subsection{Maximum Principle }
A non-constant harmonic function on a connected open set has no interior maximum or minimum. Extremes occur only on the boundary.

\bigskip\noindent\rule{\textwidth}{0.4pt}\bigskip

\section{Unit 5: Harnack's Inequality and Dirichlet Problem}

\begin{prereqbox}
 \textbf{Quick Refresher Before Starting:}
\textbf{Mean value property:} For harmonic u: u(center) = average of u on any circle around the center.
\textbf{Maximum principle:} A harmonic function on a region attains its max/min on the BOUNDARY, never inside.
\textbf{Harnack's inequality} (new): Bounds harmonic functions from above and below on compact subsets.
\textbf{Green's function:} Solves the Dirichlet problem (find harmonic u with given boundary values).
\end{prereqbox}


\subsection{Harnack's Inequality}

If u is harmonic and non-negative in |z − z$_0$| $\leq$ R, then for |z − z$_0$| = r < R:

\textbf{(R−r)/(R+r) $\cdot$ u(z$_0$) $\leq$ u(z) $\leq$ (R+r)/(R−r) $\cdot$ u(z$_0$)}

\subsubsection{Intuition}
A positive harmonic function can't oscillate too wildly — it's "controlled" by its center value.

\subsection{Harnack's Theorem}
If {u$_n$} is a monotone increasing sequence of harmonic functions on a connected open set, then either u$_n\infty$ everywhere, or u$_n$ converges uniformly on compact subsets to a harmonic function.

\subsection{The Dirichlet Problem}

\textbf{Problem:} Given a region D and continuous boundary data f on $\partial$D, find a harmonic function u on D with u|_{$\partial$D} = f.

\subsubsection{Solution for the Disk (Poisson Integral Formula) }
For the unit disk, if f is continuous on |z| = 1:

\textbf{u(re^(i$\theta$)) = ($\frac{1}{2}$$\pi$) $\int_0^2\pi$ [(1−r$^2$) / (1 − 2r cos($\theta$−t) + r$^2$)] $\cdot$ f(e^(it)) dt}

The kernel P(r,$\theta$−t) = (1−r$^2$)/(1 − 2r cos($\theta$−t) + r$^2$) is the \textbf{Poisson kernel}.

\subsection{Green's Function}
G(z, z$_0$) for a region D is harmonic in D \ {z$_0$}, zero on $\partial$D, and G(z,z$_0$) + log|z−z$_0$| is harmonic near z$_0$.

For the unit disk: \textbf{G(z, z$_0$) = −log|(z−z$_0$)/(1−z̄$_0$z)|}

\bigskip\noindent\rule{\textwidth}{0.4pt}\bigskip

\section{Unit 6: Canonical Products, Jensen & Poisson-Jensen Formulas}

\begin{prereqbox}
 \textbf{Quick Refresher Before Starting:}
\textbf{Infinite product:} $\Pi$(1 + a$_n$) = (1+a$_1$)(1+a$_2$)(1+a$_3$)... Like infinite series but with multiplication.
Converges if $\Sigma$|a$_n$| converges (analogous to absolute convergence of series).
\textbf{Canonical product:} Builds an entire function with prescribed zeros {a$_n$}:
f(z) = $\Pi$ Eₚ(z/a$_n$) where Eₚ is an "elementary factor" that ensures convergence.
\end{prereqbox}


\subsection{Weierstrass Factorization Theorem }

Every entire function f with zeros at a$_1$, a$_2$, ... (with f(0) $\neq$ 0) can be written as:

\textbf{f(z) = e^(g(z)) $\cdot\Pi$_{n=1}^$\infty$ E_p$_n$(z/a$_n$)}

where E_p(z) = (1−z)exp(z + z$^2$/2 + ... + zᵖ/p) are \textbf{elementary factors} and g is entire.

\subsubsection{Elementary Factors}
\begin{itemize}[leftmargin=*]
  \item E$_0$(z) = 1 − z
  \item E$_1$(z) = (1−z)eᶻ
  \item E$_2$(z) = (1−z)e^(z+z$^2$/2)
\end{itemize}

\subsubsection{Example: sin($\pi$z)}
\textbf{sin($\pi$z) = $\pi$z $\cdot\Pi$_{n=1}^$\infty$ (1 − z$^2$/n$^2$)}

This expresses sin as an infinite product over its zeros z = 0, $\pm$1, $\pm$2, ...

\subsection{Canonical Product}
If {a$_n$} are the zeros, the \textbf{canonical product} is:

P(z) = $\Pi$ E_p(z/a$_n$) where p is chosen minimally to ensure convergence.

The \textbf{genus} of the canonical product is the smallest p that works.

\subsection{Jensen's Formula }

If f is analytic in |z| $\leq$ R, f(0) $\neq$ 0, and a$_1$,...,a$_n$ are the zeros in |z| < R, then:

\textbf{log|f(0)| + $\Sigma$ log(R/|a$_k$|) = ($\frac{1}{2}$$\pi$) $\int_0^2\pi$ log|f(Re^(i$\theta$))| d$\theta$}

\subsubsection{Intuition}
Jensen's formula connects the growth of f on circles to the number and location of its zeros inside.

\subsection{Poisson-Jensen Formula}
Generalizes Jensen's formula to give the value of log|f(z)| at any point inside the disk, not just the center.

\bigskip\noindent\rule{\textwidth}{0.4pt}\bigskip

\section{Unit 7: Hadamard's Three Circles Theorem}

\begin{prereqbox}
 \textbf{Quick Refresher Before Starting:}
\textbf{M(r) = max|f(z)| on |z| = r} — the maximum modulus on a circle of radius r.
As r grows, M(r) tells you how fast the entire function grows.
\textbf{Hadamard's three circles theorem:} log M(r) is a CONVEX function of log r.
This means: knowing M on two circles constrains M on any circle between them.
\end{prereqbox}


\subsection{Statement}

Let f be analytic in the annulus r$_1\leq$ |z| $\leq$ r$_3$. Let M(r) = max_{|z|=r} |f(z)|.

For r$_1$ < r$_2$ < r$_3$:

\textbf{log M(r$_2$) $\leq$ [log(r$_3$/r$_2$)/log(r$_3$/r$_1$)] $\cdot$ log M(r$_1$) + [log(r$_2$/r$_1$)/log(r$_3$/r$_1$)] $\cdot$ log M(r$_2$)}

Equivalently: \textbf{log M(r) is a convex function of log r}.

\subsubsection{Geometric Meaning}
\begin{verbatim}
log M(r)
    │         ╱
    │       ╱  ← actual curve lies BELOW the straight line
    │     ●    
    │   ╱  ╲
    │ ●      ● 
    └──────────── log r
     r₁  r₂  r₃
\end{verbatim}

The maximum modulus on intermediate circles can't exceed the linear interpolation between the maximum on the inner and outer circles.

\subsubsection{Application}
Used to prove results about the growth rate of entire functions and in the theory of Phragmén-Lindelöf.

\bigskip\noindent\rule{\textwidth}{0.4pt}\bigskip

\section{Unit 8: Borel, Hadamard Factorization, Range of Analytic Functions}

\begin{prereqbox}
 \textbf{Quick Refresher Before Starting:}
\textbf{Order of an entire function:} $\rho$ = lim sup (log log M(r))/(log r). Measures growth rate.
- Polynomials: $\rho$ = 0. eᶻ: $\rho$ = 1. e^(z$^2$): $\rho$ = 2.
\textbf{Hadamard factorization:} An entire function of finite order can be completely written as:
f(z) = z^m $\cdot$ e^(polynomial) $\cdot\Pi$ Eₚ(z/a$_n$) where a$_n$ are its zeros.
\end{prereqbox}


\subsection{Borel's Theorem}
An entire function of finite order $\rho$ cannot have more than a "sparse" set of values that it avoids.

\textbf{Borel's version of Picard:} If f is entire of order $\rho$, then f takes every value with at most one exception, and the exponent of convergence of the a-points equals $\rho$ for all a (with at most one exception).

\subsection{Hadamard Factorization Theorem }

Every entire function f of finite order $\rho$ with zeros a$_1$, a$_2$, ... can be written as:

\textbf{f(z) = z^m $\cdot$ e^(Q(z)) $\cdot\Pi$ E_p(z/a$_n$)}

where:
\begin{itemize}[leftmargin=*]
  \item m = order of zero at origin
  \item Q(z) is a polynomial of degree $\leq\rho$
  \item p $\leq\rho$ (the genus is at most the order)
\end{itemize}

\subsubsection{Example}
eᶻ − 1 has order 1, zeros at z = 2$\pi$in (n $\in\mathbb{Z}$, n $\neq$ 0):

eᶻ − 1 = z $\cdot$ e^(z/2) $\cdot\Pi$_{n$\neq$0} E$_1$(z/(2$\pi$in))

\subsection{Range of Analytic Functions}

\subsubsection{Open Mapping Theorem }
A non-constant analytic function is an open map (maps open sets to open sets).

\textbf{Consequence:} A non-constant analytic function cannot have a local maximum of |f| in the interior of its domain (Maximum Modulus Principle).

\bigskip\noindent\rule{\textwidth}{0.4pt}\bigskip

\section{Unit 9: Bloch, Schottky, Little Picard Theorem }

\begin{prereqbox}
 \textbf{Quick Refresher Before Starting:}
\textbf{Liouville's theorem:} Bounded entire function = constant. (If |f(z)| $\leq$ M for ALL z, then f is constant.)
\textbf{Picard goes much further:}
- \textbf{Little Picard:} A non-constant entire function takes EVERY value, with at most ONE exception.
- Example: eᶻ is entire and takes every value except 0. That one exception is the maximum allowed.
\textbf{Bloch's theorem:} Every analytic function on the unit disk contains a disk of a certain minimum size in its image.
\end{prereqbox}


\subsection{Bloch's Theorem}
There exists a universal constant B > 0 such that: if f is analytic in |z| < 1 with f'(0) = 1, then f(|z| < 1) contains a disk of radius B.

The exact value B $\geq$ 1/(2$\sqrt{3}$) is known; the sharp constant is still an open problem.

\subsection{Schottky's Theorem}
If f is analytic in |z| < 1 and omits the values 0 and 1, then |f(z)| is bounded by a quantity depending only on |f(0)| and |z| (not on f itself).

This is key to proving Picard's theorems.

\subsection{Little Picard Theorem }

\textbf{A non-constant entire function takes every complex value with at most ONE exception.}

\subsubsection{Examples}
\begin{itemize}[leftmargin=*]
  \item eᶻ omits 0 (eᶻ $\neq$ 0 for any z), but takes every other value. One exception — OK!
  \item sin z takes EVERY value (no exceptions)
  \item A polynomial of degree n takes every value (by FTA)
\end{itemize}

\subsubsection{ CSIR NET Question}
\textbf{Q:} True or False: There exists an entire function that omits both 0 and 1.
\textbf{A:} FALSE — by Little Picard, at most one value can be omitted.

\subsection{Great Picard Theorem}
In every neighborhood of an essential singularity, an analytic function takes every complex value (with at most one exception) infinitely many times.

\subsubsection{Example}
e^(1/z) near z = 0: for any w $\neq$ 0, the equation e^(1/z) = w has solutions z = 1/(log w + 2$\pi$in) accumulating at 0.

\bigskip\noindent\rule{\textwidth}{0.4pt}\bigskip

\section{Unit 10: Univalent Function Theory}

\begin{prereqbox}
 \textbf{Quick Refresher Before Starting:}
\textbf{Univalent (schlicht) function:} Analytic AND injective (one-to-one). f(z$_1$) = f(z$_2$) $\Longrightarrow$ z$_1$ = z$_2$.
\textbf{Schwarz lemma:} If f maps the unit disk to itself with f(0) = 0, then |f(z)| $\leq$ |z|.
\textbf{Conformal mapping:} Analytic + non-zero derivative. Preserves angles between curves.
\textbf{Riemann mapping theorem:} Any simply connected region ($\neq\mathbb{C}$) can be conformally mapped to the unit disk.
\end{prereqbox}


\subsection{Univalent (Schlicht) Functions}

A function f is \textbf{univalent} (schlicht) on a domain D if it is one-to-one (injective) on D.

The class \textbf{S} consists of all univalent functions f on the unit disk with f(0) = 0 and f'(0) = 1:

f(z) = z + a$_2$z$^2$ + a$_3$z$^3$ + ...

\subsection{Bieberbach Conjecture (de Branges' Theorem)}

\textbf{|a$_n$| $\leq$ n for all n $\geq$ 2}, with equality only for the Koebe function k(z) = z/(1−z)$^2$.

This was conjectured in 1916 and proved by de Branges in 1985.

\subsection{Key Results in S}

\subsubsection{Koebe $\frac{1}{4}$ Theorem}
If f $\in$ S, then f(D) contains the disk |w| < $\frac{1}{4}$.

\subsubsection{Growth and Distortion Theorems}
For f $\in$ S and |z| = r < 1:

\textbf{r/(1+r)$^2\leq$ |f(z)| $\leq$ r/(1−r)$^2$} (Growth)

\textbf{(1−r)/(1+r)$^3\leq$ |f'(z)| $\leq$ (1+r)/(1−r)$^3$} (Distortion)

\subsection{Area Theorem}
If g(z) = 1/z + $\Sigma$ b$_n$z$^n$ is univalent in |z| > 1, then $\Sigma$ n|b$_n$|$^2\leq$ 1.

\bigskip\noindent\rule{\textwidth}{0.4pt}\bigskip


\textbf{Q1.} If f is entire and |f(z)| $\leq$ e^|z| for all z, what is the maximum order of f?
\textbf{Answer:} Order $\rho\leq$ 1.

\textbf{Q2.} How many values can a non-constant entire function omit?
\textbf{Answer:} At most 1 (Little Picard Theorem).

\textbf{Q3.} Find the harmonic conjugate of u(x,y) = e$^x$ cos y.
\textbf{Solution:} v_x = −u_y = e$^x$ sin y  v = e$^x$ sin y + $\varphi$(y)
v_y = e$^x$ cos y + $\varphi$'(y) = u_x = e$^x$ cos y  $\varphi$'(y) = 0  v = e$^x$ sin y
So f(z) = e$^x$(cos y + i sin y) = eᶻ 

\textbf{Q4.} Express sin($\pi$z) as an infinite product.
\textbf{Answer:} sin($\pi$z) = $\pi$z $\cdot\Pi$_{n=1}^$\infty$ (1 − z$^2$/n$^2$)

\textbf{Q5.} State whether true/false: The Poisson integral of a continuous function on $\partial$D is harmonic in D.
\textbf{Answer:} TRUE.

\bigskip\noindent\rule{\textwidth}{0.4pt}\bigskip

\begin{quote}
\textbf{ Recommended Reading:} Ahlfors "Complex Analysis", Conway "Functions of One Complex Variable II"
\textbf{Next Subject:} [Advanced Differential Equations ](./05_Advanced_Differential_Equations.md)
\end{quote}

\newpage

\chapter{Complex Analysis --- Proofs \& Problems}


\begin{quote}
\textbf{Supplement to 04_Advanced_Complex_Analysis.md}
\end{quote}

\bigskip\noindent\rule{\textwidth}{0.4pt}\bigskip


\subsection{Proof 1: Liouville's Theorem ⚡⚡⚡}

\begin{tcolorbox}[colback=lightblue,colframe=thmblue,title=\textbf{Theorem},breakable,boxrule=1.5pt]
\textbf{Theorem:} If f is entire and bounded (|f(z)| $\leq$ M for all z), then f is constant.
\end{tcolorbox}

\begin{proof}

By Cauchy's integral formula for derivatives:
f'(z$_0$) = ($\frac{1}{2}$$\pi$i) ∮_{|z-z$_0$|=R} f(z)/(z-z$_0$)$^2$ dz
\end{proof}

Taking modulus:
\begin{center}
\begin{tabular}{|l|}
\hline
\textbf{f'(z$_0$)} \\
\hline
\end{tabular}
\end{center}

This holds for ALL R > 0. Letting R  $\infty$: |f'(z$_0$)| = 0 for every z$_0$.

So f' ≡ 0 everywhere, meaning f is constant. $\blacksquare$

\textbf{Application: Fundamental Theorem of Algebra}
If p(z) is a non-constant polynomial with no zeros, then 1/p(z) is entire and bounded (since |p(z)|  $\infty$ as |z|  $\infty$). By Liouville, 1/p is constant  p is constant. Contradiction! So every non-constant polynomial has a zero. $\blacksquare$

\bigskip\noindent\rule{\textwidth}{0.4pt}\bigskip

\subsection{Proof 2: Maximum Modulus Principle ⚡⚡⚡}

\begin{tcolorbox}[colback=lightblue,colframe=thmblue,title=\textbf{Theorem},breakable,boxrule=1.5pt]
\textbf{Theorem:} If f is analytic and non-constant on a connected open set D, then |f| has no local maximum in D.
\end{tcolorbox}

\begin{proof}

Suppose |f| attains maximum at z$_0\in$ D. By mean value property:
f(z$_0$) = ($\frac{1}{2}$$\pi$) $\int_0^2\pi$ f(z$_0$ + re^(i$\theta$)) d$\theta$
\end{proof}

So |f(z$_0$)| $\leq$ ($\frac{1}{2}$$\pi$) $\int_0^2\pi$ |f(z$_0$ + re^(i$\theta$))| d$\theta\leq$ max|f| on circle = |f(z$_0$)|

Equality throughout forces |f(z$_0$ + re^(i$\theta$))| = |f(z$_0$)| for all $\theta$. This holds for all small r, so |f| is constant near z$_0$. By the identity theorem, f is constant on D. Contradiction! $\blacksquare$

\bigskip\noindent\rule{\textwidth}{0.4pt}\bigskip

\subsection{Proof 3: Little Picard Theorem (Statement + Proof Idea) ⚡⚡⚡}

\begin{tcolorbox}[colback=lightblue,colframe=thmblue,title=\textbf{Theorem},breakable,boxrule=1.5pt]
\textbf{Theorem:} A non-constant entire function takes every value with at most one exception.
\end{tcolorbox}

\begin{proof}
\textbf{Proof idea (using Schottky's theorem):}
Suppose f omits both a and b. WLOG a=0, b=1 (compose with Möbius transformation).
Then g(z) = f(Rz) (for large R) also omits 0 and 1, with g(0) = f(0).
By Schottky's theorem, |g| is bounded on |z| $\leq$ $\frac{1}{2}$, hence |f| is bounded on |z| $\leq$ R/2.
Since R is arbitrary, f is bounded on $\mathbb{C}$. By Liouville, f is constant. Contradiction! $\blacksquare$
\end{proof}

\bigskip\noindent\rule{\textwidth}{0.4pt}\bigskip

\subsection{Proof 4: Identity Theorem ⚡⚡}

\begin{tcolorbox}[colback=lightblue,colframe=thmblue,title=\textbf{Theorem},breakable,boxrule=1.5pt]
\textbf{Theorem:} If f, g are analytic on connected D and f = g on a set S with a limit point in D, then f = g on all of D.
\end{tcolorbox}

\begin{proof}

Let h = f − g. Then h has a zero set Z with a limit point z$_0\in$ D.
Since h is analytic at z$_0$ and h vanishes on a sequence converging to z$_0$, by the power series representation, h ≡ 0 in a neighborhood of z$_0$.
The set {z $\in$ D : h ≡ 0 near z} is open (by analyticity) and closed (by continuity).
Since D is connected and this set is nonempty, it equals D. So h ≡ 0 on D. $\blacksquare$
\end{proof}

\bigskip\noindent\rule{\textwidth}{0.4pt}\bigskip

\subsection{Proof 5: Open Mapping Theorem ⚡⚡}

\begin{tcolorbox}[colback=lightblue,colframe=thmblue,title=\textbf{Theorem},breakable,boxrule=1.5pt]
\textbf{Theorem:} A non-constant analytic function maps open sets to open sets.
\end{tcolorbox}

\begin{proof}
\textbf{Proof sketch:}
Let w$_0$ = f(z$_0$) and U a neighborhood of z$_0$. Need to show w$_0$ is interior to f(U).
Since f − w$_0$ has an isolated zero at z$_0$, on a small circle C around z$_0$: |f(z) − w$_0$| $\geq\delta$ > 0.
For |w − w$_0$| < $\delta$: by Rouché's theorem, f(z) − w and f(z) − w$_0$ have the same number of zeros inside C.
So f(z) = w has a solution inside C $\subset$ U. Hence B(w$_0$, $\delta$) $\subset$ f(U). $\blacksquare$
\end{proof}

\bigskip\noindent\rule{\textwidth}{0.4pt}\bigskip


\subsection{Liouville Applications — 5 Variations}

\textbf{Ex 1:} f entire, |f(z)| $\leq$ 5|z|$^3$ + 7 for all z. Then f is a polynomial of degree $\leq$ 3.
\textit{Proof:} g(z) = f(z)/z$^4$  0 as |z|  $\infty$. By extended Liouville, f⁽$^4$⁾ ≡ 0, so f is polynomial of deg $\leq$ 3.

\textbf{Ex 2:} f entire, Re(f) $\leq$ 100 for all z. Then f is constant.
\textit{Proof:} g(z) = e^(f(z)). Then |g| = e^(Re f) $\leq$ e$^{10}^0$. So g is entire and bounded  g constant  f constant.

\textbf{Ex 3:} f entire, |f(z)| $\geq$ 1 for all z. Then f is constant.
\textit{Proof:} 1/f is entire with |1/f| $\leq$ 1. By Liouville, 1/f constant  f constant.

\textbf{Ex 4:} f entire, f maps $\mathbb{R}$ to $\mathbb{R}$ and maps i$\mathbb{R}$ to i$\mathbb{R}$. Show f(z) = f(z̄)̄ and the Taylor coefficients are real.
\textit{Proof:} g(z) = f(z̄)̄ is also entire and agrees with f on $\mathbb{R}$. By identity theorem, f = g on $\mathbb{C}$.

\textbf{Ex 5 (NET style):} f entire with f(n) = n/(n+1) for n = 1,2,3,... What is f?
\textit{Solution:} The function g(z) = z/(z+1) agrees with f at z = 1,2,3,... which has limit point $\infty$. But both are analytic where defined. Since f is entire and g has a pole at z = −1, we need: if f is entire and f(n) = n/(n+1), then f(z) = z/(z+1)... but this isn't entire! So consider h(z) = f(z) − z/(z+1). This has zeros at z = 1,2,3,... However f is entire, so (z+1)f(z) − z is entire with zeros at n=1,2,3,... and also at z=−1: (0)f(−1)−(−1) = 1 $\neq$ 0 unless f(−1) is defined differently. Actually f(z) = 1 − 1/(z+1) = z/(z+1) is NOT entire. So no entire function satisfies this! The question tests understanding that entire functions can't behave like rational functions.

\subsection{Harmonic Conjugate — 4 Examples}

\textbf{Ex 1:} u = x$^3$ − 3xy$^2$. Find v.
vₓ = −uᵧ = 6xy  v = 3x$^2$y + $\varphi$(y)
vᵧ = uₓ = 3x$^2$ − 3y$^2$  3x$^2$ + $\varphi$'(y) = 3x$^2$ − 3y$^2\varphi$'(y) = −3y$^2\varphi$ = −y$^3$
v = 3x$^2$y − y$^3$. So f(z) = z$^3$ + C. 

\textbf{Ex 2:} u = log(x$^2$+y$^2$) (= 2 log|z|). Harmonic for (x,y) $\neq$ (0,0).
v = 2 arctan(y/x). f(z) = 2 log z.

\textbf{Ex 3:} u = e$^x$ cos y. v = e$^x$ sin y. f(z) = eᶻ. (Standard result.)

\textbf{Ex 4:} u = x/(x$^2$+y$^2$). Find v.
uₓ = (y$^2$−x$^2$)/(x$^2$+y$^2$)$^2$, uᵧ = −2xy/(x$^2$+y$^2$)$^2$
vₓ = 2xy/(x$^2$+y$^2$)$^2$, integrate: v = −y/(x$^2$+y$^2$) + $\varphi$(y)
Check vᵧ = uₓ: confirms $\varphi$'(y) = 0. v = −y/(x$^2$+y$^2$).
f(z) = 1/z. 

\bigskip\noindent\rule{\textwidth}{0.4pt}\bigskip


\textbf{P1.} f is entire with |f(z)| $\leq$ e^(|Re z|). What can you conclude about the order of f?
\textbf{Answer:} log M(r) $\leq$ r (since |Re z| $\leq$ |z| = r). Order $\rho\leq$ 1.

\textbf{P2.} Find the order of f(z) = sin z.
\textbf{Answer:} |sin z| $\leq$ e^|z|. And sin(iy) = i sinh(y) grows like e^|y|/2. So $\rho$ = 1, type $\sigma$ = 1.

\textbf{P3.} How many zeros does z⁷ − 5z$^3$ + 12 have in |z| < 1?
\textbf{Answer:} On |z|=1: |12| = 12 > |z⁷−5z$^3$| $\leq$ 1+5 = 6. By Rouché: same as 12 (constant)  \textbf{0 zeros}.

\textbf{P4.} How many zeros does z$^4$ − 6z + 3 have in |z| < 2?
\textbf{Answer:} On |z|=2: |z$^4$| = 16 > |−6z+3| $\leq$ 12+3 = 15. By Rouché: same as z$^4$  \textbf{4 zeros}.

\textbf{P5.} True/False: An analytic function that maps the unit disk to itself and fixes two points is the identity.
\textbf{Answer:} TRUE (by Schwarz-Pick lemma — if f fixes two points, it's a rotation that fixes them, hence identity).

\textbf{P6.} Prove that a non-constant polynomial takes every complex value.
\begin{proof}
If p(z) $\neq$ w$_0$ for all z, then 1/(p(z)−w$_0$) is entire and bounded (since |p(z)|  $\infty$). By Liouville, constant. Contradiction.
\end{proof}

\textbf{P7.} Find $\int_0^2\pi$ d$\theta$/(2+cos$\theta$) using residue calculus.
\textbf{Answer:} z = e^(i$\theta$), cos$\theta$ = (z+z$^{-}^1$)/2. Integral = ∮ dz/[iz(2+(z+z$^{-}^1$)/2)] = ∮ 2dz/[i(z$^2$+4z+1)].
Poles at z = −2$\pm$$\sqrt{3}$. Only z = −2+$\sqrt{3}$ $\in$ unit disk. Residue = 1/$\sqrt{3}$. Answer = 2$\pi$/$\sqrt{3}$.

\textbf{P8.} Construct an entire function with zeros exactly at z = n$^2$ (n = 1,2,3,...).
\textbf{Answer:} Use Weierstrass: f(z) = $\Pi$ E$_0$(z/n$^2$) = $\Pi$(1−z/n$^2$) converges since $\Sigma$1/n$^2$ < $\infty$.

\newpage

\chapter{Advanced Differential Equations}


\begin{quote}
\textbf{CSIR NET Priority: ⭐⭐⭐⭐ | ODE theory appears in Part B; stability/qualitative theory in Part C}
\end{quote}

\bigskip\noindent\rule{\textwidth}{0.4pt}\bigskip

\subsection{️ Subject Mind Map}

\begin{verbatim}
              ADVANCED DIFFERENTIAL EQUATIONS
                         │
       ┌─────────┬───────┴───────┬──────────┐
       │         │               │          │
   Existence  Systems       Nonlinear   Stability
       │         │               │          │
   ┌───┴───┐  ┌──┴──┐      ┌───┴───┐  ┌───┴───┐
   ε-Approx Linear  Nonlin  Poincaré  Lyapunov
   Picard   Systems Systems Bendixson Critical
   Uniquen. Matrix          Limit     Points
              Exp.          Cycles
\end{verbatim}

\bigskip\noindent\rule{\textwidth}{0.4pt}\bigskip

\section{Unit 1: $\varepsilon$-Approximate Solutions}

\begin{prereqbox}
 \textbf{Quick Refresher Before Starting:}
\textbf{ODE} (Ordinary Differential Equation): An equation involving a function y and its derivatives y', y'', etc.
\textbf{Example:} y' = 2y. Solution: y = Ce$^{2x}$ (check: (Ce$^{2x}$)' = 2Ce$^{2x}$ = 2y ).
\textbf{Separation of variables} (Class 12): If dy/dx = f(x)g(y), then $\int$dy/g(y) = $\int$f(x)dx.
\textbf{$\varepsilon$-approximate solution:} A function $\varphi$ with |$\varphi$' - f(t,$\varphi$)| < $\varepsilon$ — it "almost" satisfies the ODE.
\end{prereqbox}



\subsection{Setting the Stage}

We study the initial value problem (IVP):

\textbf{y' = f(t, y), y(t$_0$) = y$_0$}

Before proving solutions exist, we need the concept of "almost-solutions."

\subsection{Definition}

A continuous function $\varphi$(t) is an \textbf{$\varepsilon$-approximate solution} of y' = f(t,y) on [t$_0$, t$_0$+a] if:

\begin{itemize}
  \item $\varphi$(t$_0$) = y$_0$
  \item $\varphi$'(t) exists almost everywhere
  \item |$\varphi$'(t) − f(t, $\varphi$(t))| < $\varepsilon$ wherever $\varphi$'(t) exists
\end{itemize}

\textbf{In plain English:} $\varphi$ almost satisfies the equation — the error is less than $\varepsilon$ everywhere.

\subsubsection{Example}
For y' = y, y(0) = 1, the exact solution is e$^t$.

The function $\varphi$(t) = 1 + t + t$^2$/2 is an $\varepsilon$-approximate solution for small t because:
$\varphi$'(t) = 1 + t, and f(t,$\varphi$(t)) = 1 + t + t$^2$/2

\begin{center}
\begin{tabular}{|l|l|l|}
\hline
\textbf{$\varphi$'(t) − f(t,$\varphi$(t))} & \textbf{=} & \textbf{−t$^2$/2} \\
\hline
\end{tabular}
\end{center}

\subsection{Cauchy-Euler Polygonal Method}

Construct $\varepsilon$-approximate solutions by piecewise linear interpolation:

\begin{verbatim}
Given: y' = f(t,y), y(t₀) = y₀
Step 1: From (t₀, y₀), draw line with slope f(t₀, y₀) to t₁ = t₀ + h
Step 2: At (t₁, y₁), draw line with slope f(t₁, y₁) to t₂ = t₁ + h
... Continue
\end{verbatim}

This gives a piecewise linear function. As h  0, these converge to the actual solution (under appropriate conditions).

\bigskip\noindent\rule{\textwidth}{0.4pt}\bigskip

\section{Unit 2: Basic Existence and Uniqueness Theorems}

\begin{prereqbox}
 \textbf{Quick Refresher Before Starting:}
\textbf{Existence:} Does a solution exist? \textbf{Uniqueness:} Is there exactly one solution?
\textbf{Lipschitz condition:} |f(t,y$_1$) - f(t,y$_2$)| $\leq$ L|y$_1$-y$_2$|. It means f doesn't change too wildly in y.
\textbf{Picard iteration:} Start with $\varphi_0$ = y$_0$, then $\varphi_n$₊$_1$(t) = y$_0$ + $\int$f(s,$\varphi_n$(s))ds. Keep iterating — the sequence converges to the unique solution.
Example: y' = y, y(0) = 1. $\varphi_0$ = 1, $\varphi_1$ = 1+t, $\varphi_2$ = 1+t+t$^2$/2, ...  e$^t$ 
\end{prereqbox}


\subsection{Peano's Existence Theorem}

\textbf{If f(t,y) is continuous on a rectangle R = {|t−t$_0$| $\leq$ a, |y−y$_0$| $\leq$ b}, then the IVP has at least one solution on [t$_0$−$\alpha$, t$_0$+$\alpha$]} where $\alpha$ = min(a, b/M), M = max|f| on R.

\textbf{Note:} Continuity gives existence but NOT uniqueness!

\subsubsection{Example of Non-Uniqueness}
y' = y^($\frac{2}{3}$), y(0) = 0

Both y(t) = 0 and y(t) = (t/3)$^3$ are solutions! (f is continuous but not Lipschitz at y = 0)

\subsection{Picard-Lindelöf Theorem  (CSIR NET HIGH PRIORITY)}

\textbf{If f is continuous AND satisfies a Lipschitz condition:}

\textbf{|f(t,y$_1$) − f(t,y$_2$)| $\leq$ L|y$_1$ − y$_2$|}

\textbf{then the IVP has a UNIQUE solution.}

\subsubsection{Picard's Iteration Method}

Starting with $\varphi_0$(t) = y$_0$, define:

\textbf{$\varphi_n$₊$_1$(t) = y$_0$ + $\int$_{t$_0$}^{t} f(s, $\varphi_n$(s)) ds}

These converge uniformly to the unique solution.

\subsubsection{Worked Example }
Solve y' = y, y(0) = 1 by Picard iteration:

\begin{verbatim}
φ₀(t) = 1
φ₁(t) = 1 + ∫₀ᵗ 1 ds = 1 + t
φ₂(t) = 1 + ∫₀ᵗ (1+s) ds = 1 + t + t²/2
φ₃(t) = 1 + ∫₀ᵗ (1+s+s²/2) ds = 1 + t + t²/2 + t³/6
...
φₙ(t) = Σₖ₌₀ⁿ tᵏ/k! → eᵗ  ✓
\end{verbatim}

\subsection{Gronwall's Inequality }

If u(t) $\leq\alpha$ + $\int$_{t$_0$}^{t} $\beta$(s)u(s)ds where u, $\beta\geq$ 0, then:

\textbf{u(t) $\leq\alpha\cdot$ exp($\int$_{t$_0$}^{t} $\beta$(s) ds)}

Used extensively in proving uniqueness and continuous dependence.

\bigskip\noindent\rule{\textwidth}{0.4pt}\bigskip

\section{Unit 3: Systems of Differential Equations}

\begin{prereqbox}
 \textbf{Quick Refresher Before Starting:}
\textbf{System of ODEs:} Multiple equations coupled together.
x' = 2x + y
y' = x - y
Written as: X' = AX where X = (x,y) and A = [[2,1],[1,-1]].
\textbf{Why matrices matter:} Solving X' = AX requires eigenvalues of A!
\end{prereqbox}


\subsection{From Single ODE to Systems}

An nth-order ODE can always be converted to a first-order SYSTEM:

\textbf{y⁽$^n$⁾ = f(t, y, y', ..., y⁽$^n^{-}^1$⁾)}

Setting x$_1$ = y, x$_2$ = y', ..., x$_n$ = y⁽$^n^{-}^1$⁾:

\begin{verbatim}
x₁' = x₂
x₂' = x₃
...
xₙ' = f(t, x₁, x₂, ..., xₙ)
\end{verbatim}

In vector form: \textbf{x' = F(t, x)} where x = (x$_1$,...,x$_n$)ᵀ

\subsubsection{Example}
y'' + 3y' + 2y = 0 becomes:

\begin{verbatim}
x₁' = x₂
x₂' = −2x₁ − 3x₂

Matrix form: x' = |0   1| x = Ax
                   |−2 −3|
\end{verbatim}

\subsection{Linear Systems: x' = A(t)x + g(t)}

\subsubsection{Homogeneous case: x' = Ax (constant coefficients)}

\textbf{Solution: x(t) = eᴬ$^t$ x$_0$}

where the \textbf{matrix exponential} eᴬ$^t$ = I + At + A$^2$t$^2$/2! + A$^3$t$^3$/3! + ...

\subsubsection{Fundamental Matrix}
A matrix $\Phi$(t) whose columns are n linearly independent solutions. Then:

\textbf{General solution: x(t) = $\Phi$(t)c} where c is an arbitrary constant vector.

\textbf{Wronskian:} W(t) = det($\Phi$(t)). By Abel's formula:

W(t) = W(t$_0$) $\cdot$ exp($\int$_{t$_0$}^{t} tr(A(s)) ds)

\bigskip\noindent\rule{\textwidth}{0.4pt}\bigskip

\section{Unit 4: Solutions of Systems}

\begin{prereqbox}
 \textbf{Quick Refresher Before Starting:}
\textbf{Matrix exponential:} e^(At) = I + At + A$^2$t$^2$/2! + A$^3$t$^3$/3! + ... (like the scalar e$^x$ but for matrices).
\textbf{Solution of X' = AX:} X(t) = e^(At) $\cdot$ X(0).
\textbf{If A is diagonalizable} (A = PDP$^{-}^1$): e^(At) = P $\cdot$ diag(e^($\lambda_1$t), ..., e^($\lambda_n$t)) $\cdot$ P$^{-}^1$.
This is why eigenvalues determine everything about the system's behavior!
\end{prereqbox}


\subsection{Matrix Exponential Method }

For x' = Ax with constant A, the key is computing eᴬ$^t$.

\subsubsection{Case 1: A is Diagonalizable}
If A = PDP$^{-}^1$ where D = diag($\lambda_1$,...,$\lambda_n$), then:

\textbf{eᴬ$^t$ = P $\cdot$ diag(e^($\lambda_1$t),...,e^($\lambda_n$t)) $\cdot$ P$^{-}^1$}

\subsubsection{Example}
\begin{verbatim}
A = |1  1|, eigenvalues λ = 0, 2
    |1  1|

P = | 1  1|, D = |0  0|
    |−1  1|      |0  2|

eᴬᵗ = P |1    0  | P⁻¹ = ½|1+e²ᵗ   e²ᵗ−1|
         |0   e²ᵗ|         |e²ᵗ−1   1+e²ᵗ|
\end{verbatim}

\subsubsection{Case 2: A has Repeated Eigenvalues (Jordan Form)}
If A = PJP$^{-}^1$ where J = Jordan form, compute eᴶ$^t$ block by block:

\begin{verbatim}
e^(Jₖ(λ)t) = e^(λt) |1   t   t²/2!  ...|
                      |0   1   t      ...|
                      |0   0   1      ...|
                      |⋮               ⋱ |
\end{verbatim}

\subsubsection{Case 3: Complex Eigenvalues}
If $\lambda$ = $\alpha\pm\beta$i, real solutions involve e^($\alpha$t)cos($\beta$t) and e^($\alpha$t)sin($\beta$t).

\subsubsection{Variation of Parameters (Non-homogeneous)}
For x' = Ax + g(t):

\textbf{x(t) = eᴬ$^t$x$_0$ + $\int_0^t$ eᴬ⁽$^t^{-}$ˢ⁾ g(s) ds}

\bigskip\noindent\rule{\textwidth}{0.4pt}\bigskip

\section{Unit 5: Nonlinear Differential Systems}

\begin{prereqbox}
 \textbf{Quick Refresher Before Starting:}
\textbf{Linear system:} X' = AX. Behavior fully determined by eigenvalues of A.
\textbf{Nonlinear system:} X' = F(X) where F is NOT just a matrix times X.
Example: x' = x - x$^3$ (the -x$^3$ term is nonlinear).
\textbf{Linearization:} Near an equilibrium point x$_0$ (where F(x$_0$) = 0), approximate: F(x) $\approx$ DF(x$_0$)$\cdot$(x - x$_0$).
So near equilibrium, the nonlinear system behaves like a linear one!
\end{prereqbox}


\subsection{Phase Plane Analysis}

For a 2D autonomous system:
\begin{verbatim}
x' = f(x, y)
y' = g(x, y)
\end{verbatim}

\textbf{Equilibrium points} (critical points): where f(x,y) = 0 AND g(x,y) = 0.

\subsubsection{Linearization at Equilibrium }
Near an equilibrium (x$_0$, y$_0$), the system behaves like:

\textbf{u' = J $\cdot$ u} where J is the Jacobian matrix evaluated at (x$_0$, y$_0$):

\begin{verbatim}
J = |∂f/∂x  ∂f/∂y|
    |∂g/∂x  ∂g/∂y| evaluated at (x₀, y₀)
\end{verbatim}

\subsubsection{Example: Predator-Prey (Lotka-Volterra)}
\begin{verbatim}
x' = x(a − by)    (prey)
y' = y(−c + dx)   (predator)
\end{verbatim}

Equilibria: (0,0) and (c/d, a/b)

At (c/d, a/b): J = |    0      −bc/d|, eigenvalues = $\pm$i$\sqrt{ac}$
                     |  ad/b      0  |

Pure imaginary eigenvalues  center (periodic orbits around equilibrium).

\bigskip\noindent\rule{\textwidth}{0.4pt}\bigskip

\section{Unit 6: Poincaré-Bendixson Theory }

\begin{prereqbox}
 \textbf{Quick Refresher Before Starting:}
\textbf{Phase plane:} Plot (x,y) as t varies. The trajectory traces a curve in 2D.
\textbf{Limit cycle:} A closed trajectory that nearby solutions spiral toward (or away from).
\textbf{Poincaré-Bendixson theorem} (key result): In 2D, if a trajectory stays in a bounded region with no equilibria, it must approach a limit cycle.
This is one of the few tools for proving periodic solutions exist in nonlinear systems!
\end{prereqbox}


\subsection{Limit Sets}

For a trajectory $\gamma$(t) in $\mathbb{R}^2$:
\begin{itemize}[leftmargin=*]
  \item \textbf{$\omega$-limit set:} points that $\gamma$(t) approaches as t  +$\infty$
  \item \textbf{$\alpha$-limit set:} points that $\gamma$(t) approaches as t  −$\infty$
\end{itemize}

\subsection{Poincaré-Bendixson Theorem  (VERY IMPORTANT)}

\textbf{If a trajectory of a 2D system stays in a bounded region containing no equilibrium points, then its $\omega$-limit set is a periodic orbit (limit cycle).}

\subsubsection{Why Only 2D?}
In 3D, trajectories can exhibit chaos (strange attractors). The Poincaré-Bendixson theorem says this CANNOT happen in 2D — the only long-term behaviors are:
\begin{itemize}
  \item Approach an equilibrium
  \item Approach a periodic orbit
  \item Be a periodic orbit
\end{itemize}

\subsubsection{Application: Proving Existence of Limit Cycles}
\textbf{Strategy:} Find a "trapping region" (bounded, no equilibria inside)  limit cycle must exist.

\subsubsection{Example: Van der Pol Oscillator}
x'' − $\mu$(1−x$^2$)x' + x = 0 ($\mu$ > 0)

As system: x' = y, y' = $\mu$(1−x$^2$)y − x

For large $\mu$, one can construct an annular trapping region  unique stable limit cycle exists.

\subsection{Bendixson's Criterion (Negative Result)}
If $\partial$f/$\partial$x + $\partial$g/$\partial$y is of constant sign in a simply connected region D, then there are no periodic orbits in D.

\bigskip\noindent\rule{\textwidth}{0.4pt}\bigskip

\section{Unit 7: Critical Points, Stability, and Lyapunov Methods }

\begin{prereqbox}
 \textbf{Quick Refresher Before Starting:}
\textbf{Equilibrium:} A point x$_0$ where F(x$_0$) = 0 (the system stays still).
\textbf{Stable:} Small perturbations stay small. \textbf{Unstable:} Small perturbations grow.
\textbf{Asymptotically stable:} Perturbations not only stay small but decay to zero.
\textbf{Lyapunov's method:} Find an "energy function" V(x) > 0 that DECREASES along solutions (V̇ $\leq$ 0).
If such V exists  stable. If V̇ < 0  asymptotically stable.
Think of V like a ball on a hill: if V always decreases, the ball rolls to the bottom (equilibrium).
\end{prereqbox}


\subsection{Classification of Critical Points (2D Linear Systems)}

For x' = Ax, classify by eigenvalues of A:

\begin{center}
\begin{tabular}{|l|l|l|}
\hline
\textbf{Eigenvalues} & \textbf{Type} & \textbf{Stability} \\
\hline
$\lambda_1$, $\lambda_2$ real, both < 0 & Stable node & Asymptotically stable \\
\hline
$\lambda_1$, $\lambda_2$ real, both > 0 & Unstable node & Unstable \\
\hline
$\lambda_1$ < 0 < $\lambda_2$ & Saddle point & Unstable \\
\hline
$\alpha\pm\beta$i, $\alpha$ < 0 & Stable spiral (focus) & Asymptotically stable \\
\hline
$\alpha\pm\beta$i, $\alpha$ > 0 & Unstable spiral & Unstable \\
\hline
$\pm\beta$i (pure imaginary) & Center & Stable (not asymptotic) \\
\hline
$\lambda_1$ = $\lambda_2$ < 0, one eigenvector & Stable improper node & Asymptotically stable \\
\hline
\end{tabular}
\end{center}

\begin{verbatim}
Stability Summary (trace-determinant plane):

     det(A)
      ↑
      |  Stable    |  Unstable
      |  Spiral    |  Spiral
      |            |
------+---Saddles---+------→ tr(A)
      |            |
      |  Stable    |  Unstable
      |  Node      |  Node
\end{verbatim}

\subsection{Lyapunov's Direct Method }

Instead of solving the equation, find a \textbf{Lyapunov function} V(x):

\subsubsection{Lyapunov Stability Theorem}
If there exists V: D  $\mathbb{R}$ with:
\begin{itemize}
  \item V(0) = 0 and V(x) > 0 for x $\neq$ 0 (positive definite)
  \item V̇(x) = $\nabla$V $\cdot$ f(x) $\leq$ 0 (negative semi-definite)
\end{itemize}

Then the origin is \textbf{stable}.

If additionally V̇(x) < 0 for x $\neq$ 0 (negative definite), then \textbf{asymptotically stable}.

\subsubsection{How to Find Lyapunov Functions}
Common choices:
\begin{itemize}[leftmargin=*]
  \item V = x$^2$ + y$^2$ (energy-like)
  \item V = xᵀPx where P is positive definite (for linear systems, solve AᵀP + PA = −Q)
\end{itemize}

\subsubsection{Example}
x' = −x$^3$, y' = −y. Try V = x$^2$ + y$^2$:

V̇ = 2x(−x$^3$) + 2y(−y) = −2x$^4$ − 2y$^2$ < 0 for (x,y) $\neq$ (0,0)

So origin is \textbf{asymptotically stable}. 

\bigskip\noindent\rule{\textwidth}{0.4pt}\bigskip

\section{Unit 8: Paths of Linear Systems}

\begin{prereqbox}
 \textbf{Quick Refresher Before Starting:}
\textbf{Phase portrait:} A picture showing ALL trajectories of a 2D linear system X' = AX.
The eigenvalues of A determine the portrait type:
- Both negative  stable node (all trajectories  origin)
- Both positive  unstable node (all trajectories  away)
- Opposite signs  saddle (some approach, some flee)
- Complex with negative real part  stable spiral
- Pure imaginary  center (circles)
\end{prereqbox}


\subsection{Phase Portraits for 2D Linear Systems}

\subsubsection{Stable Node ($\lambda_1$ < $\lambda_2$ < 0)}
All trajectories approach origin along the eigenvector of the smaller (more negative) eigenvalue.

\begin{verbatim}
        ╲   │   ╱
         ╲  │  ╱
          ╲ │ ╱
           ╲│╱
    ────────●────────  (origin)
           ╱│╲
          ╱ │ ╲
         ╱  │  ╲
        ╱   │   ╲
    Arrows point INWARD
\end{verbatim}

\subsubsection{Saddle Point ($\lambda_1$ < 0 < $\lambda_2$)}
\begin{verbatim}
         ↗  │  ↖
          ╲ │ ╱
           ╲│╱
    ←───────●───────→  (origin)
           ╱│╲
          ╱ │ ╲
         ↙  │  ↘
    Stable direction: eigenvector of λ₁
    Unstable direction: eigenvector of λ₂
\end{verbatim}

\subsubsection{Spiral ($\alpha\pm\beta$i)}
Trajectories spiral inward ($\alpha$ < 0) or outward ($\alpha$ > 0) around origin.

\subsubsection{Center ($\pm\beta$i)}
Closed elliptical orbits around origin. No spiral — periodic motion.

\bigskip\noindent\rule{\textwidth}{0.4pt}\bigskip

\section{Unit 9: Dependence on Parameters}

\begin{prereqbox}
 \textbf{Quick Refresher Before Starting:}
\textbf{Parameter dependence:} If the ODE is y' = f(t, y, $\mu$) where $\mu$ is a parameter, how does the solution change when $\mu$ changes slightly?
\textbf{Key question:} Is the solution a continuous/differentiable function of the parameter $\mu$?
\textbf{Answer:} Yes, under mild conditions! The solution varies smoothly with parameters.
\end{prereqbox}


\subsection{Continuous Dependence}

Solutions depend continuously on:
\begin{itemize}
  \item \textbf{Initial conditions:} Small change in y$_0$  small change in solution
  \item \textbf{Parameters:} If f depends on parameter $\mu$, solutions vary continuously with $\mu$
\end{itemize}

\subsubsection{Theorem (Continuous Dependence on Initial Data)}
If f satisfies Lipschitz condition with constant L, and $\varphi$, $\psi$ are solutions with different initial values:

\textbf{|$\varphi$(t) − $\psi$(t)| $\leq$ |$\varphi$(t$_0$) − $\psi$(t$_0$)| $\cdot$ e^(L|t−t$_0$|)}

\subsubsection{Sensitivity to Parameters}
For y' = f(t, y, $\mu$), the \textbf{variational equation} describes how the solution changes with $\mu$:

$\partial$/$\partial\mu$ [y(t,$\mu$)] satisfies a linear ODE obtained by differentiating the original equation w.r.t. $\mu$.

\subsection{Bifurcation Theory (Introduction)}

A \textbf{bifurcation} occurs when the qualitative behavior of solutions changes as a parameter crosses a critical value.

\subsubsection{Saddle-Node Bifurcation}
x' = $\mu$ − x$^2$

\begin{itemize}[leftmargin=*]
  \item $\mu$ < 0: no equilibria
  \item $\mu$ = 0: one equilibrium at x = 0 (semi-stable)
  \item $\mu$ > 0: two equilibria at x = $\pm$√$\mu$ (one stable, one unstable)
\end{itemize}

\bigskip\noindent\rule{\textwidth}{0.4pt}\bigskip

\section{Unit 10: Poincaré-Bendixson Theorem (Detailed)}

\begin{prereqbox}
 \textbf{Quick Refresher Before Starting:}
\textbf{Poincaré-Bendixson (full version):} In a 2D autonomous system, if:
(1) A trajectory stays in a bounded region, AND
(2) The region contains no equilibrium points,
THEN the trajectory approaches a periodic orbit (limit cycle).
\textbf{Why only 2D?} In 3D and higher, trajectories can exhibit chaos (like the Lorenz attractor). The theorem fails!
\end{prereqbox}


\subsection{Full Statement}

Let $\gamma$⁺ be a positive semi-orbit of a C$^1$ planar system that is contained in a compact set K. Then the $\omega$-limit set $\omega$($\gamma$⁺) is one of:

\begin{itemize}
  \item A single equilibrium point
  \item A periodic orbit
  \item A set of equilibria and orbits connecting them (homoclinic/heteroclinic orbits)
\end{itemize}

If K contains no equilibria, then $\omega$($\gamma$⁺) is a periodic orbit.

\subsection{Dulac's Criterion (Ruling Out Periodic Orbits)}

If there exists B(x,y) such that $\partial$(Bf)/$\partial$x + $\partial$(Bg)/$\partial$y has constant sign in D (simply connected), then no periodic orbits exist in D.

\subsubsection{Example}
x' = y, y' = −x − y$^3$. Try B = 1:
$\partial$f/$\partial$x + $\partial$g/$\partial$y = 0 + (−3y$^2$) = −3y$^2\leq$ 0

Since this is $\leq$ 0 (and = 0 only on y = 0, not on a whole orbit), no periodic orbits.

\subsection{Index Theory}

The \textbf{index} of a closed curve C around equilibria:

I(C) = ($\frac{1}{2}$$\pi$) ∮_C d$\theta$ (net rotation of the vector field)

Key facts:
\begin{itemize}[leftmargin=*]
  \item Index of a node, focus, or center = +1
  \item Index of a saddle = −1
  \item A periodic orbit must enclose equilibria whose indices sum to +1
\end{itemize}

\bigskip\noindent\rule{\textwidth}{0.4pt}\bigskip


\textbf{Q1.} Apply Picard iteration to y' = 2ty, y(0) = 1. Find $\varphi_3$(t).

\textbf{Solution:}
$\varphi_0$ = 1
$\varphi_1$ = 1 + $\int_0^t$ 2s$\cdot$1 ds = 1 + t$^2$
$\varphi_2$ = 1 + $\int_0^t$ 2s(1+s$^2$) ds = 1 + t$^2$ + t$^4$/2
$\varphi_3$ = 1 + $\int_0^t$ 2s(1+s$^2$+s$^4$/2) ds = 1 + t$^2$ + t$^4$/2 + t⁶/6
(Converging to e^(t$^2$))

\textbf{Q2.} Classify the critical point of x' = −2x + y, y' = x − 2y.

\textbf{Solution:} A = |−2  1|, eigenvalues: (−2−$\lambda$)$^2$ − 1 = 0  $\lambda$ = −1, −3
                  | 1 −2|
Both negative  \textbf{Stable node}, asymptotically stable.

\textbf{Q3.} Show x' = x$^2$ + y$^2$, y' = −2xy has no periodic orbits in $\mathbb{R}^2$.

\textbf{Solution:} $\partial$f/$\partial$x + $\partial$g/$\partial$y = 2x + (−2x) = 0. Bendixson's criterion fails (zero). Try Dulac with B = 1/y$^2$: $\partial$(f/y$^2$)/$\partial$x + $\partial$(g/y$^2$)/$\partial$y = 2x/y$^2$ + $\partial$(−2x/y)/$\partial$y = 2x/y$^2$ + 2x/y$^2$ = 4x/y$^2$ which is NOT of constant sign. Need different approach for this one.

\textbf{Q4.} Find a Lyapunov function for x' = −x + y$^2$, y' = −y.

\textbf{Solution:} Try V = x$^2$ + y$^2$. V̇ = 2x(−x+y$^2$) + 2y(−y) = −2x$^2$ + 2xy$^2$ − 2y$^2$.
For small (x,y): V̇ $\approx$ −2x$^2$ − 2y$^2$ < 0. So origin is locally asymptotically stable.

\bigskip\noindent\rule{\textwidth}{0.4pt}\bigskip

\begin{quote}
\textbf{ Recommended Reading:} Coddington & Levinson "Theory of ODEs", Perko "Differential Equations and Dynamical Systems"
\textbf{Next Subject:} [Dynamics of a Rigid Body ](./06_Dynamics_of_Rigid_Body.md)
\end{quote}

\newpage

\chapter{Differential Equations --- Proofs \& Problems}


\begin{quote}
\textbf{Supplement to 05_Advanced_Differential_Equations.md}
\end{quote}

\bigskip\noindent\rule{\textwidth}{0.4pt}\bigskip


\subsection{Proof 1: Picard-Lindelöf (Existence & Uniqueness) ⚡⚡⚡}

\begin{tcolorbox}[colback=lightblue,colframe=thmblue,title=\textbf{Theorem},breakable,boxrule=1.5pt]
\textbf{Theorem:} If f(t,y) is continuous and Lipschitz in y (|f(t,y$_1$)−f(t,y$_2$)| $\leq$ L|y$_1$−y$_2$|), then y' = f(t,y), y(t$_0$) = y$_0$ has a unique solution.
\end{tcolorbox}

\begin{proof}
\textbf{Proof (Banach Fixed Point):}
Define operator T: C[t$_0$−$\alpha$, t$_0$+$\alpha$]  C by:
(T$\varphi$)(t) = y$_0$ + $\int$_{t$_0$}^{t} f(s, $\varphi$(s)) ds
\end{proof}

\textbf{T is a contraction:} 
\begin{center}
\begin{tabular}{|l|l|l|l|l|l|l|l|l|}
\hline
\textbf{T$\varphi_1$(t) − T$\varphi_2$(t)} & \textbf{$\leq\int$} & \textbf{f(s,$\varphi_1$)−f(s,$\varphi_2$)} & \textbf{ds $\leq$ L$\int$} & \textbf{$\varphi_1$−$\varphi_2$} & \textbf{ds $\leq$ L$\alpha\cdot$} & \textbf{} & \textbf{$\varphi_1$−$\varphi_2$} & \textbf{} \\
\hline
\end{tabular}
\end{center}

For $\alpha$ < 1/L, T is a contraction with constant L$\alpha$ < 1.

By Banach fixed point theorem, T has a unique fixed point $\varphi$. This $\varphi$ satisfies:
$\varphi$(t) = y$_0$ + $\int$_{t$_0$}^{t} f(s, $\varphi$(s)) ds

Differentiating: $\varphi$'(t) = f(t, $\varphi$(t)) and $\varphi$(t$_0$) = y$_0$. So $\varphi$ is the unique solution. $\blacksquare$

\bigskip\noindent\rule{\textwidth}{0.4pt}\bigskip

\subsection{Proof 2: Gronwall's Inequality ⚡⚡⚡}

\begin{tcolorbox}[colback=lightblue,colframe=thmblue,title=\textbf{Theorem},breakable,boxrule=1.5pt]
\textbf{Theorem:} If u(t) $\leq\alpha$ + $\int$_{t$_0$}^{t} $\beta$(s)u(s)ds with u,$\beta\geq$ 0, $\alpha\geq$ 0 constant, then u(t) $\leq\alpha\cdot$exp($\int$_{t$_0$}^{t} $\beta$(s)ds).
\end{tcolorbox}

\begin{proof}

Let v(t) = $\int$_{t$_0$}^{t} $\beta$(s)u(s)ds. Then v' = $\beta$u $\leq\beta$($\alpha$+v) = $\alpha\beta$ + $\beta$v.
\end{proof}

So v' − $\beta$v $\leq\alpha\beta$. Multiply by e^(−$\int\beta$): d/dt[v$\cdot$e^(−$\int\beta$)] $\leq\alpha\beta\cdot$e^(−$\int\beta$)

Integrate from t$_0$ to t:
v(t)$\cdot$e^(−$\int\beta$) $\leq\alpha$(1 − e^(−$\int\beta$))

So v(t) $\leq\alpha$(e^($\int\beta$) − 1), and u(t) $\leq\alpha$ + v(t) $\leq\alpha\cdot$e^($\int\beta$). $\blacksquare$

\bigskip\noindent\rule{\textwidth}{0.4pt}\bigskip

\subsection{Proof 3: Lyapunov Stability Theorem ⚡⚡⚡}

\begin{tcolorbox}[colback=lightblue,colframe=thmblue,title=\textbf{Theorem},breakable,boxrule=1.5pt]
\textbf{Theorem:} If V: D  $\mathbb{R}$ satisfies V(0)=0, V(x)>0 for x$\neq$0, and V̇(x) $\leq$ 0, then origin is stable. If V̇ < 0 for x$\neq$0, then asymptotically stable.
\end{tcolorbox}

\begin{proof}
\textbf{Proof (Stability):}
Given $\varepsilon$ > 0, let m = min{V(x) : ||x|| = $\varepsilon$} > 0 (V positive on sphere).
Choose $\delta$ > 0 such that ||x$_0$|| < $\delta\Longrightarrow$ V(x$_0$) < m.
\end{proof}

If ||x(t$_0$)|| < $\delta$, then V(x(t$_0$)) < m. Since V̇ $\leq$ 0, V(x(t)) $\leq$ V(x(t$_0$)) < m for all t $\geq$ t$_0$.
But V(x) $\geq$ m on ||x|| = $\varepsilon$, so ||x(t)|| < $\varepsilon$ for all t $\geq$ t$_0$. This is stability.

\textbf{Asymptotic stability:} V̇ < 0 $\Longrightarrow$ V(x(t)) is strictly decreasing. V is bounded below by 0, so V(x(t))  c $\geq$ 0. If c > 0, x(t) stays away from origin, so V̇ $\leq$ −$\delta$ < 0 in that region, forcing V  −$\infty$. Contradiction. So c = 0, meaning x(t)  0. $\blacksquare$

\bigskip\noindent\rule{\textwidth}{0.4pt}\bigskip

\subsection{Proof 4: Bendixson's Criterion ⚡⚡}

\begin{tcolorbox}[colback=lightblue,colframe=thmblue,title=\textbf{Theorem},breakable,boxrule=1.5pt]
\textbf{Theorem:} If $\partial$f/$\partial$x + $\partial$g/$\partial$y is of one sign (positive or negative, never zero) in a simply connected region D, then x' = f, y' = g has no periodic orbits in D.
\end{tcolorbox}

\begin{proof}
\textbf{Proof (by contradiction):}
Suppose $\Gamma$ is a periodic orbit in D enclosing region R. By Green's theorem:
∮_$\Gamma$ (f dy − g dx) = ∬_R ($\partial$f/$\partial$x + $\partial$g/$\partial$y) dA
\end{proof}

LHS: On $\Gamma$, dx = f dt, dy = g dt. So ∮(f$\cdot$g − g$\cdot$f)dt = 0.
RHS: $\neq$ 0 since integrand has constant sign and is not identically zero.
Contradiction! $\blacksquare$

\bigskip\noindent\rule{\textwidth}{0.4pt}\bigskip


\subsection{Picard Iteration — 4 Complete Examples}

\textbf{Ex 1:} y' = 1 + y$^2$, y(0) = 0.
$\varphi_0$ = 0, $\varphi_1$ = $\int_0^t$ (1+0)ds = t, $\varphi_2$ = $\int_0^t$ (1+s$^2$)ds = t + t$^3$/3
$\varphi_3$ = $\int_0^t$ (1+(s+s$^3$/3)$^2$)ds = t + t$^3$/3 + 2t⁵/15 + t⁷/63
Converging to tan(t) (which blows up at t = $\pi$/2).

\textbf{Ex 2:} y' = −2ty, y(0) = 1.
$\varphi_0$ = 1, $\varphi_1$ = 1 + $\int_0^t$ (−2s)ds = 1 − t$^2$
$\varphi_2$ = 1 + $\int_0^t$ −2s(1−s$^2$)ds = 1 − t$^2$ + t$^4$/2
Converging to e^(−t$^2$).

\textbf{Ex 3:} y' = y + t, y(0) = 1.
$\varphi_0$ = 1, $\varphi_1$ = 1 + $\int_0^t$ (1+s)ds = 1 + t + t$^2$/2
$\varphi_2$ = 1 + $\int_0^t$ (1+s+s$^2$/2+s)ds = 1 + t + t$^2$ + t$^3$/6
Converging to 2e$^t$ − t − 1.

\textbf{Ex 4 (Non-unique):} y' = 3y^($\frac{2}{3}$), y(0) = 0.
Both y ≡ 0 and y = t$^3$ are solutions. Lipschitz fails at y = 0 since |3y^($\frac{2}{3}$)| has infinite slope there.

\subsection{Critical Point Classification — 6 Complete Examples}

\textbf{Ex 1:} x' = −x, y' = −2y. A = diag(−1,−2). $\lambda$ = −1,−2. \textbf{Stable node.}

\textbf{Ex 2:} x' = x, y' = −y. A = diag(1,−1). $\lambda$ = 1,−1. \textbf{Saddle point (unstable).}

\textbf{Ex 3:} x' = −y, y' = x. A = [[0,−1],[1,0]]. $\lambda$ = $\pm$i. \textbf{Center (stable, not asymptotic).}

\textbf{Ex 4:} x' = −x+y, y' = −x−y. A = [[−1,1],[−1,−1]]. $\lambda$ = −1$\pm$i. $\alpha$ = −1 < 0. \textbf{Stable spiral.}

\textbf{Ex 5:} x' = 2x−y, y' = x+2y. $\lambda$ = 2$\pm$i (from det: $\lambda^2$−4$\lambda$+5=0). $\alpha$ = 2 > 0. \textbf{Unstable spiral.}

\textbf{Ex 6 (Nonlinear):} x' = −x+y$^2$, y' = −y. Linearize at (0,0): A = [[−1,0],[0,−1]]. $\lambda$ = −1,−1. Linear says stable node. Lyapunov confirms: V = x$^2$+y$^2$, V̇ = 2x(−x+y$^2$)+2y(−y) = −2x$^2$−2y$^2$+2xy$^2$ < 0 near origin. \textbf{Asymptotically stable.} 

\subsection{Lyapunov Function Construction — 5 Examples}

\textbf{Ex 1:} x' = −x$^3$, y' = −y. V = x$^2$+y$^2$. V̇ = −2x$^4$−2y$^2$ < 0. Asymp. stable. 

\textbf{Ex 2:} x' = y, y' = −x−y$^3$. V = x$^2$+y$^2$. V̇ = 2xy+2y(−x−y$^3$) = −2y$^4\leq$ 0. Stable (not asymp. by this V alone, but by LaSalle: V̇=0 only when y=0, then x'=0 forces x=0. So asymp. stable). 

\textbf{Ex 3:} x' = −x+2y$^2$, y' = −y. V = x$^2$+y$^4$. V̇ = 2x(−x+2y$^2$)+4y$^3$(−y) = −2x$^2$+4xy$^2$−4y$^4$ = −2(x−y$^2$)$^2$−2y$^4\leq$ 0. 

\textbf{Ex 4:} For linear system x' = Ax, find V = xᵀPx where AᵀP+PA = −I. Then V̇ = −xᵀx < 0. Solve the \textbf{Lyapunov equation} for P.

\textbf{Ex 5 (Showing instability):} x' = x$^3$, y' = y. Try V = x$^2$+y$^2$. V̇ = 2x$^4$+2y$^2$ > 0 near origin. Origin is \textbf{unstable}. (Chetaev's theorem.)

\bigskip\noindent\rule{\textwidth}{0.4pt}\bigskip


\textbf{P1.} Classify: x' = −3x+4y, y' = −2x+3y. 
\textbf{Solution:} A = [[−3,4],[−2,3]]. tr = 0, det = −9+8 = −1 < 0. \textbf{Saddle point.}

\textbf{P2.} Is the system x' = y, y' = −sin x − y a center at origin?
\textbf{Solution:} Linearize: A = [[0,1],[−1,−1]]. $\lambda$ = (−1$\pm$$\sqrt{1−4}$)/2 = (−1$\pm$i$\sqrt{3}$)/2. $\alpha$ = −$\frac{1}{2}$ < 0. \textbf{Stable spiral} (not center).

\textbf{P3.} Find eᴬ$^t$ for A = [[0,1],[−1,0]].
\textbf{Solution:} $\lambda$ = $\pm$i. eᴬ$^t$ = [[cos t, sin t],[−sin t, cos t]] (rotation matrix).

\textbf{P4.} The system x' = x(1−x−y), y' = y(1−2x−y). Find all equilibria.
\textbf{Solution:} (0,0), (1,0), (0,1), and from 1−x−y=0, 1−2x−y=0: x=0... wait, subtract: x=0, then y=1. So (0,1) already found. Check ($\frac{1}{2}$, $\frac{1}{2}$): 1−$\frac{1}{2}$−$\frac{1}{2}$=0 , 1−1−$\frac{1}{2}$ $\neq$ 0 . Actually: from x$\neq$0,y$\neq$0: 1−x−y=0 and 1−2x−y=0 gives x=0 (subtract). Only non-trivial equilibria: \textbf{(0,0), (1,0), (0,1).}

\textbf{P5.} Show x' = −x(x$^2$+y$^2$), y' = −y(x$^2$+y$^2$) has asymptotically stable origin.
\textbf{Solution:} V = x$^2$+y$^2$. V̇ = −2(x$^2$+y$^2$)$^2$ < 0 for (x,y)$\neq$0. 

\textbf{P6.} Does x' = x$^2$+y$^2$, y' = xy have periodic orbits?
\textbf{Solution:} $\partial$f/$\partial$x + $\partial$g/$\partial$y = 2x+x = 3x. Not of constant sign (changes at x=0). Bendixson inconclusive. Try Dulac with B = 1/y$^2$...

\textbf{P7.} True/False: Every continuous function f: $\mathbb{R}\mathbb{R}$ satisfies Lipschitz condition on compact sets.
\textbf{Answer:} FALSE. f(x) = √|x| is continuous but not Lipschitz at x=0.

\newpage

\chapter{Dynamics of a Rigid Body}


\begin{quote}
\textbf{CSIR NET Priority: ⭐⭐⭐ | Classical Mechanics appears in CSIR NET (Lagrangian/Hamiltonian overlap)}
\end{quote}

\bigskip\noindent\rule{\textwidth}{0.4pt}\bigskip

\subsection{️ Subject Mind Map}

\begin{verbatim}
                DYNAMICS OF A RIGID BODY
                         │
       ┌─────────┬───────┴───────┬──────────┐
       │         │               │          │
   Inertia    Principles     Motion      Impulsive
       │         │               │          │
   ┌───┴───┐  ┌──┴──┐      ┌───┴───┐  ┌───┴───┐
   Moments D'Alembert Fixed  2D      Conservation
   Products Angular   Axis  Motion   Initial
   Theorems Momentum  Rotation       Motion
\end{verbatim}

\bigskip\noindent\rule{\textwidth}{0.4pt}\bigskip

\section{Unit 1: Moments and Products of Inertia}

\begin{prereqbox}
 \textbf{Quick Refresher Before Starting:}
\textbf{Moment of inertia I} = resistance to rotation (like mass is resistance to linear motion).
For a point mass m at distance r from axis: I = mr$^2$.
For a continuous body: I = $\int$r$^2$dm (integrate over the body).
\textbf{Common results you'll derive:}
- Thin rod about center: I = ML$^2$/12
- Thin rod about end: I = ML$^2$/3
- Solid disk about center: I = MR$^2$/2
- Solid sphere about diameter: I = 2MR$^2$/5
\end{prereqbox}



\subsection{What is a Rigid Body?}

A \textbf{rigid body} is a system of particles where the distance between any two particles remains constant during motion. Unlike a single particle, a rigid body can \textbf{rotate}.

\subsection{Moment of Inertia}

The \textbf{moment of inertia} about an axis is the rotational equivalent of mass:

\textbf{I = $\Sigma$ m$_i$r$_i^2$} (discrete) or \textbf{I = $\int$ r$^2$ dm} (continuous)

where r$_i$ is the perpendicular distance from particle i to the axis.

\subsubsection{ Memory Aid}
Mass tells you how hard it is to push something (linear). Moment of inertia tells you how hard it is to SPIN something (rotational).

\subsubsection{Standard Results (MEMORIZE THESE)}

\begin{center}
\begin{tabular}{|l|l|l|}
\hline
\textbf{Body} & \textbf{Axis} & \textbf{Moment of Inertia} \\
\hline
Thin rod (mass M, length L) & Through center, $\perp$ to rod & ML$^2$/12 \\
\hline
Thin rod & Through end, $\perp$ to rod & ML$^2$/3 \\
\hline
Circular ring (M, radius R) & Through center, $\perp$ to plane & MR$^2$ \\
\hline
Circular disk (M, radius R) & Through center, $\perp$ to plane & MR$^2$/2 \\
\hline
Solid sphere (M, radius R) & Through center & 2MR$^2$/5 \\
\hline
Hollow sphere (M, radius R) & Through center & 2MR$^2$/3 \\
\hline
Rectangular plate (M, a$\times$b) & Through center, $\perp$ to plate & M(a$^2$+b$^2$)/12 \\
\hline
\end{tabular}
\end{center}

\subsubsection{Worked Example}
Find the moment of inertia of a uniform rod of mass M, length L about its center.

\begin{verbatim}
Take the rod along x-axis from −L/2 to L/2.
Linear density: λ = M/L

I = ∫_{-L/2}^{L/2} x² · λ dx = (M/L) · [x³/3]_{-L/2}^{L/2}
  = (M/L) · (L³/24 + L³/24) = (M/L) · (L³/12) = ML²/12 ✓
\end{verbatim}

\subsection{Products of Inertia}

\textbf{I_{xy} = −$\Sigma$ m$_i$x$_i$y$_i$} (or −$\int$ xy dm)

Similarly I_{xz}, I_{yz}. These measure the "asymmetry" of mass distribution.

\subsection{Inertia Tensor}

The full picture is captured by the \textbf{inertia tensor} (3$\times$3 symmetric matrix):

\begin{verbatim}
[I] = | I_{xx}   -I_{xy}  -I_{xz} |
      | -I_{xy}   I_{yy}  -I_{yz} |
      | -I_{xz}  -I_{yz}   I_{zz} |
\end{verbatim}

where I_{xx} = $\Sigma$ m$_i$(y$_i^2$+z$_i^2$), etc.

\bigskip\noindent\rule{\textwidth}{0.4pt}\bigskip

\section{Unit 2: Special Cases and Theorems}

\begin{prereqbox}
 \textbf{Quick Refresher Before Starting:}
\textbf{Parallel axis theorem:} I = I_cm + Md$^2$ (shift axis by distance d from center of mass).
Example: Rod about center: ML$^2$/12. Rod about end (d = L/2): ML$^2$/12 + M(L/2)$^2$ = ML$^2$/3.
\textbf{Perpendicular axis theorem} (flat bodies only): I_z = I_x + I_y
Example: Disk: I_x = I_y = MR$^2$/4 (by symmetry), so I_z = MR$^2$/2 
\end{prereqbox}


\subsection{Parallel Axis Theorem (Steiner's Theorem) }

\textbf{I = I_cm + Md$^2$}

where I_cm = moment about axis through center of mass, d = distance between parallel axes.

\subsubsection{Example}
Moment of inertia of a rod about one end:
I_end = I_cm + M(L/2)$^2$ = ML$^2$/12 + ML$^2$/4 = ML$^2$/3 

\subsection{Perpendicular Axis Theorem}

For a \textbf{planar body} (lying in xy-plane):

\textbf{I_z = I_x + I_y}

\subsubsection{Example}
For a circular disk: I_x = I_y (by symmetry) and I_z = MR$^2$/2.
So I_x = I_y = MR$^2$/4.

\subsection{Principal Axes}

Since the inertia tensor is real and symmetric, it can be diagonalized. The eigenvectors are the \textbf{principal axes}, and the eigenvalues are the \textbf{principal moments of inertia} I$_1$, I$_2$, I$_3$.

About principal axes, all products of inertia vanish: I_{xy} = I_{xz} = I_{yz} = 0.

\subsubsection{Example}
For a rectangular plate (a$\times$b, mass M) with sides along x,y axes:
\begin{itemize}[leftmargin=*]
  \item I$_1$ = Ma$^2$/12 (about x-axis)
  \item I$_2$ = Mb$^2$/12 (about y-axis)  
  \item I$_3$ = M(a$^2$+b$^2$)/12 (about z-axis)
  \item All products of inertia = 0 (symmetry)
\end{itemize}

These are already principal axes.

\subsection{Ellipsoid of Inertia}

The set of points satisfying I_{xx}x$^2$ + I_{yy}y$^2$ + I_{zz}z$^2$ − 2I_{xy}xy − 2I_{xz}xz − 2I_{yz}yz = 1 forms an ellipsoid. Its principal axes are the principal axes of inertia.

\bigskip\noindent\rule{\textwidth}{0.4pt}\bigskip

\section{Unit 3: D'Alembert's Principle}

\begin{prereqbox}
 \textbf{Quick Refresher Before Starting:}
\textbf{Newton's second law:} F = ma (force = mass $\times$ acceleration).
\textbf{D'Alembert's principle:} Rewrite as F - ma = 0, treating -ma as a "fictitious force." Now it looks like a statics problem!
\textbf{Virtual work:} If we imagine a small displacement $\delta$r, the work done by all forces (including -ma) is zero: $\Sigma$(F$_i$ - m$_i$a$_i$)$\cdot\delta$r$_i$ = 0.
This powerful idea leads directly to Lagrangian mechanics.
\end{prereqbox}


\subsection{From Newton to D'Alembert}

Newton's law: \textbf{F = ma}  rearrange as \textbf{F − ma = 0}

D'Alembert's idea: treat \textbf{−ma} as a fictitious "inertial force." Then the system is in "equilibrium":

\textbf{F + (−ma) = 0}  Apply methods of statics!

\subsubsection{D'Alembert's Principle (General Form) }

For a system of particles with constraints:

\textbf{$\Sigma$ (F$_i$ − m$_i$a$_i$) $\cdot\delta$r$_i$ = 0}

where $\delta$r$_i$ are virtual displacements consistent with constraints.

\subsubsection{ Why is this useful?}
Constraint forces (like normal forces, tension) do no virtual work for ideal constraints, so they DROP OUT of the equation!

\subsubsection{Example: Particle on Inclined Plane}
Mass m on a frictionless plane inclined at angle $\theta$.

Virtual displacement $\delta$r along the plane:
(mg sin $\theta$ − ma) $\cdot\delta$s = 0  a = g sin $\theta$ 

The normal force N never appeared — it does no work along the allowed motion.

\subsection{Connection to Lagrangian Mechanics }

D'Alembert's principle leads directly to the \textbf{Euler-Lagrange equations}:

\textbf{d/dt($\partial$L/$\partial$q̇$_i$) − $\partial$L/$\partial$q$_i$ = 0}

where L = T − V (Lagrangian = Kinetic − Potential energy).

\subsubsection{Example: Simple Pendulum}
q = $\theta$ (angle from vertical), length l.

T = ½ml$^2\theta$̇$^2$, V = −mgl cos $\theta$, L = ½ml$^2\theta$̇$^2$ + mgl cos $\theta$

$\partial$L/$\partial\theta$̇ = ml$^2\theta$̇, d/dt(ml$^2\theta$̇) = ml$^2\theta$̈
$\partial$L/$\partial\theta$ = −mgl sin $\theta$

Euler-Lagrange: ml$^2\theta$̈ + mgl sin $\theta$ = 0  \textbf{$\theta$̈ + (g/l) sin $\theta$ = 0} 

\bigskip\noindent\rule{\textwidth}{0.4pt}\bigskip

\section{Unit 4: Angular Momentum and Equations of Motion}

\begin{prereqbox}
 \textbf{Quick Refresher Before Starting:}
\textbf{Angular momentum} L = I$\omega$ (analogous to linear momentum p = mv).
\textbf{Torque} $\tau$ = I$\alpha$ = dL/dt (analogous to F = dp/dt).
\textbf{Cross product} (Class 12): a $\times$ b is perpendicular to both a and b. |a $\times$ b| = |a||b|sin$\theta$.
\textbf{Euler's equations:} Govern rotation of a 3D rigid body about its center of mass. They couple the three components of angular velocity.
\end{prereqbox}


\subsection{Angular Momentum }

For a rigid body rotating with angular velocity \textbf{$\omega$}:

\textbf{L = [I] $\cdot\omega$} (angular momentum = inertia tensor $\times$ angular velocity)

For rotation about a principal axis: L = I$\omega$ (scalar form).

\subsection{Euler's Equations of Motion }

In the body-fixed frame (rotating with the body), for principal axes:

\begin{verbatim}
I₁ω̇₁ − (I₂ − I₃)ω₂ω₃ = N₁
I₂ω̇₂ − (I₃ − I₁)ω₃ω₁ = N₂
I₃ω̇₃ − (I₁ − I₂)ω₁ω₂ = N₃
\end{verbatim}

where N$_i$ are the components of external torque.

\subsubsection{ Memory Aid for Signs}
Each equation has the pattern: I_i $\cdot\omega$̇_i − (I_j − I_k) $\cdot\omega$_j $\cdot\omega$_k = N_i with cyclic indices (1231).

\subsubsection{Example: Torque-Free Motion of Symmetric Top}
I$_1$ = I$_2\neq$ I$_3$, N = 0:

$\omega$̇$_3$ = 0  $\omega_3$ = const
$\omega$̇$_1$ = [(I$_2$−I$_3$)/I$_1$]$\omega_2\omega_3$ = $\Omega\omega_2$
$\omega$̇$_2$ = [(I$_3$−I$_1$)/I$_2$]$\omega_3\omega_1$ = −$\Omega\omega_1$

where $\Omega$ = (I$_3$−I$_1$)$\omega_3$/I$_1$. Solution: $\omega_1$ = A cos($\Omega$t), $\omega_2$ = A sin($\Omega$t).

The angular velocity vector \textbf{precesses} around the symmetry axis!

\subsection{Kinetic Energy of Rotation}

\textbf{T = ½$\omega\cdot$ L = ½$\omega$ᵀ[I]$\omega$ = ½(I$_1\omega_1^2$ + I$_2\omega_2^2$ + I$_3\omega_3^2$)}

\bigskip\noindent\rule{\textwidth}{0.4pt}\bigskip

\section{Unit 5: Motion Under Impulsive Forces}

\begin{prereqbox}
 \textbf{Quick Refresher Before Starting:}
\textbf{Impulse} (Class 12): J = F$\Delta$t = change in momentum ($\Delta$p).
For very short time (like a bat hitting a ball), force is huge but time is tiny. The product is finite.
\textbf{For rigid bodies:} Impulsive force  sudden change in velocity.
Impulsive torque  sudden change in angular velocity. $\tau\Delta$t = $\Delta$L = I$\cdot\Delta\omega$.
\end{prereqbox}


\subsection{What is an Impulse?}

An \textbf{impulse} is a very large force acting for a very short time:

\textbf{J = $\int_0$^$\Delta$t F dt = $\Delta$p} (change in momentum)

During an impulse, positions don't change ($\Delta$t $\approx$ 0), but velocities change instantaneously.

\subsection{Impulsive Motion of Rigid Bodies}

\subsubsection{Principle}
\textbf{Change in angular momentum = Impulsive torque}

[I]$\cdot$($\omega$₊ − $\omega$₋) = Impulsive moment

where $\omega$₋ = angular velocity just before, $\omega$₊ = just after.

\subsubsection{Example: Striking a Rod}
A rod of mass M, length L, is at rest. An impulse J is applied perpendicular to the rod at distance d from center.

Change in linear momentum: Mv_cm = J  v_cm = J/M
Change in angular momentum: I_cm $\cdot\omega$ = J$\cdot$d  (ML$^2$/12)$\omega$ = Jd  $\omega$ = 12Jd/(ML$^2$)

\subsubsection{Centre of Percussion}
The point where you can strike without producing a reaction at the pivot. For a rod pivoted at one end, the center of percussion is at distance 2L/3 from the pivot.

\textbf{Condition:} The impulse at the striking point produces zero impulse at the pivot.

\bigskip\noindent\rule{\textwidth}{0.4pt}\bigskip

\section{Unit 6: Motion About a Fixed Axis}

\begin{prereqbox}
 \textbf{Quick Refresher Before Starting:}
\textbf{Fixed axis rotation:} Body can ONLY rotate about one fixed line (like a door on hinges).
\textbf{Key equation:} $\tau$ = I$\alpha$ where I is moment of inertia about the fixed axis, $\alpha$ = angular acceleration.
\textbf{Compound pendulum:} A rigid body swinging about a fixed point (not just a point mass on a string).
Period: T = 2$\pi$$\sqrt{I/(Mgh}$) where h = distance from pivot to center of mass.
\end{prereqbox}


\subsection{The Equation}

For rotation about a fixed axis (say z-axis):

\textbf{I_z $\cdot\theta$̈ = N_z} (moment of inertia $\times$ angular acceleration = torque)

This is the rotational analogue of F = ma.

\subsubsection{Example 1: Compound Pendulum }
A rigid body pivoted at point O, center of mass at distance h from O.

I_O $\cdot\theta$̈ = −Mgh sin $\theta$

For small oscillations (sin $\theta\approx\theta$):

\textbf{Period T = 2$\pi$$\sqrt{I_O/(Mgh}$) = 2$\pi$$\sqrt{k$^2$+h$^2$}$/(gh)}

where k = radius of gyration about center of mass (I_cm = Mk$^2$).

\subsubsection{Equivalent Simple Pendulum Length}
l_eq = I_O/(Mh) = (k$^2$+h$^2$)/h = k$^2$/h + h

\subsubsection{Example 2: Flywheel}
A flywheel (I = 2 kg$\cdot$m$^2$) has a torque of 10 N$\cdot$m applied. Find angular acceleration.

$\theta$̈ = N/I = $\frac{10}{2}$ = 5 rad/s$^2$

After 4 seconds from rest: $\omega$ = $\theta$̈$\cdot$t = 20 rad/s, KE = ½I$\omega^2$ = ½(2)(400) = 400 J.

\subsection{Reactions at the Axis}

For rotation about a fixed axis, the pivot must supply forces to maintain the constraint. These are found from:

Ma_cm = F_external + R (where R = reaction)

\bigskip\noindent\rule{\textwidth}{0.4pt}\bigskip

\section{Unit 7: Rotational Dynamics}

\begin{prereqbox}
 \textbf{Quick Refresher Before Starting:}
\textbf{Kinetic energy of rolling:} T = ½mv$^2$ + ½I$\omega^2$ (translational + rotational).
\textbf{Rolling without slipping:} v = R$\omega$ (the contact point has zero velocity).
\textbf{Energy conservation:} For a body rolling down an incline: mgh = ½mv$^2$ + ½I$\omega^2$.
Substituting v = R$\omega$: v = $\sqrt{2gh/(1 + I/mR$^2$}$). Higher I  slower rolling!
\end{prereqbox}


\subsection{General Motion = Translation + Rotation}

Any rigid body motion can be decomposed as:

\textbf{v_P = v_cm + $\omega\times$ r_{P/cm}}

(velocity of any point = velocity of center of mass + rotation about center of mass)

\subsection{Rolling Motion }

For a body rolling without slipping: \textbf{v_cm = R$\omega$} (no slip condition)

\subsubsection{Example: Disk Rolling Down Incline}
For a uniform disk (I_cm = MR$^2$/2) on a plane inclined at $\theta$:

Energy method: Mgh = ½Mv$^2$ + ½I$\omega^2$ = ½Mv$^2$ + ½(MR$^2$/2)(v/R)$^2$ = ¾Mv$^2$

So v = $\sqrt{4gh/3}$ — slower than a sliding block (v = $\sqrt{2gh}$) because energy goes into rotation!

\subsubsection{Acceleration on incline}
a = g sin $\theta$ / (1 + I_cm/(MR$^2$))

\begin{center}
\begin{tabular}{|l|l|l|}
\hline
\textbf{Body} & \textbf{I_cm/(MR$^2$)} & \textbf{a} \\
\hline
Hollow cylinder & 1 & g sin $\theta$ / 2 (slowest) \\
\hline
Solid cylinder & $\frac{1}{2}$ & 2g sin $\theta$ / 3 \\
\hline
Solid sphere & $\frac{2}{5}$ & 5g sin $\theta$ / 7 \\
\hline
Hollow sphere & $\frac{2}{3}$ & 3g sin $\theta$ / 5 \\
\hline
\end{tabular}
\end{center}

The solid sphere reaches the bottom first!

\subsection{Gyroscopic Precession}

A spinning top (gyroscope) precesses about the vertical at rate:

\textbf{$\Omega$_precession = Mgh/(I_spin $\cdot\omega$_spin)}

Faster spin  slower precession (counterintuitive!).

\bigskip\noindent\rule{\textwidth}{0.4pt}\bigskip

\section{Unit 8: Motion in Two Dimensions}

\begin{prereqbox}
 \textbf{Quick Refresher Before Starting:}
\textbf{2D rigid body motion = translation + rotation.} Three equations govern it:
M$\cdot\ddot{x}$ = $\Sigma$Fₓ (horizontal forces)
M$\cdot\ddot{y}$ = $\Sigma$Fᵧ (vertical forces)
I$\cdot\theta$̈ = $\Sigma$N (torques about center of mass)
Plus constraint equations (e.g., rolling: $\ddot{x}$ = R$\theta$̈).
\end{prereqbox}


\subsection{Equations of Motion (2D) }

For a rigid body moving in a plane:

\begin{verbatim}
M·ẍ_cm = ΣF_x    (translation in x)
M·ÿ_cm = ΣF_y    (translation in y)
I_cm·θ̈ = ΣN_cm   (rotation about center of mass)
\end{verbatim}

Three equations for three unknowns (x_cm, y_cm, $\theta$).

\subsubsection{Example: Cylinder Rolling on Rough Surface}
A solid cylinder (M, R) rolls without slipping on a horizontal surface under force F applied at center.

\begin{verbatim}
Mẍ = F − f        (translation)
Iθ̈ = fR           (rotation about center)
ẍ = Rθ̈             (rolling condition)

From rotation: f = Iθ̈/R = (MR/2)θ̈ = (M/2)ẍ
Substituting: Mẍ = F − (M/2)ẍ → (3M/2)ẍ = F → ẍ = 2F/(3M)
And f = F/3
\end{verbatim}

\subsection{Energy Methods}

\textbf{Total KE = ½Mv$^2$_cm + ½I_cm$\cdot\omega^2$} (König's theorem)

Work-energy theorem: W_net = $\Delta$KE

These often provide the easiest route to solving problems.

\bigskip\noindent\rule{\textwidth}{0.4pt}\bigskip

\section{Unit 9: Impulsive Forces in Two Dimensions}

\begin{prereqbox}
 \textbf{Quick Refresher Before Starting:}
\textbf{Combining impulsive forces with 2D motion:}
Linear impulse: J = M$\cdot\Delta$v (change in linear velocity of CM)
Angular impulse: Moment of J about CM = I$\cdot\Delta\omega$ (change in angular velocity)
\textbf{Centre of percussion:} The point where a blow produces NO reaction at the pivot.
\end{prereqbox}


\subsection{General Equations}

For an impulse applied to a rigid body in 2D:

\begin{verbatim}
M·Δv_x = J_x        (change in x-momentum)
M·Δv_y = J_y        (change in y-momentum)
I_cm·Δω = moment of impulse about CM
\end{verbatim}

\subsubsection{Example: Billiard Ball}
A billiard ball (solid sphere, mass M, radius R) is struck horizontally at height h above the table surface. Find h for pure rolling immediately after the strike.

Impulse J applied horizontally at height h above surface (center at height R):

\begin{verbatim}
M·v = J              (translation)
I_cm·ω = J·(h − R)   (rotation about center)
\end{verbatim}

For pure rolling: v = R$\omega$

So MR$\omega$ = J and (2MR$^2$/5)$\omega$ = J(h−R)

Dividing: (2R/5) = (h−R)  h = R + 2R/5 = \textbf{7R/5}

Strike at height 7R/5 above the table for immediate rolling!

\bigskip\noindent\rule{\textwidth}{0.4pt}\bigskip

\section{Unit 10: Conservation Principles and Initial Motion}

\begin{prereqbox}
 \textbf{Quick Refresher Before Starting:}
\textbf{Conservation laws} (from Class 12 physics):
- \textbf{Energy:} T + V = constant (if no non-conservative forces)
- \textbf{Linear momentum:} If no external force, p = constant
- \textbf{Angular momentum:} If no external torque, L = constant
\textbf{Initial motion problems:} At t = 0, find the initial accelerations and reactions just as motion begins from rest.
\end{prereqbox}


\subsection{Conservation Laws }

\subsubsection{Conservation of Energy}
If all forces are conservative: \textbf{T + V = constant}

\subsubsection{Conservation of Linear Momentum}
If no external forces: \textbf{Mv_cm = constant}

\subsubsection{Conservation of Angular Momentum}
If no external torque about an axis: \textbf{L = I$\omega$ = constant}

\subsubsection{Example: Figure Skater}
Skater pulls arms in: I decreases  $\omega$ increases (since L = I$\omega$ = const).

If I$_1$ = 4 kg$\cdot$m$^2$ at $\omega_1$ = 2 rad/s, and arms pulled to I$_2$ = 1 kg$\cdot$m$^2$:
$\omega_2$ = I$_1\omega_1$/I$_2$ = 4(2)/1 = 8 rad/s (spins 4$\times$ faster!)

\subsection{Initial Motion Problems}

These involve finding accelerations at the instant a system is released from rest.

\textbf{Method:}
\begin{itemize}
  \item Set up equations of motion
  \item Substitute initial conditions (velocities = 0 at start)
  \item Solve the resulting algebraic equations for accelerations
\end{itemize}

\subsubsection{Example}
A uniform rod of length 2a and mass M is held horizontally with one end hinged. It is released. Find initial angular acceleration.

At the instant of release ($\omega$ = 0):
I_O $\cdot\theta$̈ = Mga (torque about hinge)

I_O = M(2a)$^2$/3 = 4Ma$^2$/3

$\theta$̈ = Mga/(4Ma$^2$/3) = \textbf{3g/(4a)}

Acceleration of the free end: a_end = 2a $\cdot\theta$̈ = 3g/2 > g!

The free end accelerates faster than free fall initially! (Surprising but true.)

\bigskip\noindent\rule{\textwidth}{0.4pt}\bigskip


\textbf{Q1.} A solid sphere and a hollow sphere, both of same mass and radius, roll down an incline from the same height. Which reaches the bottom first?

\textbf{Answer:} Solid sphere. a_solid = 5g sin$\theta$/7 > a_hollow = 3g sin$\theta$/5.

\textbf{Q2.} A rod of mass M, length L is free to rotate about one end. A bullet of mass m hits the other end with velocity v and embeds. Find the angular velocity after impact.

\textbf{Solution:} Conservation of angular momentum about pivot:
mvL = (ML$^2$/3 + mL$^2$)$\omega\omega$ = 3mv/((M+3m)L)

\textbf{Q3.} Find the principal moments of inertia of a uniform cube of mass M, side a, about axes through the center parallel to edges.

\textbf{Answer:} By symmetry, I$_1$ = I$_2$ = I$_3$ = Ma$^2$/6.

\textbf{Q4.} A compound pendulum has period 2$\pi$$\sqrt{5/3g}$. If the body is a uniform rod of length L pivoted at one end, find L.

\textbf{Solution:} T = 2$\pi$$\sqrt{I/(Mgh}$) where I = ML$^2$/3, h = L/2.
T = 2$\pi$$\sqrt{2L/(3g}$). Setting equal: 2L/(3g) = 5/(3g)  L = $\frac{5}{2}$ = 2.5 m.

\bigskip\noindent\rule{\textwidth}{0.4pt}\bigskip

\begin{quote}
\textbf{ Recommended Reading:} Goldstein "Classical Mechanics", S.L. Loney "Dynamics of a Particle and Rigid Bodies"
\textbf{ Back to [Master Index](./00_MASTER_INDEX.md)}
\end{quote}

\newpage

\chapter{Dynamics --- Proofs \& Problems}


\begin{quote}
\textbf{Supplement to 06_Dynamics_of_Rigid_Body.md}
\end{quote}

\bigskip\noindent\rule{\textwidth}{0.4pt}\bigskip


\subsection{Proof 1: Parallel Axis Theorem ⚡⚡⚡}

\begin{tcolorbox}[colback=lightblue,colframe=thmblue,title=\textbf{Theorem},breakable,boxrule=1.5pt]
\textbf{Theorem:} I = I_cm + Md$^2$ where d = distance between parallel axes.
\end{tcolorbox}

\begin{proof}

Let axis through CM be the origin. Parallel axis at distance d.
I = $\Sigma$ m$_i$(r$_i$ + d)$^2$ = $\Sigma$ m$_i$r$_i^2$ + 2d$\cdot\Sigma$ m$_i$r$_i$ + d$^2\cdot\Sigma$ m$_i$
\end{proof}

\begin{itemize}[leftmargin=*]
  \item First term = I_cm
  \item Second term: $\Sigma$ m$_i$r$_i$ = M$\cdot$r̄ = 0 (r$_i$ measured from CM, so r̄ = 0)
  \item Third term = Md$^2$
\end{itemize}

Therefore I = I_cm + Md$^2$. $\blacksquare$

\bigskip\noindent\rule{\textwidth}{0.4pt}\bigskip

\subsection{Proof 2: Perpendicular Axis Theorem ⚡⚡}

\begin{tcolorbox}[colback=lightblue,colframe=thmblue,title=\textbf{Theorem},breakable,boxrule=1.5pt]
\textbf{Theorem:} For a planar body in the xy-plane: I_z = I_x + I_y.
\end{tcolorbox}

\begin{proof}

I_x = $\Sigma$ m$_i$y$_i^2$, I_y = $\Sigma$ m$_i$x$_i^2$ (distances from x and y axes)
I_z = $\Sigma$ m$_i$(x$_i^2$ + y$_i^2$) = $\Sigma$ m$_i$x$_i^2$ + $\Sigma$ m$_i$y$_i^2$ = I_y + I_x. $\blacksquare$
\end{proof}

\bigskip\noindent\rule{\textwidth}{0.4pt}\bigskip

\subsection{Proof 3: König's Theorem (KE Decomposition) ⚡⚡⚡}

\begin{tcolorbox}[colback=lightblue,colframe=thmblue,title=\textbf{Theorem},breakable,boxrule=1.5pt]
\textbf{Theorem:} T = ½Mv$^2$_cm + ½$\Sigma$ m$_i$v'$_i^2$ where v'$_i$ = velocity relative to CM.
\end{tcolorbox}

\begin{proof}

v$_i$ = v_cm + v'$_i$
\end{proof}

T = ½$\Sigma$ m$_i$|v$_i$|$^2$ = ½$\Sigma$ m$_i$|v_cm + v'$_i$|$^2$
= ½$\Sigma$ m$_i$v$^2$_cm + $\Sigma$ m$_i$(v_cm $\cdot$ v'$_i$) + ½$\Sigma$ m$_i$v'$_i^2$
= ½Mv$^2$_cm + v_cm $\cdot$ ($\Sigma$ m$_i$v'$_i$) + ½$\Sigma$ m$_i$v'$_i^2$

But $\Sigma$ m$_i$v'$_i$ = d/dt($\Sigma$ m$_i$r'$_i$) = 0 (since CM is origin in body frame).

So T = ½Mv$^2$_cm + T_rot. For rigid body: T_rot = ½I_cm$\cdot\omega^2$. $\blacksquare$

\bigskip\noindent\rule{\textwidth}{0.4pt}\bigskip

\subsection{Proof 4: Euler-Lagrange Equations from D'Alembert ⚡⚡}

\textbf{Derivation:}
D'Alembert: $\Sigma$(F$_i$ − m$_i$r̈$_i$)$\cdot\delta$r$_i$ = 0

Express r$_i$ = r$_i$(q$_1$,...,q$_n$,t). Then $\delta$r$_i$ = $\Sigma_j$ ($\partial$r$_i$/$\partial$q$_j$)$\delta$q$_j$

After substitution and using:
\begin{itemize}[leftmargin=*]
  \item $\partial$ṙ$_i$/$\partial$q̇$_j$ = $\partial$r$_i$/$\partial$q$_j$
  \item d/dt($\partial$r$_i$/$\partial$q$_j$) = $\partial$ṙ$_i$/$\partial$q$_j$
\end{itemize}

We get: $\Sigma_j$ [d/dt($\partial$T/$\partial$q̇$_j$) − $\partial$T/$\partial$q$_j$ − Q$_j$]$\delta$q$_j$ = 0

For independent $\delta$q$_j$: \textbf{d/dt($\partial$T/$\partial$q̇$_j$) − $\partial$T/$\partial$q$_j$ = Q$_j$}

For conservative Q$_j$ = −$\partial$V/$\partial$q$_j$ with L = T − V:

\textbf{d/dt($\partial$L/$\partial$q̇$_j$) − $\partial$L/$\partial$q$_j$ = 0} $\blacksquare$

\bigskip\noindent\rule{\textwidth}{0.4pt}\bigskip


\textbf{P1.} Find MI of a uniform disk (M, R) about a tangent in its plane.

\textbf{Solution:} I_diameter = MR$^2$/4 (by perpendicular axis: I_z = 2I_d  I_d = MR$^2$/4).
By parallel axis: I_tangent = MR$^2$/4 + MR$^2$ = \textbf{5MR$^2$/4}.

\textbf{P2.} A uniform rod (M=2kg, L=1m) rotates about one end. Torque = 6 N$\cdot$m. Find angular acceleration.

\textbf{Solution:} I = ML$^2$/3 = 2(1)/3 = $\frac{2}{3}$ kg$\cdot$m$^2$. $\alpha$ = $\tau$/I = 6/($\frac{2}{3}$) = \textbf{9 rad/s$^2$}.

\textbf{P3.} Solid sphere vs solid cylinder race down incline (same M, R). Who wins?

\textbf{Solution:} a_sphere = 5g sin$\theta$/7, a_cylinder = 2g sin$\theta$/3. Ratio: ($\frac{5}{7}$)/($\frac{2}{3}$) = $\frac{15}{14}$ > 1. \textbf{Sphere wins} (higher acceleration).

\textbf{P4.} A rod (M, L) hangs from one end. It's displaced by angle $\theta_0$ and released. Find period of small oscillations.

\textbf{Solution:} I_end = ML$^2$/3, h = L/2 (CM distance from pivot).
T = 2$\pi$$\sqrt{I/(Mgh}$) = 2$\pi$$\sqrt{ML$^2$/3 / (MgL/2}$) = 2$\pi$$\sqrt{2L/3g}$.
Equivalent simple pendulum: l = 2L/3.

\textbf{P5.} A flywheel (I = 5 kg$\cdot$m$^2$) rotating at 60 rpm is brought to rest by friction in 30 s. Find frictional torque.

\textbf{Solution:} $\omega_0$ = 60$\cdot$2$\pi$/60 = 2$\pi$ rad/s. $\alpha$ = −$\omega_0$/t = −2$\pi$/30 = −$\pi$/15.
$\tau$ = I$\alpha$ = 5$\cdot\pi$/15 = \textbf{$\pi$/3 $\approx$ 1.05 N$\cdot$m}.

\textbf{P6.} Find MI of a hollow sphere (M, R) about diameter using integration.

\textbf{Solution:} Use spherical shell element at angle $\varphi$: dm = M$\cdot$sin $\varphi$ d$\varphi$/2, distance from axis = R sin$\varphi$.
I = $\int_0\pi$ (R sin$\varphi$)$^2\cdot$ M sin$\varphi$ d$\varphi$/2 = MR$^2$/2 $\cdot\int_0\pi$ sin$^3\varphi$ d$\varphi$ = MR$^2$/2 $\cdot$ $\frac{4}{3}$ = \textbf{2MR$^2$/3}. 

\textbf{P7.} Billiard ball struck at height h = 7R/5 above table. Show it rolls without slipping immediately.

\textbf{Solution:} Impulse J at height h, contact at R from center.
Linear: Mv = J, Angular: (2MR$^2$/5)$\omega$ = J(h−R) = J$\cdot$2R/5.
So $\omega$ = J/(MR) = v/R. Rolling condition v = R$\omega$ satisfied. 

\textbf{P8.} A uniform square plate (side a, mass M). MI about diagonal?

\textbf{Solution:} I_x = I_y = Ma$^2$/12 (about sides through center). I_z = Ma$^2$/6 (perp axis).
By symmetry, I_diagonal = I_x = \textbf{Ma$^2$/12} (diagonal is equivalent to side axis for square).

\textbf{P9.} Two particles of mass m at (a,0) and (0,a). Find products of inertia.

\textbf{Solution:} I_xy = −$\Sigma$m$_i$x$_i$y$_i$ = −[m$\cdot$a$\cdot$0 + m$\cdot$0$\cdot$a] = 0.
I_xx = $\Sigma$m$_i$(y$_i^2$+z$_i^2$) = m(0+0) + m(a$^2$+0) = ma$^2$.
I_yy = m(a$^2$) + m(0) = ma$^2$. Inertia tensor: diag(ma$^2$, ma$^2$, 2ma$^2$).

\textbf{P10.} A compound pendulum has minimum period when h = k (radius of gyration about CM). Prove.

\textbf{Solution:} T = 2$\pi$$\sqrt{(k$^2$+h$^2$}$/(gh)). Minimize (k$^2$+h$^2$)/h = k$^2$/h + h.
d/dh(k$^2$/h + h) = −k$^2$/h$^2$ + 1 = 0  h = k.
T_min = 2$\pi$$\sqrt{2k/g}$. $\blacksquare$

\textbf{P11.} Centre of percussion: Rod pivoted at A, struck at distance x from A. Find x so no impulse at pivot.

\textbf{Solution:} Let impulse J at x from pivot. Linear: M$\Delta$v_cm = J − R (R = reaction at pivot).
Angular about A: I_A$\cdot\Delta\omega$ = Jx.
Rolling: $\Delta$v_cm = (L/2)$\cdot\Delta\omega$.

For R = 0: M$\Delta$v_cm = J and I_A$\cdot\Delta\omega$ = Jx, with M(L/2)$\Delta\omega$ = J.
So x = I_A$\cdot\Delta\omega$/J = I_A/(M$\cdot$L/2) = (ML$^2$/3)/(ML/2) = \textbf{2L/3}.

\textbf{P12.} Derive the acceleration of a body rolling down an incline using energy methods.

\textbf{Solution:} Energy: Mgh = ½Mv$^2$ + ½I$\omega^2$ = ½v$^2$(M + I/R$^2$) where h = s sin$\theta$.
Differentiate: Mg sin$\theta\cdot$ v = v$\cdot$(M + I/R$^2$)$\cdot$a.
So \textbf{a = g sin$\theta$/(1 + I/(MR$^2$))}. $\blacksquare$

\bigskip\noindent\rule{\textwidth}{0.4pt}\bigskip


\begin{verbatim}
Parallel Axis:       I = I_cm + Md²
Perpendicular Axis:  I_z = I_x + I_y  (planar body only)
Euler-Lagrange:      d/dt(∂L/∂q̇) − ∂L/∂q = 0,  L = T − V
Euler's Equations:   I₁ω̇₁ = (I₂−I₃)ω₂ω₃ + N₁  (cyclic)
Rolling no-slip:     v = Rω,  a = Rα
KE:                  T = ½Mv²_cm + ½I_cm·ω²
Compound Pendulum:   T = 2π√((k²+h²)/(gh))
Precession:          Ω = Mgh/(I_spin·ω_spin)
Angular impulse:     I·Δω = ∫τ dt = moment of impulse
Centre of perc.:     x = I_pivot/(M·d_cm)  from pivot
Rolling accel.:      a = g sinθ/(1 + I_cm/(MR²))
\end{verbatim}

\newpage

\chapter{CSIR NET Quick Revision}


\begin{quote}
\textbf{Use this sheet for last-minute revision before CSIR NET or semester exams}
\textbf{Format: Theorem/Concept  One-line statement  Key formula}
\end{quote}

\bigskip\noindent\rule{\textwidth}{0.4pt}\bigskip

\subsection{ ADVANCED LINEAR ALGEBRA — Flash Cards}

\subsubsection{Must-Know Theorems}
\begin{center}
\begin{tabular}{|l|l|l|}
\hline
\textbf{#} & \textbf{Theorem} & \textbf{Statement} \\
\hline
1 & Rank-Nullity & dim(V) = rank(T) + nullity(T) \\
\hline
2 & Cayley-Hamilton & Every matrix satisfies its characteristic polynomial: p(A) = 0 \\
\hline
3 & Spectral Theorem & Real symmetric matrices are orthogonally diagonalizable \\
\hline
4 & Sylvester's Law & Signature (p, q) of a quadratic form is invariant under change of basis \\
\hline
5 & Jordan Form & Every matrix over $\mathbb{C}$ is similar to a block-diagonal of Jordan blocks \\
\hline
6 & Riesz Representation & Every linear functional on inner product space is ⟨$\cdot$, v$_0$⟩ for unique v$_0$ \\
\hline
\end{tabular}
\end{center}

\subsubsection{Key Formulas}
\begin{verbatim}
Eigenvalues:        det(A − λI) = 0
Diagonalizable iff: minimal polynomial has no repeated roots
Orthogonal matrix:  QᵀQ = I, det(Q) = ±1
Cauchy-Schwarz:     |⟨u,v⟩| ≤ ‖u‖·‖v‖
Gram-Schmidt:       uₖ = vₖ − Σ ⟨vₖ,eᵢ⟩eᵢ
dim(W) + dim(W⊥) = dim(V)
dim(W) + dim(W⁰) = dim(V)  (annihilator)
\end{verbatim}

\subsubsection{ Common Traps}
\begin{itemize}[leftmargin=*]
  \item Eigenvectors of DIFFERENT eigenvalues are linearly independent
  \item Symmetric $\neq$ orthogonal (symmetric: Aᵀ=A; orthogonal: AᵀA=I)
  \item Nilpotent matrix: all eigenvalues = 0, but NOT the zero matrix
  \item Similar matrices have same eigenvalues, BUT same eigenvalues $\neq$ similar
\end{itemize}

\bigskip\noindent\rule{\textwidth}{0.4pt}\bigskip

\subsection{ MATHEMATICAL ANALYSIS — Flash Cards}

\subsubsection{Must-Know Theorems}
\begin{center}
\begin{tabular}{|l|l|l|}
\hline
\textbf{#} & \textbf{Theorem} & \textbf{Statement} \\
\hline
1 & Weierstrass M-test &  & f$_n$(x) & $\leq$ M$_n$, $\Sigma$M$_n$ < $\infty\Longrightarrow\Sigma$f$_n$ converges uniformly \\
\hline
2 & Stone-Weierstrass & Continuous functions on [a,b] can be uniformly approximated by polynomials \\
\hline
3 & Arzelà-Ascoli & Uniformly bounded + equicontinuous $\Longrightarrow$ has uniformly convergent subsequence \\
\hline
4 & Implicit Function Thm & F(a,b)=0 and $\partial$F/$\partial$y invertible $\Longrightarrow$ y=g(x) locally \\
\hline
5 & Inverse Function Thm & Df(a) invertible $\Longrightarrow$ f is local diffeomorphism \\
\hline
\end{tabular}
\end{center}

\subsubsection{Key Formulas}
\begin{verbatim}
R-S integral:       ∫f dα = ∫f(x)α'(x)dx  (when α is differentiable)
Integration by parts: ∫f dα + ∫α df = f(b)α(b) − f(a)α(a)
Uniform convergence: sup|fₙ(x) − f(x)| → 0
Jacobian:           J = det[∂fᵢ/∂xⱼ], area element: |J|du dv
Lagrange multiplier: ∇f = λ∇g at constrained extremum
Second deriv test:  D = fₓₓfᵧᵧ − fₓᵧ² (D>0,fₓₓ>0: min; D>0,fₓₓ<0: max; D<0: saddle)
\end{verbatim}

\subsubsection{ Common Traps}
\begin{itemize}[leftmargin=*]
  \item Pointwise convergence does NOT preserve continuity (x$^n$ on [0,1])
  \item Partial derivatives exist does NOT imply differentiable
  \item Continuous partials DOES imply differentiable
  \item Clairaut: fₓᵧ = fᵧₓ only when both are continuous
\end{itemize}

\bigskip\noindent\rule{\textwidth}{0.4pt}\bigskip

\subsection{ TOPOLOGY — Flash Cards}

\subsubsection{Must-Know Theorems}
\begin{center}
\begin{tabular}{|l|l|l|}
\hline
\textbf{#} & \textbf{Theorem} & \textbf{Statement} \\
\hline
1 & Heine-Borel & Subset of $\mathbb{R}^n$ is compact $\Longleftrightarrow$ closed and bounded \\
\hline
2 & Tychonoff & Product of compact spaces is compact \\
\hline
3 & Bolzano-Weierstrass & Infinite subset of compact space has limit point \\
\hline
4 & IVT (topological) & Continuous image of connected space is connected \\
\hline
5 & Urysohn's Lemma & Normal space: disjoint closed sets separated by continuous function \\
\hline
6 & Extreme Value Thm & Continuous real function on compact space attains max and min \\
\hline
\end{tabular}
\end{center}

\subsubsection{Key Facts}
\begin{verbatim}
Topology axioms:     ∅,X ∈ τ; finite ∩; arbitrary ∪
Closed = complement of open
Ā = A ∪ A' (closure = set ∪ limit points)
∂A = Ā \ Int(A)
Connected: no partition into 2 nonempty disjoint open sets
Path connected ⟹ Connected (converse FALSE)
Compact + Hausdorff ⟹ Normal
In metric spaces: Separable ⟺ Second Countable ⟺ Lindelöf
\end{verbatim}

\subsubsection{Separation Axioms Quick Reference}
\begin{verbatim}
T₀: distinguishable    T₁: singletons closed    T₂: Hausdorff (disjoint nbhds)
T₃: Regular (pt & closed set separated)    T₄: Normal (closed sets separated)
\end{verbatim}

\subsubsection{ Common Traps}
\begin{itemize}[leftmargin=*]
  \item Open and closed are NOT opposites (a set can be both, or neither)
  \item Compact does NOT mean closed (only in Hausdorff spaces)
  \item $\mathbb{Q}$ is NOT connected, NOT compact, NOT locally compact
  \item Cofinite topology: T$_1$ but NOT T$_2$; every infinite set is dense; space is compact
\end{itemize}

\bigskip\noindent\rule{\textwidth}{0.4pt}\bigskip

\subsection{ ADVANCED COMPLEX ANALYSIS — Flash Cards}

\subsubsection{Must-Know Theorems}
\begin{center}
\begin{tabular}{|l|l|l|}
\hline
\textbf{#} & \textbf{Theorem} & \textbf{Statement} \\
\hline
1 & Liouville & Entire + bounded $\Longrightarrow$ constant \\
\hline
2 & Little Picard & Non-constant entire function omits at most 1 value \\
\hline
3 & Great Picard & Near essential singularity: takes every value (except maybe 1) infinitely often \\
\hline
4 & Hadamard Factorization & Entire function of finite order = polynomial $\times$ canonical product $\times$ exponential \\
\hline
5 & Maximum Principle &  & f & has no interior maximum for non-constant analytic f \\
\hline
6 & Open Mapping & Non-constant analytic function is an open map \\
\hline
\end{tabular}
\end{center}

\subsubsection{Key Formulas}
\begin{verbatim}
Order of entire function:     ρ = lim sup [log log M(r)]/[log r]
Jensen's formula:             log|f(0)| + Σlog(R/|aₖ|) = (1/2π)∫log|f(Re^iθ)|dθ
Poisson integral:             u(re^iθ) = (1/2π)∫ P(r,θ−t)f(e^it)dt
Weierstrass product:          f(z) = e^g(z) · Π Eₚ(z/aₙ)
sin(πz) = πz · Π(1−z²/n²)
Hadamard 3 circles:           log M(r) is convex in log r
\end{verbatim}

\bigskip\noindent\rule{\textwidth}{0.4pt}\bigskip

\subsection{⚙️ ADVANCED DIFFERENTIAL EQUATIONS — Flash Cards}

\subsubsection{Must-Know Theorems}
\begin{center}
\begin{tabular}{|l|l|l|}
\hline
\textbf{#} & \textbf{Theorem} & \textbf{Statement} \\
\hline
1 & Picard-Lindelöf & f continuous + Lipschitz $\Longrightarrow$ unique solution exists \\
\hline
2 & Peano & f continuous $\Longrightarrow$ solution exists (may not be unique) \\
\hline
3 & Gronwall & u $\leq\alpha$ + $\int\beta$u $\Longrightarrow$ u $\leq\alpha\cdot$exp($\int\beta$) \\
\hline
4 & Poincaré-Bendixson & Bounded orbit in 2D with no equilibria $\Longrightarrow$ limit cycle \\
\hline
5 & Bendixson Criterion & div(f,g) of constant sign $\Longrightarrow$ no periodic orbits \\
\hline
6 & Lyapunov Stability & V>0, V̇$\leq$0 $\Longrightarrow$ stable; V̇<0 $\Longrightarrow$ asymptotically stable \\
\hline
\end{tabular}
\end{center}

\subsubsection{Critical Point Classification (2D)}
\begin{verbatim}
λ₁,λ₂ < 0:           Stable node
λ₁,λ₂ > 0:           Unstable node
λ₁ < 0 < λ₂:         Saddle (unstable)
α±βi, α<0:           Stable spiral
α±βi, α>0:           Unstable spiral
±βi:                  Center (stable, not asymptotic)
\end{verbatim}

\subsubsection{Key Formulas}
\begin{verbatim}
Matrix exponential:    eᴬᵗ = Σ (At)ⁿ/n!
If A = PDP⁻¹:         eᴬᵗ = P·diag(eλᵢt)·P⁻¹
Variation of params:   x(t) = eᴬᵗx₀ + ∫₀ᵗ eᴬ⁽ᵗ⁻ˢ⁾g(s)ds
Wronskian:             W(t) = W(t₀)·exp(∫tr(A)ds)
\end{verbatim}

\bigskip\noindent\rule{\textwidth}{0.4pt}\bigskip

\subsection{ DYNAMICS OF A RIGID BODY — Flash Cards}

\subsubsection{Must-Know Formulas}
\begin{verbatim}
Moment of inertia:     I = ∫r²dm
Parallel axis:         I = I_cm + Md²
Perpendicular axis:    Iᵤ = Iₓ + Iᵧ (planar bodies only)
Euler's equations:     I₁ω̇₁ − (I₂−I₃)ω₂ω₃ = N₁ (cyclic)
Rolling (no slip):     v = Rω,  a = Rα
KE of rigid body:      T = ½Mv²_cm + ½I_cm·ω²
Compound pendulum:     T = 2π√(I_O/(Mgh))
Angular momentum:      L = Iω
\end{verbatim}

\subsubsection{Standard Moments of Inertia}
\begin{verbatim}
Rod (center):    ML²/12       Rod (end):      ML²/3
Disk (center):   MR²/2        Ring:           MR²
Solid sphere:    2MR²/5       Hollow sphere:  2MR²/3
\end{verbatim}

\bigskip\noindent\rule{\textwidth}{0.4pt}\bigskip

\subsection{ CSIR NET EXAM DAY STRATEGY}

\subsubsection{Time Allocation (3 hours = 180 min)}
\begin{verbatim}
Part A (General Aptitude):  30 min — attempt 15/20
Part B (Subject MCQs):      80 min — attempt 20/25
Part C (Advanced MSQs):     60 min — attempt 8-10/20
Buffer/Review:              10 min
\end{verbatim}

\subsubsection{Subject-wise Question Distribution (Typical)}
\begin{verbatim}
Linear Algebra:        4-5 questions (Part B+C)  ← YOUR STRONGEST BET
Real Analysis:         4-5 questions (Part B+C)  ← HIGH ROI
Complex Analysis:      2-3 questions (Part B)
Topology:              2-3 questions (Part B+C)
ODE/PDE:               2-3 questions (Part B)
Algebra (Groups/Rings): 3-4 questions           ← Semester 2/3 topic
Classical Mechanics:    1-2 questions (Part B)
\end{verbatim}

\subsubsection{Golden Rules}
\begin{itemize}
  \item \textbf{Attempt Part A fully} — easy marks, low negative marking
  \item \textbf{Never guess in Part C} — high negative marking (-1.0)
  \item \textbf{Start with your strongest subject} in Part B
  \item \textbf{Skip and return} — don't get stuck on one problem
  \item \textbf{Use elimination} — plug in options for MCQs
\end{itemize}

\newpage

\chapter{Previous Year Questions --- Solved}


\begin{quote}
\textbf{Curated PYQs mapped to your Semester 1 syllabus}
\textbf{Sources: CSIR NET (2019-2025), SET/SLET exams, University recruitment tests}
\end{quote}

\bigskip\noindent\rule{\textwidth}{0.4pt}\bigskip

\subsection{ LINEAR ALGEBRA PYQs}

\subsubsection{CSIR NET June 2024}
\textbf{Q1.} Let A be a 4$\times$4 real matrix with characteristic polynomial (x−1)$^2$(x−2)$^2$ and minimal polynomial (x−1)(x−2). Then A is:
(a) Diagonalizable  (b) Not diagonalizable  (c) Nilpotent  (d) Idempotent

\textbf{Answer: (a)} Minimal polynomial has no repeated roots  diagonalizable.

\bigskip\noindent\rule{\textwidth}{0.4pt}\bigskip

\textbf{Q2.} Let V be the vector space of 2$\times$2 real matrices. Define T(A) = AᵀA. Is T linear?

\textbf{Answer: NO.} T(A+B) = (A+B)ᵀ(A+B) = AᵀA + AᵀB + BᵀA + BᵀB $\neq$ T(A) + T(B).

\bigskip\noindent\rule{\textwidth}{0.4pt}\bigskip

\subsubsection{CSIR NET Dec 2023}
\textbf{Q3.} Let A be an n$\times$n real matrix with A$^n$ = 0 but A$^n^{-}^1\neq$ 0. The Jordan form of A is:
(a) Single Jordan block J$_n$(0)  (b) Diagonal  (c) nJ$_1$(0)  (d) Cannot determine

\textbf{Answer: (a)} A is nilpotent with nilpotency index n  single Jordan block.

\bigskip\noindent\rule{\textwidth}{0.4pt}\bigskip

\textbf{Q4.} If A is a 3$\times$3 orthogonal matrix with det(A) = 1, which eigenvalue is guaranteed?
(a) 0  (b) 1  (c) −1  (d) i

\textbf{Answer: (b)} det(A) = product of eigenvalues = 1. Eigenvalues of orthogonal matrices have |$\lambda$|=1. Complex eigenvalues come in conjugate pairs. For odd-sized matrix, at least one real eigenvalue. Since det=1 and |$\lambda$|=1, one eigenvalue must be 1.

\bigskip\noindent\rule{\textwidth}{0.4pt}\bigskip

\subsubsection{CSIR NET June 2023}
\textbf{Q5.} Dimension of the space of 3$\times$3 real symmetric matrices with trace 0?

\textbf{Answer: 5.} Symmetric 3$\times$3 has dim = 6. Trace = 0 is one linear constraint. So dim = 5.

\bigskip\noindent\rule{\textwidth}{0.4pt}\bigskip

\textbf{Q6.} T: $\mathbb{R}^3\mathbb{R}^3$ defined by T(x,y,z) = (x+2y, y−z, x+2y−z). Find nullity(T).

\textbf{Answer:} Matrix: [[1,2,0],[0,1,−1],[1,2,−1]]. Row reduce  rank = 2. Nullity = 3−2 = \textbf{1}.

\bigskip\noindent\rule{\textwidth}{0.4pt}\bigskip

\subsubsection{SET/SLET Exam Questions}
\textbf{Q7.} The number of linearly independent eigenvectors of [[2,1],[0,2]] is:
(a) 0  (b) 1  (c) 2  (d) 3

\textbf{Answer: (b)} Eigenvalue $\lambda$=2 (repeated). Eigenspace: (A−2I)v=0  [[0,1],[0,0]]v=0  span{(1,0)}. Dimension = 1.

\bigskip\noindent\rule{\textwidth}{0.4pt}\bigskip

\textbf{Q8.} A quadratic form Q(x,y) = 5x$^2$ + 2y$^2$ + 6xy. Classify it.

\textbf{Answer:} Matrix: [[5,3],[3,2]]. Eigenvalues: (5−$\lambda$)(2−$\lambda$)−9 = $\lambda^2$−7$\lambda$+1 = 0. $\lambda$ = (7$\pm$$\sqrt{45}$)/2 $\approx$ 6.85, 0.15. Both positive  \textbf{Positive definite.}

\bigskip\noindent\rule{\textwidth}{0.4pt}\bigskip

\subsection{ REAL ANALYSIS PYQs}

\subsubsection{CSIR NET Dec 2024}
\textbf{Q9.} The sequence f$_n$(x) = nx(1−x$^2$)$^n$ on [0,1]. Find the pointwise limit and determine if convergence is uniform.

\textbf{Answer:} For x=0: f$_n$(0) = 0. For x $\in$ (0,1]: f$_n$(x) = nx(1−x$^2$)$^n$  0 (exponential beats linear). Pointwise limit f ≡ 0. But $\int_0^1$ f$_n$ = n$\cdot$[−(1−x$^2$)$^n$⁺$^1$/(2(n+1))]$_0^1$ = n/(2(n+1))  $\frac{1}{2}$ $\neq$ 0. So convergence is \textbf{NOT uniform}.

\bigskip\noindent\rule{\textwidth}{0.4pt}\bigskip

\subsubsection{CSIR NET June 2024}
\textbf{Q10.} f(x,y) = xy/$\sqrt{x$^2$+y$^2$}$ for (x,y)$\neq$(0,0), f(0,0)=0. Is f differentiable at origin?

\textbf{Answer:} fₓ(0,0) = lim_{h0} f(h,0)/h = 0. Similarly fᵧ(0,0) = 0. Check: |f(h,k)−0−0|/$\sqrt{h$^2$+k$^2$}$ = |hk|/(h$^2$+k$^2$). Along h=k: this = $\frac{1}{2}$ $\neq$ 0. \textbf{Not differentiable.}

\bigskip\noindent\rule{\textwidth}{0.4pt}\bigskip

\subsubsection{CSIR NET Dec 2023}
\textbf{Q11.} Which of the following series converges uniformly on $\mathbb{R}$?
(a) $\Sigma$ x$^n$/n!  (b) $\Sigma$ sin(nx)/n$^2$  (c) $\Sigma$ x$^n$  (d) $\Sigma$ 1/(n+x$^2$)

\textbf{Answer: (b)} By M-test: |sin(nx)/n$^2$| $\leq$ 1/n$^2$ and $\Sigma$1/n$^2$ < $\infty$. Series (a) converges uniformly on bounded sets but not all of $\mathbb{R}$. Series (c) diverges for |x|$\geq$1. Series (d): terms don't go to 0 uniformly.

\bigskip\noindent\rule{\textwidth}{0.4pt}\bigskip

\textbf{Q12.} The Jacobian of the transformation x = r sin$\varphi$ cos$\theta$, y = r sin$\varphi$ sin$\theta$, z = r cos$\varphi$ is:

\textbf{Answer: r$^2$ sin$\varphi$} (spherical coordinates). Volume element: dV = r$^2$ sin$\varphi$ dr d$\varphi$ d$\theta$.

\bigskip\noindent\rule{\textwidth}{0.4pt}\bigskip

\subsubsection{SET Exam Questions}
\textbf{Q13.} True/False: A uniformly convergent sequence of differentiable functions converges to a differentiable function.

\textbf{Answer: FALSE.} Uniform convergence preserves continuity but NOT differentiability. Need uniform convergence of derivatives too.

\bigskip\noindent\rule{\textwidth}{0.4pt}\bigskip

\subsection{ TOPOLOGY PYQs}

\subsubsection{CSIR NET June 2024}
\textbf{Q14.} In the standard topology on $\mathbb{R}$, which of the following is/are connected?
(a) $\mathbb{Q}$  (b) $\mathbb{R}$\$\mathbb{Q}$  (c) [0,1]$\cup$[2,3]  (d) (0,1)

\textbf{Answer: (b) and (d).} $\mathbb{Q}$ is totally disconnected. [0,1]$\cup$[2,3] is disconnected (split at any point in (1,2)). $\mathbb{R}$\$\mathbb{Q}$ (irrationals) IS connected. (0,1) is an interval, hence connected.

\bigskip\noindent\rule{\textwidth}{0.4pt}\bigskip

\subsubsection{CSIR NET Dec 2023}
\textbf{Q15.} Which of the following subsets of $\mathbb{R}^2$ is compact?
(a) {(x,y): x$^2$+y$^2$ < 1}  (b) {(x,y): x$^2$+y$^2\leq$ 1}  (c) {(x,y): x$^2$+y$^2$ = 1}  (d) Both (b) and (c)

\textbf{Answer: (d)} Heine-Borel: compact $\Longleftrightarrow$ closed and bounded. (a) is open (not closed). Both (b) and (c) are closed and bounded.

\bigskip\noindent\rule{\textwidth}{0.4pt}\bigskip

\textbf{Q16.} The cofinite topology on an uncountable set X is:
(a) Hausdorff  (b) T$_1$ but not T$_2$  (c) Connected  (d) Both (b) and (c)

\textbf{Answer: (d)} Singletons are closed (T$_1$). Any two non-empty open sets intersect (not T$_2$). Cannot split X into two non-empty cofinite open sets (connected).

\bigskip\noindent\rule{\textwidth}{0.4pt}\bigskip

\subsubsection{CSIR NET June 2023}
\textbf{Q17.} A topological space X is such that every subset is compact. Then X must be:
(a) Finite  (b) Countable  (c) Discrete  (d) Indiscrete

\textbf{Answer:} If X has the discrete topology and every subset is compact, then X must be \textbf{(a) Finite}.

\bigskip\noindent\rule{\textwidth}{0.4pt}\bigskip

\subsubsection{SET Exam Questions}
\textbf{Q18.} Give an example of a space that is compact but not sequentially compact.

\textbf{Answer:} The product space {0,1}^[0,1] with product topology is compact (Tychonoff) but not sequentially compact (not first countable).

\bigskip\noindent\rule{\textwidth}{0.4pt}\bigskip

\subsection{ COMPLEX ANALYSIS PYQs}

\subsubsection{CSIR NET Dec 2024}
\textbf{Q19.} If f is entire with f(1/n) = 1/n$^2$ for all n $\in\mathbb{N}$, then f(z) = ?

\textbf{Answer: f(z) = z$^2$.} By identity theorem: f and g(z)=z$^2$ agree on {1/n} which has limit point 0. So f = g on all of $\mathbb{C}$.

\bigskip\noindent\rule{\textwidth}{0.4pt}\bigskip

\subsubsection{CSIR NET June 2024}
\textbf{Q20.} The number of zeros of eᶻ − z in the disk |z| < 1 is:

\textbf{Answer:} By Rouché's theorem: |eᶻ| $\leq$ e < 3 and |−z| $\leq$ 1 on |z|=1. Since eᶻ dominates... actually we compare: let f(z) = eᶻ, g(z) = −z. |f(z)| = eᴿᵉ⁽ᶻ⁾ $\geq$ e$^{-}^1$ > 0 and we need more careful analysis. On |z|=1: |eᶻ| $\geq$ e$^{-}^1\approx$ 0.37 but |z| = 1. So |g| > |f| on |z|=1, and −z has 1 zero inside. By Rouché: eᶻ − z has \textbf{1 zero} in |z| < 1.

\bigskip\noindent\rule{\textwidth}{0.4pt}\bigskip

\subsubsection{CSIR NET Dec 2023}
\textbf{Q21.} The order of the entire function f(z) = cos($\sqrt{z}$) is:

\textbf{Answer:} cos($\sqrt{z}$) = $\Sigma$ (−1)$^n$ z$^n$/(2n)!. Growth: |cos($\sqrt{z}$)| $\leq$ e^|$\sqrt{z}$| = e^√|z|. So log M(r) $\approx$ $\sqrt{r}$, log log M(r) $\approx$ ½ log r. Order $\rho$ = ½.

\bigskip\noindent\rule{\textwidth}{0.4pt}\bigskip

\textbf{Q22.} True/False: There exists a non-constant analytic function from $\mathbb{C}$ to the open unit disk.

\textbf{Answer: FALSE.} Such a function would be entire and bounded by 1. By Liouville's theorem, it must be constant.

\bigskip\noindent\rule{\textwidth}{0.4pt}\bigskip

\subsection{⚙️ DIFFERENTIAL EQUATIONS PYQs}

\subsubsection{CSIR NET June 2024}
\textbf{Q23.} The system x' = x − y, y' = x + y has the origin as:
(a) Stable node  (b) Unstable spiral  (c) Saddle  (d) Center

\textbf{Answer: (b)} A = [[1,−1],[1,1]], eigenvalues: (1−$\lambda$)$^2$+1=0  $\lambda$ = 1$\pm$i. Real part $\alpha$ = 1 > 0  \textbf{Unstable spiral.}

\bigskip\noindent\rule{\textwidth}{0.4pt}\bigskip

\subsubsection{CSIR NET Dec 2023}
\textbf{Q24.} For which value(s) of $\alpha$ does y' = |y|^$\alpha$, y(0) = 0 have a unique solution?
(a) $\alpha$ = $\frac{1}{2}$  (b) $\alpha$ = 1  (c) $\alpha$ = 2  (d) $\alpha\geq$ 1

\textbf{Answer: (d)} Uniqueness requires Lipschitz condition. |y|^$\alpha$ is Lipschitz near y=0 iff $\alpha\geq$ 1. For $\alpha$ = $\frac{1}{2}$, y = (t/2)$^2$ is also a solution (non-unique).

\bigskip\noindent\rule{\textwidth}{0.4pt}\bigskip

\textbf{Q25.} The critical point (0,0) of the system x' = −y − x$^3$, y' = x − y$^3$ is:
(a) Stable  (b) Unstable  (c) Asymptotically stable  (d) Center

\textbf{Answer: (c)} Try V = x$^2$ + y$^2$. V̇ = 2x(−y−x$^3$) + 2y(x−y$^3$) = −2x$^4$ − 2y$^4$ < 0 for (x,y) $\neq$ 0. \textbf{Asymptotically stable} by Lyapunov.

\bigskip\noindent\rule{\textwidth}{0.4pt}\bigskip

\subsection{ DYNAMICS PYQs}

\subsubsection{University Recruitment Exam}
\textbf{Q26.} A uniform solid cylinder of mass M and radius R rolls down a rough incline of angle 30°. The acceleration is:

\textbf{Answer:} a = g sin30°/(1 + I/(MR$^2$)) = (g/2)/(1 + $\frac{1}{2}$) = \textbf{g/3 $\approx$ 3.27 m/s$^2$}

\bigskip\noindent\rule{\textwidth}{0.4pt}\bigskip

\textbf{Q27.} Two equal masses are attached to the ends of a light rod of length 2a. The moment of inertia about an axis through the center perpendicular to the rod is:

\textbf{Answer:} I = ma$^2$ + ma$^2$ = \textbf{2ma$^2$}

\bigskip\noindent\rule{\textwidth}{0.4pt}\bigskip

\textbf{Q28.} A compound pendulum oscillates with period T. The distance between pivot and center of mass is h, and I_O = Mh$^2$(1 + k$^2$/h$^2$). Express T.

\textbf{Answer:} T = 2$\pi$$\sqrt{I_O/(Mgh}$) = 2$\pi$$\sqrt{(k$^2$ + h$^2$}$/(gh))

The equivalent simple pendulum length = (k$^2$ + h$^2$)/h.

\bigskip\noindent\rule{\textwidth}{0.4pt}\bigskip

\subsection{ Topic-wise Frequency Analysis (CSIR NET 2019-2025)}

\begin{center}
\begin{tabular}{|l|l|l|l|}
\hline
\textbf{Topic} & \textbf{Avg Questions/Exam} & \textbf{Difficulty} & \textbf{Your Priority} \\
\hline
Linear Algebra & 4-5 & Medium &  Maximum \\
\hline
Real Analysis & 4-5 & Medium-Hard &  Maximum \\
\hline
Abstract Algebra & 3-4 & Hard &  Sem 2-3 topic \\
\hline
Complex Analysis & 2-3 & Medium &  High \\
\hline
Topology & 2-3 & Hard &  High \\
\hline
ODE/PDE & 2-3 & Medium &  High \\
\hline
Probability/Stats & 1-2 & Easy-Medium &  Sem 3-4 topic \\
\hline
Numerical Analysis & 1-2 & Easy &  Quick prep \\
\hline
Classical Mechanics & 1-2 & Medium &  Via Dynamics \\
\hline
\end{tabular}
\end{center}

\bigskip\noindent\rule{\textwidth}{0.4pt}\bigskip

\begin{quote}
\textbf{ Study Tip:} Solve at least 10 years of CSIR NET papers. Focus on Linear Algebra and Real Analysis first — they give you the best score-per-hour-studied ratio.
\end{quote}

\begin{quote}
\textbf{ PYQ Resources:}
- CSIR NET official papers: csirnet.nta.ac.in
- DIPS Academy question bank: dipsacademy.com
- IFAS solved papers: ifasonline.com
- pkalika.in (free solved PYQs for math)
\end{quote}

\newpage

\chapter{Additional Proofs \& Gap Filler}


\begin{quote}
\textbf{This file plugs every remaining gap across all 6 subjects}
\end{quote}

\bigskip\noindent\rule{\textwidth}{0.4pt}\bigskip


\subsection{Arzelà-Ascoli Theorem (Full Proof) ⚡⚡}

\begin{tcolorbox}[colback=lightblue,colframe=thmblue,title=\textbf{Theorem},breakable,boxrule=1.5pt]
\textbf{Theorem:} A sequence {f$_n$} in C[a,b] has a uniformly convergent subsequence iff it is uniformly bounded and equicontinuous.
\end{tcolorbox}

\begin{proof}
\textbf{Proof ($\Longrightarrow$):} If f$_{nk}$  f uniformly, then {f$_{nk}$} is uniformly bounded (f bounded + uniform tail < $\varepsilon$). Equicontinuity follows since f is uniformly continuous on [a,b] and the convergence is uniform.
\end{proof}

\begin{proof}
\textbf{Proof ($\Longleftarrow$ — the hard direction):}
\textbf{Step 1:} Let {r$_k$} = $\mathbb{Q}\cap$ [a,b] (countable dense set).
At r$_1$: {f$_n$(r$_1$)} is bounded  extract convergent subsequence {f$_1$,$_n$}.
At r$_2$: {f$_1$,$_n$(r$_2$)} is bounded  extract subsequence {f$_2$,$_n$}.
Continue diagonally: {f$_n$,$_n$} converges at EVERY rational.
\end{proof}

\textbf{Step 2:} Show {f$_n$,$_n$} is Cauchy in sup norm.
Given $\varepsilon$ > 0, by equicontinuity: $\exists\delta$ such that |x−y| < $\delta\Longrightarrow$ |f$_n$(x)−f$_n$(y)| < $\varepsilon$/3 for ALL n.

Cover [a,b] by finitely many intervals of length $\delta$, picking rationals r$_1$,...,r$_k$ in each.
Choose N so that |f$_m$,$_m$(r$_j$) − f$_n$,$_n$(r$_j$)| < $\varepsilon$/3 for all j and m,n $\geq$ N.

For any x $\in$ [a,b], pick r$_j$ with |x−r$_j$| < $\delta$:
\begin{center}
\begin{tabular}{|l|l|l|l|l|l|l|}
\hline
\textbf{f$_m$,$_m$(x) − f$_n$,$_n$(x)} & \textbf{$\leq$} & \textbf{f$_m$,$_m$(x)−f$_m$,$_m$(r$_j$)} & \textbf{+} & \textbf{f$_m$,$_m$(r$_j$)−f$_n$,$_n$(r$_j$)} & \textbf{+} & \textbf{f$_n$,$_n$(r$_j$)−f$_n$,$_n$(x)} \\
\hline
\end{tabular}
\end{center}
< $\varepsilon$/3 + $\varepsilon$/3 + $\varepsilon$/3 = $\varepsilon$

So {f$_n$,$_n$} is uniformly Cauchy  uniformly convergent (C[a,b] is complete). $\blacksquare$

\bigskip\noindent\rule{\textwidth}{0.4pt}\bigskip

\subsection{Stone-Weierstrass Theorem (Proof for polynomials) ⚡⚡}

\textbf{Weierstrass version:} Every continuous function on [a,b] can be uniformly approximated by polynomials.

\begin{proof}
\textbf{Proof (Bernstein's constructive proof):}
Define the nth Bernstein polynomial:
\end{proof}

B$_n$(f;x) = $\Sigma_k$₌$_0^n$ f(k/n) $\cdot$ C(n,k) $\cdot$ xᵏ(1−x)$^n^{-}$ᵏ

\textbf{Claim:} B$_n$(f;x)  f(x) uniformly on [0,1].

Key identity: $\Sigma_k$₌$_0^n$ C(n,k)xᵏ(1−x)$^n^{-}$ᵏ = 1 (binomial theorem).

Using f's uniform continuity: |f(x)−f(k/n)| < $\varepsilon$ when |x−k/n| < $\delta$.

Split the sum: terms with |k/n − x| < $\delta$ contribute < $\varepsilon$ (by uniform continuity), terms with |k/n − x| $\geq\delta$ are controlled by Chebyshev (their total weight is $\leq$ 1/(4n$\delta^2$)  0).

So |B$_n$(f;x) − f(x)| < $\varepsilon$ + 2M/(4n$\delta^2$)  $\varepsilon$ as n  $\infty$. Since $\varepsilon$ arbitrary, convergence is uniform. $\blacksquare$

\bigskip\noindent\rule{\textwidth}{0.4pt}\bigskip

\subsection{Inverse Function Theorem (Complete) ⚡⚡}

\begin{tcolorbox}[colback=lightblue,colframe=thmblue,title=\textbf{Theorem},breakable,boxrule=1.5pt]
\textbf{Theorem:} If f: $\mathbb{R}^n\mathbb{R}^n$ is C$^1$ and det(Df(a)) $\neq$ 0, then f is a local diffeomorphism near a: there exist neighborhoods U of a and V of f(a) such that f: U  V is bijective with C$^1$ inverse.
\end{tcolorbox}

\begin{proof}
\textbf{Proof sketch:}
1. WLOG f(a) = 0, Df(a) = I (compose with linear map).
2. Define g(x) = x − f(x). Then Dg(a) = 0, so ||Dg(x)|| < $\frac{1}{2}$ near a.
3. For fixed y near 0, define Tᵧ(x) = y + g(x) = y + x − f(x). Show Tᵧ is a contraction on a small ball.
4. By Banach fixed point, Tᵧ has unique fixed point x = Tᵧ(x), giving f(x) = y.
5. Continuity of f$^{-}^1$: follows from contraction estimates.
6. Differentiability of f$^{-}^1$: D(f$^{-}^1$)(y) = [Df(f$^{-}^1$(y))]$^{-}^1$ (by chain rule). $\blacksquare$
\end{proof}

\bigskip\noindent\rule{\textwidth}{0.4pt}\bigskip

\subsection{R-S Integral: Integration by Parts (Proof) ⚡}

\begin{tcolorbox}[colback=lightblue,colframe=thmblue,title=\textbf{Theorem},breakable,boxrule=1.5pt]
\textbf{Theorem:} $\int$ₐᵇ f d$\alpha$ + $\int$ₐᵇ $\alpha$ df = f(b)$\alpha$(b) − f(a)$\alpha$(a)
\end{tcolorbox}

\begin{proof}

For partition P, the Riemann-Stieltjes sum for $\int$f d$\alpha$ is:
S$_1$ = $\Sigma$ f(t$_i$)[$\alpha$(x$_i$) − $\alpha$(x$_i$₋$_1$)]
\end{proof}

Similarly for $\int\alpha$ df:
S$_2$ = $\Sigma\alpha$(s$_i$)[f(x$_i$) − f(x$_i$₋$_1$)]

Now S$_1$ + S$_2$ = $\Sigma$ [f(t$_i$)$\alpha$(x$_i$) − f(t$_i$)$\alpha$(x$_i$₋$_1$) + $\alpha$(s$_i$)f(x$_i$) − $\alpha$(s$_i$)f(x$_i$₋$_1$)]

By Abel summation (summation by parts), this telescopes to f(b)$\alpha$(b) − f(a)$\alpha$(a) in the limit as mesh  0. $\blacksquare$

\bigskip\noindent\rule{\textwidth}{0.4pt}\bigskip


\subsection{Tychonoff's Theorem (Statement + Proof Idea) ⚡⚡}

\begin{tcolorbox}[colback=lightblue,colframe=thmblue,title=\textbf{Theorem},breakable,boxrule=1.5pt]
\textbf{Theorem:} The product of compact spaces is compact (in product topology).
\end{tcolorbox}

\begin{proof}
\textbf{Proof idea (using Alexander subbasis theorem):}
1. \textbf{Alexander's Lemma:} If every cover by subbasis elements has a finite subcover, then X is compact.
2. Product topology has subbasis: $\pi$_$\alpha^{-}^1$(U$\alpha$) where U$\alpha$ is open in X$\alpha$.
3. Suppose a cover by subbasis elements has NO finite subcover. For each $\alpha$, the projections $\pi$_$\alpha^{-}^1$(U$\alpha$) that appear can't cover X$\alpha$ (otherwise finitely many would cover X by compactness of X$\alpha$ and the projection).
4. So for each $\alpha$, pick x$\alpha\notin\cup$(relevant U$\alpha$'s). Then x = (x$\alpha$) is not in any set of the cover — contradiction! $\blacksquare$
\end{proof}

\textbf{Note:} Full proof requires Zorn's lemma (equivalent to Axiom of Choice).

\bigskip\noindent\rule{\textwidth}{0.4pt}\bigskip

\subsection{Urysohn's Lemma (Complete) ⚡⚡}

\begin{tcolorbox}[colback=lightblue,colframe=thmblue,title=\textbf{Theorem},breakable,boxrule=1.5pt]
\textbf{Theorem:} If X is normal and A, B are disjoint closed sets, then $\exists$ continuous f: X  [0,1] with f(A) = {0}, f(B) = {1}.
\end{tcolorbox}

\begin{proof}
\textbf{Proof sketch:}
1. By normality, separate A from B: A $\subset$ U$_1$/$_2\subset$ Ū$_1$/$_2\subset$ X\B.
2. Iterate: separate A from X\U$_1$/$_2$ to get U$_1$/₄, and separate Ū$_1$/$_2$ from B to get U$_3$/₄.
3. Continue for all dyadic rationals r = p/2$^n\in$ [0,1]: get open sets Uᵣ with:
   A $\subset$ Uᵣ $\subset$ Ūᵣ $\subset$ Uₛ whenever r < s.
4. Define f(x) = inf{r : x $\in$ Uᵣ} (or 1 if x $\notin$ any Uᵣ).
5. f is continuous (check preimages of (a,$\infty$) and (−$\infty$,a) are open). $\blacksquare$
\end{proof}

\bigskip\noindent\rule{\textwidth}{0.4pt}\bigskip

\subsection{Bolzano-Weierstrass $\Longleftrightarrow$ Compactness in Metric Spaces ⚡⚡}

\begin{tcolorbox}[colback=lightblue,colframe=thmblue,title=\textbf{Theorem},breakable,boxrule=1.5pt]
\textbf{Theorem:} In a metric space, X is compact $\Longleftrightarrow$ every sequence has a convergent subsequence.
\end{tcolorbox}

\begin{proof}
\textbf{Proof ($\Longrightarrow$):} X compact. Let {x$_n$} be a sequence. If it has no convergent subsequence, then {x$_n$} has no limit point. So for each x $\in$ X, $\exists$Bₓ containing finitely many x$_n$. These Bₓ cover X. By compactness, finitely many cover X. But finitely many balls contain finitely many x$_n$ — can't cover infinitely many. Contradiction! $\blacksquare$
\end{proof}

\begin{proof}
\textbf{Proof ($\Longleftarrow$):} Use: sequentially compact $\Longrightarrow$ totally bounded + complete $\Longrightarrow$ compact (in metric spaces). $\blacksquare$
\end{proof}

\bigskip\noindent\rule{\textwidth}{0.4pt}\bigskip


\subsection{Schwarz Lemma ⚡⚡⚡ (frequently tested!)}

\begin{tcolorbox}[colback=lightblue,colframe=thmblue,title=\textbf{Theorem},breakable,boxrule=1.5pt]
\textbf{Theorem:} If f: D  D is analytic (D = unit disk), f(0) = 0, then:
1. |f(z)| $\leq$ |z| for all z $\in$ D
2. |f'(0)| $\leq$ 1
3. If |f(z$_0$)| = |z$_0$| for some z$_0\neq$ 0, or |f'(0)| = 1, then f(z) = e^(i$\theta$)z (rotation).
\end{tcolorbox}

\begin{proof}

g(z) = f(z)/z is analytic in D (removable singularity at 0, g(0) = f'(0)).
On |z| = r < 1: |g(z)| = |f(z)|/r $\leq$ 1/r.
By maximum principle: |g(z)| $\leq$ 1/r for |z| $\leq$ r.
Let r  1: |g(z)| $\leq$ 1 for all z $\in$ D.
So |f(z)| $\leq$ |z| and |f'(0)| = |g(0)| $\leq$ 1.
Equality case: |g| attains maximum 1 inside D $\Longrightarrow$ g is constant e^(i$\theta$) $\Longrightarrow$ f(z) = e^(i$\theta$)z. $\blacksquare$
\end{proof}

\bigskip\noindent\rule{\textwidth}{0.4pt}\bigskip

\subsection{Rouché's Theorem ⚡⚡⚡ (for counting zeros)}

\begin{tcolorbox}[colback=lightblue,colframe=thmblue,title=\textbf{Theorem},breakable,boxrule=1.5pt]
\textbf{Theorem:} If f, g are analytic inside and on a simple closed curve C, and |g(z)| < |f(z)| on C, then f and f+g have the same number of zeros inside C.
\end{tcolorbox}

\begin{tcolorbox}[colback=examplebg!30,colframe=exampleborder,title=\textbf{Worked Example},breakable,boxrule=1pt]
\textbf{Example:} Zeros of z⁵ + 3z + 1 in |z| < 1.
Take f = 3z, g = z⁵+1. On |z|=1: |f| = 3, |g| $\leq$ 1+1 = 2 < 3. 
f = 3z has 1 zero inside. So z⁵+3z+1 has \textbf{1 zero} in |z| < 1.
\end{tcolorbox}

\bigskip\noindent\rule{\textwidth}{0.4pt}\bigskip

\subsection{Residue Theorem Application ⚡⚡⚡}

\textbf{Computing real integrals using residues:}

\textbf{Type 1:} $\int_0^2\pi$ R(cos$\theta$, sin$\theta$) d$\theta$. Substitute z = e^(i$\theta$), use ∮.

\textbf{Type 2:} $\int$_{-$\infty$}^{$\infty$} f(x)dx. Close contour in upper half-plane, sum residues.

\begin{tcolorbox}[colback=examplebg!30,colframe=exampleborder,title=\textbf{Worked Example},breakable,boxrule=1pt]
\textbf{Example:} $\int$_{-$\infty$}^{$\infty$} dx/(1+x$^4$).
Poles of 1/(1+z$^4$) at z = e^(i$\pi$/4), e^(3i$\pi$/4), e^(5i$\pi$/4), e^(7i$\pi$/4).
Upper half-plane: z$_1$ = e^(i$\pi$/4), z$_2$ = e^(3i$\pi$/4).
Res(z$_1$) = 1/(4z$_1^3$) = e^(-3i$\pi$/4)/4, Res(z$_2$) = e^(-9i$\pi$/4)/4 = e^(-i$\pi$/4)/4.
Sum = (e^(-3i$\pi$/4) + e^(-i$\pi$/4))/4 = (−1/$\sqrt{2}$ − i/$\sqrt{2}$ + 1/$\sqrt{2}$ − i/$\sqrt{2}$)/4 = −i/(2$\sqrt{2}$).
Integral = 2$\pi$i $\cdot$ (−i/(2$\sqrt{2}$)) = \textbf{$\pi$/$\sqrt{2}$}.
\end{tcolorbox}

\bigskip\noindent\rule{\textwidth}{0.4pt}\bigskip


\subsection{How to Sketch Phase Portraits}

\textbf{Step 1:} Find eigenvalues of A.
\textbf{Step 2:} Find eigenvectors.
\textbf{Step 3:} Draw eigenvector directions.
\textbf{Step 4:} Draw trajectories based on type:

\subsubsection{Stable Node ($\lambda_1$ < $\lambda_2$ < 0)}
\begin{itemize}[leftmargin=*]
  \item Both eigenvector directions point INWARD
  \item Trajectories tangent to SLOW eigenvector ($\lambda_2$, less negative) far from origin
  \item Trajectories tangent to FAST eigenvector ($\lambda_1$, more negative) near origin
  \item ALL arrows point toward origin
\end{itemize}

\subsubsection{Unstable Node (0 < $\lambda_1$ < $\lambda_2$)}
\begin{itemize}[leftmargin=*]
  \item Same as stable node but ALL arrows point AWAY from origin
\end{itemize}

\subsubsection{Saddle ($\lambda_1$ < 0 < $\lambda_2$)}
\begin{itemize}[leftmargin=*]
  \item Stable manifold: along eigenvector of $\lambda_1$ (arrows IN)
  \item Unstable manifold: along eigenvector of $\lambda_2$ (arrows OUT)
  \item Generic trajectories: approach along unstable, then veer away following stable direction — hyperbola-like curves
\end{itemize}

\subsubsection{Stable Spiral ($\alpha\pm\beta$i, $\alpha$ < 0)}
\begin{itemize}[leftmargin=*]
  \item No real eigenvectors to draw
  \item Trajectories spiral INWARD
  \item Direction (CW or CCW) determined by sign of $\beta$ and matrix entries
  \item To determine direction: check where x' > 0 and y' > 0 at a test point
\end{itemize}

\subsubsection{Center ($\pm\beta$i)}
\begin{itemize}[leftmargin=*]
  \item Closed elliptical orbits
  \item NO spiraling — orbits are periodic
  \item Direction determined same as spiral
\end{itemize}

\subsubsection{Star Node ($\lambda_1$ = $\lambda_2$, 2 eigenvectors)}
\begin{itemize}[leftmargin=*]
  \item Trajectories are straight lines through origin (all directions)
  \item Inward if $\lambda$ < 0, outward if $\lambda$ > 0
\end{itemize}

\subsubsection{Improper Node ($\lambda_1$ = $\lambda_2$, 1 eigenvector)}
\begin{itemize}[leftmargin=*]
  \item Single eigenvector direction is dominant
  \item Other trajectories curve toward this direction
  \item Like a node but with a "twist"
\end{itemize}

\bigskip\noindent\rule{\textwidth}{0.4pt}\bigskip


\subsection{Euler's Equations Derivation ⚡⚡}

Starting from L̇ = N (torque) in the fixed frame. In the body frame (rotating with $\omega$):

(dL/dt)_body + $\omega\times$ L = N

With L = [I]$\omega$ and [I] = diag(I$_1$,I$_2$,I$_3$) in principal axes:

Component 1: I$_1\omega$̇$_1$ + ($\omega\times$ I$\omega$)$_1$ = N$_1$
$\omega\times$ I$\omega$ = ($\omega_2$I$_3\omega_3$ − $\omega_3$I$_2\omega_2$) ê$_1$ + ... = (I$_3$−I$_2$)$\omega_2\omega_3$ ê$_1$ + ...

So: \textbf{I$_1\omega$̇$_1$ + (I$_3$−I$_2$)$\omega_2\omega_3$ = N$_1$}
Or equivalently: \textbf{I$_1\omega$̇$_1$ − (I$_2$−I$_3$)$\omega_2\omega_3$ = N$_1$} $\blacksquare$

\subsection{Compound Pendulum — Minimum Period ⚡}

T$^2$ = 4$\pi^2$(k$^2$+h$^2$)/(gh) where h = pivot-to-CM distance, k = radius of gyration about CM.

dT$^2$/dh = 4$\pi^2$(2h$\cdot$gh − g(k$^2$+h$^2$))/(gh)$^2$ = 4$\pi^2$(h$^2$−k$^2$)/(gh$^2$)

Setting = 0: h$^2$ = k$^2$, so h = k.

T_min = 2$\pi$$\sqrt{2k/g}$. This equals the period of a simple pendulum of length l = 2k. $\blacksquare$

\newpage

\end{document}
